% 本文件由 LaTeX 排版 Agent 根据有效 translation.json 自动生成。
% author=Augustine of Hippo; work_id=1509; title=驳斥白拉奇派两函

\chapter{驳斥白拉奇派两函}

% structure: p0001 source=h1 target=omitted reason=page-title-already-rendered-by-parent

摘自\Person{奥斯定}的\WorkTitle{再论}卷二，第六十一章。\newline{} \Emph{随后有四本书是我写给罗马教会主教\newline{}波尼法爵的，以驳斥白拉奇派的两封信\newline{}；因为这些信到了他手中，他寄给了我，发现\newline{}在它们之中对我名字的诽谤性提及。此书如此\newline{}开篇：\Quote{我确实因你闻名遐迩的声誉而认识了你。}}\par

卷一\newline{}卷二\newline{}卷三\newline{}卷四\par

\SourceRecord{Translated by Peter Holmes and Robert Ernest Wallis, and revised by Benjamin B. Warfield.；Nicene and Post-Nicene Fathers, First Series；Vol. 5.；Edited by Philip Schaff.；Buffalo, NY: Christian Literature Publishing Co.,；1887.；<http://www.newadvent.org/fathers/1509.htm>.}

% structure: p0001 source=h1 target=section reason=page-title

\section{\WorkTitle{驳斥佩拉纠派两封书信 (第一篇)}}

奥斯定答复一封据说由儒利安寄往罗马的书信；首先，为其异端指控辩护大公教会的教义；然后，揭示并驳斥那封书信作者用以对抗大公教会的信仰宣言中所隐藏的佩拉纠派异端含义。\par

% structure: p0003 source=h2 target=subsection reason=relative-heading-rank; base=h2

\subsection{第一章——引言：致博尼法爵。}

我先前确已因您那显赫声誉的赞美而认识您，并通过众多极其可靠的传信者得知您充满天主的恩宠，至福可敬的教宗博尼法爵！但在我兄弟阿利比乌亲身见到您，蒙您以全部仁爱真诚接待，并因爱德之命与您交谈，与您同住，虽只短短时日，却以恳切之爱与您联合，将他与我一同倾注于您的心中，又将您带回于他的心中之后——您的友谊越是确信，我对您圣德的信念便越大。因为您虽居于高位，却不傲慢，不嫌弃作卑微者的朋友，并回报那给予您的爱。友谊之名岂非只源于爱，且只在基督内才真实，惟有在祂内才能永恒而幸福？因此，借着那位兄弟，我得以更亲密地认识您，从而更大胆地就那些在此关键时刻以新的刺激唤起主教职分（就我们所能）对主的羊群保持警醒之事，写些东西给您的真福。\par

% structure: p0005 source=h2 target=subsection reason=relative-heading-rank; base=h2

\subsection{第二章——为何必须答复异端著作。}

因为那些新的异端者，天主恩宠的敌人——这恩宠是由主耶稣基督赐给大大小小之人的——虽然他们已更公开地显明需要以明显的不赞同来避免，但他们仍不停以著作试探那些较为粗心、学识较少之人的心。这些著作当然必须答复，以免他们使自己或朋友陷入那邪恶的错误，即使我们并不担心他们可能以花言巧语欺骗某位大公信徒。但既然他们不停地在主的羊圈入口咆哮，并从四面八方撕开入口，企图撕裂那以如此高价赎回的羊群；既然牧者的守望塔为我们所有执行主教职分的人所共有（尽管您在其上处于更高地位），我尽己所能，在我所承担的小小部分职分内，正如主因您祈祷的帮助屈尊赐我力量，以医治性和防御性的著作来对抗他们那瘟疫般狡诈的著作，好使他们所狂怒的癫狂或得治愈，或不得伤害他人。\par

% structure: p0007 source=h2 target=subsection reason=relative-heading-rank; base=h2

\subsection{第三章——为何将本书献给博尼法爵。}

但我现正答复他们两封书信的这些文字——一封据说是儒利安寄往罗马的，我相信他借此要么找到、要么制造尽量多的盟友；另一封是十八位所谓的主教，与他同陷错误者，竟敢寄往得撒洛尼，不是寄给随便什么人，而是寄给该地的主教本人，意图以狡诈试探他，若能成事便将他拉入他们的观点——我所说的这些文字，是我在就该论题答复他们那两封书信时所写的，我决定特别献给您的圣德，与其说是为了您的学习，不如说是为了您的审阅，并且若偶有令您不悦之处，为了您的指正。因为我的兄弟向我透露，您亲自屈尊将那些书信交给他；这些书信若非经由我兄弟、您儿子的最警醒勤勉，必不能到您手中。我感谢您对我的极真诚仁爱，因您不愿那些天主恩宠之敌的书信对我隐藏，况且在其中您发现我的名字既被诽谤又公开提及。但我寄望于我的主天主，那些以毁谤之齿撕裂我的人，必不致无天上赏报；我为他们而反对自己，为保护小孩子，使他们不致被留给诡诈的谄媚者佩拉纠而遭毁灭，而能被呈交给真实的救主基督以获得拯救。\par

% structure: p0009 source=h2 target=subsection reason=relative-heading-rank; base=h2

\subsection{第四章【II.】——儒利安的诽谤：大公信徒教导自由意志因亚当之罪而被消灭。}

因此，现在让我们答复儒利安的书信。\Quote{那些摩尼教徒说，}他说，\Quote{现在我们不与共融的——也就是我们与之分歧的全部人——说，因第一人，即亚当，的罪，自由意志灭亡了；并且现在没有人有善生的能力，而是众人都被他们肉身的必然性所强迫而犯罪。}他称大公信徒为摩尼教徒，仿效多年前那个约维尼安的做法，他作为新异端者，破坏了蒙福玛利亚的童贞，并将信徒的婚姻置于与她神圣的童贞同等地位。他对此向大公信徒提出的反对，不是基于其他理由，而是因为他希望他们显得是婚姻的控告者或谴责者。\par

% structure: p0011 source=h2 target=subsection reason=relative-heading-rank; base=h2

\subsection{第五章——自由选择并未随亚当之罪而消灭。哪种自由消灭了。}

但在辩护自由意志时，他们更急于信赖自由意志去行义，而非依赖天主的助佑，并各自在自己身上夸耀，而不在主内夸耀。\BibleQuote{凡夸耀的，应当因主而夸耀。}\BibleRef{格前 1:31} 然而我们中间有谁会说，因第一个人的罪，自由意志便从人类中消失了呢？因罪，自由确实消失了，但那是乐园中的自由，即拥有不朽的圆满义德；正因如此，人性需要天主的恩宠，因为主说：\BibleQuote{如果圣子使你们自由，你们便真是自由的了。}\BibleRef{若 8:36}——当然是指能度良善正直的生活。因为罪人中的自由意志并没有消失到这种程度——即所有人都能藉它犯罪，尤其是那些带着喜悦和罪爱犯罪的人；他们所乐意做的事，便令他们喜悦。因此宗徒也说：\BibleQuote{当你们作罪的奴隶时，你们在义德上是自由的。}\BibleRef{罗 6:20} 看，他们被表明绝不可能凭着另一个自由去服事罪。所以，他们藉着意志的选择，在义德上是不自由的，但他们只能藉救主的恩宠才能脱离罪。为此，那位杰出的导师也区分了这些词：\BibleQuote{因为当你们作罪的奴隶时，你们在义德上是自由的。那么，你们那时在你们现在所羞耻的事上，有什么果实呢？那些事的结局就是死亡。但如今你们脱离了罪，成了天主的奴隶，你们的果实是成圣，结局是永生。}\BibleRef{罗 6:20-22} 他称他们从义德中\Quote{自由}，而不是\Quote{被释放}；但从罪中他称他们不是\Quote{自由}，以免他们将此归于自己，而是最谨慎地选用\Quote{被释放}一词，并将此联系到主的宣告：\BibleQuote{如果圣子使你们自由，你们便真是自由的了。}\BibleRef{若 8:36} 既然人类之子除非成为天主的子女，否则不能善度生活，为何这位作者希望赋予自由意志以善度生活的能力，而这能力只藉天主的恩宠，通过我们的主耶稣基督赐予？正如福音所说：\BibleQuote{凡接受祂的，祂赐给他们权能，成为天主的子女。}\BibleRef{若 1:12}\par

% structure: p0013 source=h2 target=subsection reason=relative-heading-rank; base=h2

\subsection{第六章 [III.]——恩宠不是按功德赐予的。}

但恐怕他们也许会说，他们被助佑是为了能\BibleQuote{成为天主的子女}\BibleRef{若 1:12}，而他们之能配得这权能，是藉自由意志不靠恩宠的助佑首先\BibleQuote{接受了祂}\BibleRef{若 1:12}（因为他们的努力正是要摧毁恩宠，以便主张恩宠是按我们的功德赐予的）；恐怕他们就这样分割那福音的陈述，将功德归于其中说\BibleQuote{凡接受祂的}\BibleRef{若 1:12}那部分，然后声称在接着说的\BibleQuote{祂赐给他们权能，成为天主的子女}\BibleRef{若 1:12}中，恩宠不是白白赐予，而是回报这功德；如果问他们\Quote{接受了祂}是什么意思，除了\BibleQuote{信了祂}\BibleRef{若 1:12}之外，还能说什么呢？那么，为了让他们知道这也属于恩宠，请他们读一读宗徒所说的话：\BibleQuote{凡事不怕敌人的惊吓；这证明他们沉沦，你们得救，并且是出于天主。因为为了基督，恩赐给你们，不但是信祂，而且是为祂受苦。}\BibleRef{斐 1:28-29} 他明确说两者都是恩赐。也请他们读他所说的：\BibleQuote{愿平安、慈爱、信德，从天主父和主耶稣基督，赐给弟兄们。}\BibleRef{弗 6:23} 也请他们读主自己所说的：\BibleQuote{除非派遣我的父吸引他，谁也不能到我这里来。}\BibleRef{若 6:44} 此处，以免有人认为\Quote{到我这里来}与\Quote{信我}意思不同，稍后当祂谈论自己的体和血，许多人因祂的话而跌倒时，祂说：\BibleQuote{我对你们所说的话，就是神，就是生命。但你们中间有些人并不信。}\BibleRef{若 6:63-64} 然后福音作者补充说：\BibleQuote{因为耶稣从起头就知道谁不信祂，谁要出卖祂。祂又说：为此，我对你们说过：除非蒙父恩赐，谁也不能到我这里来。}\BibleRef{若 6:64-65} 也就是说，祂重复了祂曾说的：\BibleQuote{除非派遣我的父吸引他，谁也不能到我这里来。}祂声明祂说这话是为了信者和不信者，藉着重复同样的话，用另一种说法解释了祂所说的\BibleQuote{除非派遣我的父吸引他}，即\BibleQuote{除非蒙父恩赐}。因为那被赐予信基督之恩的人，就被吸引到基督那里。因此，权能被赐予，使那些信祂的人成为天主的子女，因为正是这信本身被赐予。除非这权能来自天主，出于自由意志便不可能有；因为若释放者未使它自由，它就不能对善自由；但在恶中，他有自由意志，因为那欺骗者——或隐密或明显——已将恶爱植入他心中，或他自己说服了自己。\par

% structure: p0015 source=h2 target=subsection reason=relative-heading-rank; base=h2

\subsection{第七章——他结论说，他并未剥夺恶人的自由意志。}

因此，以下说法是不真实的：正如有些人断言我们所说的，而且你们的那位通信者竟敢写道，人\Quote{因肉身的必然性}而被\Quote{强迫犯罪}，仿佛他们不情愿似的；但若他们已到能运用自己心智之选择的年龄，则他们既因自己的意志而被羁留在罪中，也因自己的意志而从罪匆忙地奔向罪。因为即使那说服和欺骗人的，并不在他们内行动，除非他们因自己的意志而犯罪，或是由于对真理的无知，或是由于对邪恶的喜悦，或是由于两种邪恶——既是盲目的，也是软弱的。但这意志，它在邪恶中是自由的，因为它喜好邪恶；在善事中却不自由，因为它没有被释放。除非蒙那位不能愿意邪恶者——即藉着我们的主耶稣基督的天主恩宠——的帮助，人不能愿意任何善事。因为\BibleQuote{凡不出于信德的，都是罪。}\BibleRef{罗 14:23}因此，那从罪中退出的善意志是信德的，因为义人因信德而生活。\BibleRef{哈 2:4}并且相信基督属于信德。除非蒙赐予他，没有人能相信基督——即来到祂面前。\BibleRef{罗 1:17}因此，除非从上头领受了真正的、即白白赐予的恩宠，没有任何先前的功德，没有人能拥有正义的意志。\par

% structure: p0017 source=h2 target=subsection reason=relative-heading-rank; base=h2

\subsection{第八章 [IV.]——伯拉纠派摧毁自由意志。}

这些骄傲自大的人不愿这样；然而他们并非通过净化自由意志来维护它，而是通过夸大它来摧毁它。因为他们对我们这些说这些话的人发怒，没有别的原因，只是他们不屑于在主内夸耀。然而伯拉纠惧怕巴勒斯坦主教的审判；当有人指控他说天主的恩宠是按照我们的功德赐予时，他否认说过这话，并以绝罚谴责了那些这样说的人。然而在他后来写的书中，除了辩护这种观点之外，什么也没有发现；他以为通过说谎或以他不知如何的含糊措辞隐藏其意，欺骗了审判他的人。\par

% structure: p0019 source=h2 target=subsection reason=relative-heading-rank; base=h2

\subsection{第九章 [V.]——朱利安的另一诽谤——即\Quote{据说婚姻不是天主所制定的}。}

但现在让我们看看下面的话。\Quote{他们还断言，}他说，\Quote{那些现在举行的婚姻并不是天主指定的，这可以在奥斯定的书中读到，我曾以四卷书答复该书。而这些奥斯定的话已被我们的敌人以对真理的敌意所采用。}对于这些极其诽谤性的话，我认为必须简要地回答，因为他后来在想要暗示像他们这样的人会如何说的话时重复了这些话，仿佛是针对我的话。关于这一点，在天主的助佑下，我必须与他辩论，直到问题似乎需要为止。因此，现在我回答说，婚姻是由天主所制定的，无论是在当时，当它说：\BibleQuote{因此，人要离开父亲和母亲，依附自己的妻子，二人成为一体。}\BibleRef{创 2:24}还是现在，因此经上写道：\BibleQuote{女人是藉着主与男人结合的。}\BibleRef{箴 19:24}因为即使是现在，除了男人依附他的妻子并二人成为一体之外，没有别的发生。因为关于那现在缔结的婚姻本身，犹太人曾咨询主是否可以为任何原因休妻。祂在提及律法的见证时补充说：\BibleQuote{所以，天主所结合的，人不可拆散。}宗徒保禄在劝诫丈夫应当爱妻子时也应用了律法的这个见证。\BibleRef{弗 5:25}因此，愿那人在我的书中读到任何与这些神圣见证相悖的话的想法远离我们！但他或出于不理解，或更确切地说出于诬蔑，试图将他所读到的曲解为另一种含义。但我写了我的书——他提到曾以四卷书答复该书——是在伯拉纠和塞莱斯定被定罪之后。这一点我认为必须说明，因为那人声称我的话已被他的敌人以对真理的敌意所采用，以免有人以为这些新的异端者因我这本书而被定罪为基督恩宠的敌人。但在那本书中，所发现的是对婚姻的辩护而非指责。\par

% structure: p0021 source=h2 target=subsection reason=relative-heading-rank; base=h2

\subsection{第十章——第三种诽谤——断言夫妻房事被谴责。}

\Quote{他们还断言，}他说，\Quote{性冲动和已婚者的交合是由魔鬼发明的，因此那些生来无辜的人是有罪的，并且他们生于这魔鬼般的交合是魔鬼的工作，而非天主的工作。而这一点，毫无歧义，是摩尼教。}不，正如我说婚姻是由天主为生育子女的秩序而设立的，同样我说，即使在天堂里，如果没有性冲动和夫妻的交合，要生育的子女的繁衍就不可能发生，如果在那里生育子女的话。但是，这样的冲动和交合是否会像现在这样带着可耻的情欲存在，如果没有人犯罪的话，这是一个问题，我将在以后讨论，如果天主愿意的话。\par

% structure: p0023 source=h2 target=subsection reason=relative-heading-rank; base=h2

\subsection{第十一章 [VI.]——伯拉纠派赞扬夫妻房事的清白的目的。}

然而，他们究竟想要什么，追求什么，企图达到什么结果，那位作者附加的话表明了，当他断言我说：\Quote{因此，那些生来无罪的人是有罪的，并且他们从这种魔鬼般的交合中出生，是魔鬼的工作，而非天主的工作。}既然我既不说这种夫妻间的交合（尤其是在信徒中）是魔鬼般的——这交合是为了生育后来要重生的孩子；也不说任何人是由魔鬼造的，而是就他们是人而言，是由天主造的；尽管如此，即使是信主的夫妻所生的孩子也是有罪的人（如同从橄榄树结出野橄榄一样），这是由于原罪，因此他们在魔鬼的权势之下，除非他们在基督内重生，因为魔鬼是过错的作者，而非本性的作者：那么，另一方面，那些说婴儿没有原罪、因此不在魔鬼权势之下的人，他们劳心费力地要达成什么，除了使天主在婴儿身上的恩宠归于无效？这恩宠使我们，如宗徒所说，脱离了黑暗的权势，迁到了祂爱子的国度里。\BibleRef{格前 1:13} [VII.] 当他们否认婴儿在得救之主帮助之前就已经在黑暗的权势之下时，他们如此赞美造物主在婴儿身上的工作，以至于摧毁了救赎者的仁慈。因为我承认这在成年人和婴儿身上都是如此，所以他说这毫无疑问是摩尼教，尽管这是最古老的天主教教义，而这些人新奇的异端教义正是被它推翻的。\par

% structure: p0025 source=h2 target=subsection reason=relative-heading-rank; base=h2

\subsection{第十二章——第四个诽谤——说旧约的圣人并非无罪}

他说：\Quote{他们说，旧约中的圣人并非无罪——也就是说，他们即使改过也不能摆脱罪行，而是带着罪孽被死亡抓住。}不，我说，无论在法律之前还是在旧约时代，他们都是从罪中得了释放——不是凭自己的能力，因为\BibleQuote{凡信赖世人的人，是可咒骂的}，\BibleRef{耶 17:5} 而且无疑那些神圣咏所提及的\BibleQuote{信赖自己力量的人}也属于这咒诅之下；也不是凭那旧盟约，它使人处于奴役之中，\BibleRef{迦 4:24} 尽管它是因着某种确定的安排之恩宠而由天主赐下的；也不是凭那法律本身，它虽是神圣、公义、良善的，其中写着\BibleQuote{不可贪恋}，\BibleRef{出 20:7} 因为这法律并非赐下以赐予生命，而是为了过犯而添加的，直到那蒙应许的后嗣来到；而是说，他们是由救赎者本人的血所释放的，祂是天主与人之间的唯一中保，就是降生成人的基督耶稣。\BibleRef{弟前 2:5} 但是那些天主的恩宠——这恩宠是藉着我们的主耶稣基督赐给大小之人的——的仇敌却说，古时的天主之人具有完美的义德，免得他们被认为需要基督的降生成人、受难和复活，而他们正是因相信祂而得救的。\par

% structure: p0027 source=h2 target=subsection reason=relative-heading-rank; base=h2

\subsection{第十三章 [VIII.]——第五个诽谤——说保禄和其余宗徒被情欲玷污}

他说：\Quote{他们说连宗徒\Person{保禄}，甚至所有的宗徒，都常被无度的情欲所玷污。}但无论何等亵渎的人，谁敢说这话？但无疑此人如此歪曲，是因为他们主张，宗徒所说的：\BibleQuote{我知道在我内，即在我肉身内，没有善住着，因为立志行善由得我，但行出来却由不得我} \BibleRef{罗 7:18}以及诸如此类的话，并非指自己，而是假借了某个我不知道是谁的、正在承受这些痛苦的人。因此，必须仔细考查并研究他书信中的那一段，以免他们的错误隐没在他话语的晦暗之中。因此，尽管宗徒在这里广泛辩论，并以伟大而持久的争战维护恩宠，反对那些以法律自夸的人，但我们确实接触到一些与当前问题相关的事。关于这主题，他说：\BibleQuote{因为凡有血肉的，没有一人因法律而在天主面前成义；因为法律只叫人认识罪过。但如今，天主的义德在法律之外已显示出来，为法律和先知所证明，就是天主的义德，因耶稣基督的信德，为所有相信的人。原来没有分别；因为众人都犯了罪，亏缺了天主的荣耀，如今却白白地因他的恩宠成义，因基督耶稣的救赎。} \BibleRef{罗 3:20}又说：\BibleQuote{那里有可夸的？已被排除。借什么法律？是行为的法律吗？不是的，是信德的法律。因此我们认为：人成义是因信德，而不是因遵守法律的行为。} \BibleRef{罗 3:27}又说：\BibleQuote{因为天主应许亚巴郎和其后裔，作世界的承继者，不是凭法律，而是凭由信德而来的正义。因为若是凭法律的人才是承继者，那么信德便是空虚的，恩许便是无效的。因为法律只能激起义怒，那里没有法律，那里就没有违犯。} \BibleRef{罗 4:13}在另一处：\BibleQuote{法律本是外添的，为叫过犯增多；但罪恶在那里增多，恩宠在那里也格外丰盛。} \BibleRef{罗 5:20}又在另一处：\BibleQuote{罪恶不能统辖你们，因为你们不在法律之下，而是在恩宠之下。} \BibleRef{罗 6:14}又在另一处：\BibleQuote{弟兄们！我对明白法律的人说：你们岂不知道法律管束人，只在人活着的时候吗？因为有丈夫的女人，在她丈夫活着的时候，是受法律约束的；但如果丈夫死了，她就脱离了丈夫的法律。} \BibleRef{罗 7:1}稍后：\BibleQuote{所以我的弟兄们！你们借着基督的身体，已死于法律，为使你们属于另一人，就是从死者中复活的那一位，好使我们结出果实归于天主。因为我们还在肉体时，那借着法律而在我们肢体中活动的罪恶的情欲，结果实归于死亡；但现在我们已脱离了那捆绑我们的法律，得以在圣神的新里事奉，而不在文字的旧里。} \BibleRef{罗 7:4-6}诸如此类的见证，外邦人的教师充分证明了法律不能除去罪恶，反而使之增多，而恩宠却除去罪恶；因为法律知道如何命令，而对这命令的软弱屈服，恩宠却知道如何帮助，由此灌输了爱。恐怕有人因这些见证而指责法律，并主张它是邪恶的，宗徒预见那些误解的人可能发生的问题，便亲自向自己提出了同样的问题。他说：\Quote{那么我们能说什么呢？}他说：\BibleQuote{法律是罪吗？绝对不是！但除非借着法律，我就不知道罪恶。} \BibleRef{罗 7:7}他之前已经说过：\BibleQuote{因为法律只叫人认识罪过。} \BibleRef{罗 3:20}所以，法律不是罪的除去，而是罪的知识。\par

% structure: p0029 source=h2 target=subsection reason=relative-heading-rank; base=h2

\subsection{第十四章——宗徒是以自己以及那些在恩宠之下、而非仍在法律之下的人的身份说话。}

从这里，他现在开始——正是为此我着手思考这些事——引入他自己的人格，仿佛讲论自己；而白拉奇派不愿承认宗徒这里的话是指自己，而说他将另一个人物转变成自己——即一个仍在法律之下、尚未被恩宠释放的人。的确，他们至少应当承认，如同这位宗徒在其他地方所说，\BibleQuote{在法律之下没有人成义}；但法律有助于认识罪，以及法律的过犯本身，以便罪被认识并增多，恩宠能借着信德被寻求。但他们不畏惧那些事情应当被理解为宗徒所说的，他也可能论及他的过去，但他们畏惧接下来的事。因为这里他说：\BibleQuote{若不是法律说\InnerQuote{不可贪恋}，我就不知情欲是什么。但罪恶趁着机会，借着诫命在我内生出各样的贪心。因为没有法律，罪是死的。我以前没有法律是活着的，但诫命一来，罪就活了，我却死了；那本是叫人活的诫命，反倒叫我死。因为罪恶趁着机会，借着诫命诱骗了我，并借着它杀了我。所以法律是圣的，诫命也是圣的、公义的、良善的。那么，那良善的是叫我死吗？绝对不是！但罪恶借着那良善的，在我内产生死亡，为显出罪恶真是罪，以致罪恶借着诫命变成极度的邪恶。} \BibleRef{罗 7:7-13}这一切，如我所说，宗徒似乎是在回忆他的过去生活：因此从他说\BibleQuote{我以前没有法律是活着的} \BibleRef{罗 7:9}，他可能希望被理解为他的早年时期，即从婴孩至理性年龄之前；但他加上了\BibleQuote{但诫命一来，罪就活了，我却死了} \BibleRef{罗 7:9}，他愿意表明自己能够接受诫命，却不能遵行，因此成为法律的违犯者。\par

% structure: p0031 source=h2 target=subsection reason=relative-heading-rank; base=h2

\subsection{第十五章 [IX.]——仅因惧怕而不敢犯罪的人，在意志上已犯罪。}

我们也不要因他写给斐理伯人的信而感到困扰：\BibleQuote{就法律上的正义说，是无可指摘的。}因为他在内心怀着邪恶的情欲时，可能是法律的违犯者，却仍能完成法律外在的功行，或因畏惧人，或因畏惧天主本身；但那是出于对惩罚的恐惧，而非出于对正义的爱与喜悦。因为怀着行善的意愿去行善，与怀着行恶的意愿以至于若非因惩罚阻止就会实际行恶，这是截然不同的。因为这样，他确实在内心的意志中犯罪，因为他不犯罪不是出于意志，而是出于恐惧。宗徒深知自己在这些内在情感中曾是那样的，在藉着我们的主耶稣基督而来的天主的恩宠之前，他在别处非常清楚地承认了这一点。因为他在写给厄弗所人的信中说：\BibleQuote{而你们，你们从前因你们的过犯和罪恶原是死的，那时你们在过犯和罪恶中生活，随从了这世界的潮流，顺从了空中权能的首领，即现今在悖逆之子身上运行的邪神；我们从前也都在他们中间，放纵肉体的私欲，随从肉体和思念的意愿而行，且生来就是义怒之子，和别人一样。然而富于慈悲的天主，因祂爱我们的大爱，竟在我们因过犯死了的时候，使我们同基督一起生活——你们得救是因着恩宠。}又给弟铎写信说：\BibleQuote{因为我们从前也是昏愚的、不信的、迷途的，服事各种私欲和逸乐，生活在恶毒和嫉妒中，是可憎的，又彼此仇恨。}这就是撒乌耳，当他说自己就法律上的正义而言是无可指摘时的情形。因为他在法律上并没有前进，并在这种可恨的生活之后改变性情而变得无可指摘——他在接下来的话中清楚地表明，他并非从这些恶事中转变，除非藉着救主的恩宠。因为他也在这里，如同在厄弗所书一样，加上这一点说：\BibleQuote{但当我们的救主天主的良善和慈爱显现时，祂救了我们，并不是由于我们凭义德所立的功劳，而是按照祂的怜悯，藉着圣神所施行的重生和更新的洗礼，这圣神乃是祂藉着我们的救主耶稣基督，厚厚地倾注在我们身上的，好使我们因祂的恩宠成义，而能成为永生的希望继承人。}\par

% structure: p0033 source=h2 target=subsection reason=relative-heading-rank; base=h2

\subsection{第十六章——罪如何死了，又如何活了}

他在《罗马书》那一段所说的\BibleQuote{罪，为显示它是罪，藉着善给我带来了死亡}\BibleRef{罗 7:13}，与他先前所说的\BibleQuote{我只是藉着法律才认识罪，因为如果不是法律说\InnerQuote{不可贪恋}，我就不认识贪恋}\BibleRef{罗 7:7}是一致的。而且先前有\BibleQuote{藉着法律，罪的认知}\BibleRef{罗 7:7}，因为他也在这里说了\BibleQuote{为显示它是罪}；这样我们就不会把他所说的\BibleQuote{没有法律，罪是死的}理解为仿佛罪不存在，\BibleQuote{它隐藏着，不显现，完全被忽略，如同埋在某个无知的黑暗中。}至于他说\BibleQuote{我以前没有法律是活着的}，除了是说\BibleQuote{我自以为活着}之外，还能是什么呢？至于他加上的\BibleQuote{但诫命来到，罪又活了}，这不过是罪显明了，成为了显明的。他说的不是\BibleQuote{活了}，而是\BibleQuote{又活了}。因为从前在乐园里它就活着，在违反命令时充分显现了；但当它遗传给来到世上的儿女时，它就隐藏着，仿佛死了，直到它的邪恶因禁止而抵抗正义时被人感觉到——当一件事被命令和认可，另一件事却令人喜悦和统治时，这时，罪在出生之人身上某种程度地又活了，尽管它最初在最初受造之人身上已经存活了一段时间。\par

% structure: p0035 source=h2 target=subsection reason=relative-heading-rank; base=h2

\subsection{第十七章［第十节］—— \BibleQuote{法律是属神的，但我是属肉体的}，应理解为保禄}

但接下来的话如何能理解为保禄自己说的，并不那么清楚。他说：\BibleQuote{我们知道法律是属神的，但我是属肉体的。}\BibleRef{罗 7:14}他没有说\BibleQuote{我曾是}，而是说\BibleQuote{我是}。那么，当保禄写这话时，他难道是属肉体的吗？或者他是就其身体而言说的？因为他仍在这个死亡的身体内，尚未成为他在别处所说的：\BibleQuote{播种的是属生理的身体，复活的是属神的身体。}因为到那时，他整个自己——即他所构成的这两部分——将成为属神的人，那时连身体也要成为属神的。因为在今生，那些仍思念地上的事的人，连灵也可能是属肉体的，这样，即使在今生说肉体是属肉的，也不为怪。因此，他说\BibleQuote{但我是属肉体的}，是因为保禄尚未有属神的身体，正如他可以说\BibleQuote{但我是必死的}，这无疑只能从身体方面理解，因为身体尚未披上不朽。此外，关于他加上的\BibleQuote{被卖给了罪}\BibleRef{罗 7:14}，免得有人认为他尚未被基督的血所救赎，这也可以从他所说的话来理解：\BibleQuote{我们这蒙受圣神初果的，也自己叹息，等待义子的名分，即我们身体的救赎。}\BibleRef{罗 8:23}因为如果是在这方面他说自己被卖给了罪——即他的身体尚未从朽坏中被救赎；或者他曾因最初违反诫命而被卖，以致有朽坏的身体拖累灵魂\BibleRef{智 9:15}——那么，有什么妨碍我们理解宗徒在此所说的关于他自己的话，即使他以第一人称说话时，希望不仅他自己，而且所有意识到自己从未在灵魂的喜乐中与肉体的情欲争战的人，都被理解为包括在内呢？\par

% structure: p0037 source=h2 target=subsection reason=relative-heading-rank; base=h2

\subsection{第十八章——宗徒如何说，他作了所不愿的恶。}

或者我们偶然担心接下来的话：\BibleQuote{因为我所作的，我不明白；我所愿意的，我偏不作；我所憎恨的，我反而去作。}\BibleRef{罗 7:15}以免有人从这些话怀疑宗徒是赞同肉身私欲偏情的恶行？但我们必须思量他所补充的：\BibleQuote{但我若作我所不愿意的，我便是同意法律是善的。}因为他说他宁愿同意法律，而非肉身私欲偏情。因他称此为\Quote{罪}。因此他说他行动和劳苦，并非出于同意和满足的欲望，而是出于私欲本身的冲动。所以接着他说：\BibleQuote{我同意法律是善的。}我同意，因为我不愿意法律所不愿意的。随后他说：\BibleQuote{这样，现在不是我作这事，而是在我内的罪作的。}\BibleRef{罗 7:17}他所谓\BibleQuote{这样，现在}是什么意思？无非是，如今在恩宠之下，恩宠已将我的意志的喜悦从私欲的同意中释放出来。因为\BibleQuote{不是我作这事}不能更好地理解为：他不同意把他的肢体献给不义作为罪的工具。因为他若贪欲、同意并行动，怎能说他本人没有作这事呢？即便他可能为作了这事而悲伤，并为被战胜而深深叹息？尽管那说\BibleQuote{不再是我作这事}的人并不同意这恶？\par

% structure: p0039 source=h2 target=subsection reason=relative-heading-rank; base=h2

\subsection{第十九章——什么是成就善事。}

而接下来的话岂非最清楚地显示他说话出自何处？\BibleQuote{因为我知道，在我内，即在我肉身内，没有善住着。}\BibleRef{罗 7:18}因为他若没有以\BibleQuote{即在我肉身内}的解释补充他所说的，则当他说的\BibleQuote{在我内}时，或许会被另作解释。因此他以另一种形式重复并强调同一件事：\BibleQuote{愿意善就近我，但成就善事却不在。}\BibleRef{罗 7:18}因为成就善事，就是人甚至不贪欲。因为当人贪欲时，善便不完全，即使人不同意贪欲的恶。\BibleQuote{我所愿意的善，}他说，\BibleQuote{我不去作；我所不愿意的恶，我反而去作。现在，我若作我所不愿意的，便不是我作这事，而是在我内的罪作的。}\BibleRef{罗 7:20}他庄重地重复了这一点，仿佛要唤醒最懒惰的人从沉睡中醒来：\BibleQuote{那么我发现，法律为我这愿意行善的人，恶却与我同在。}\BibleRef{罗 7:21}那么，法律是为愿意行善的人，但恶因贪欲而存在，尽管那说\BibleQuote{不再是我作这事}的人并不同意这恶。\par

% structure: p0041 source=h2 target=subsection reason=relative-heading-rank; base=h2

\subsection{第二十章——在我内，即在我肉身内。}

他在接下来的话中更清楚地表明了这两点：\BibleQuote{因为我按内在的人，喜悦天主的法律；但我看见在我肢体中另有一个法律，反抗我理性的法律，并将我掳去，顺从我肢体中罪的法律。}\BibleRef{罗 7:21}但他说\BibleQuote{将我掳去}，他可以在不赞同的情况下感受到情感。因此，由于那三件事——我们已经讨论过其中两件，即他说\BibleQuote{我是属血肉的}和\BibleQuote{被卖给罪}，以及这第三件\BibleQuote{在罪的法律中，即在我肢体中，将我掳去}——宗徒似乎是在描述一个仍在法律之下、尚未在恩宠之下的人。但正如我解释了前两句是针对仍可朽坏的肉身，同样这句也可理解为好像他说\BibleQuote{将我掳去}是在肉身中，而非在理性中；在情感中，而非在同意中；因此\BibleQuote{将我掳去}，因为即使肉身中也没有异类的本性，而是我们自己的。因此，正如他自己解释了他所说的：\BibleQuote{因为我知道，在我内，即在我肉身内，没有善住着}，同样，现在从那解释中也应学会这段落的含义，好像他说\BibleQuote{将\Emph{我}掳去}，即\Quote{我的肉身}，\BibleQuote{顺从我肢体中罪的法律}。\par

% structure: p0043 source=h2 target=subsection reason=relative-heading-rank; base=h2

\subsection{第二十一章——在基督耶稣内没有定罪。}

然后他补充他为什么要说这一切：\BibleQuote{我这个人真不幸呀！谁能救我脱离这该死的肉身呢？天主的恩宠，借着我们的主耶稣基督！}由此他得出结论：\BibleQuote{这样，我自己在理性上服从天主的法律，在情欲上却服从罪的法律。}\BibleRef{罗 7:24}即，在情欲上服从罪的法律，由于贪欲；但在理性上服从天主的法律，由于不同意那贪欲。\BibleQuote{如今，在基督耶稣内的，已没有定罪了。}\BibleRef{罗 8:1}因为那不同意肉身贪欲之恶的人，不被定罪。\BibleQuote{因为生命之神的法律在基督耶稣内释放了你，使你脱离罪与死的法律。}这样，肉身的贪欲就不能占有你的同意。而接下来的话越来越表明同一含义。但必须节制。\par

% structure: p0045 source=h2 target=subsection reason=relative-heading-rank; base=h2

\subsection{第二十二章——为何所引段落必须理解为在恩宠之下的人。}

我也曾以为宗徒在此论证中描述了一个在法律之下的人。但后来我被他所说的话迫使放弃这一想法，因为他说：\BibleQuote{这样，现在不是我作这事。}这与他随后所说的也相符：\BibleQuote{因此，如今在基督耶稣内的，已没有定罪了。}因为我不明白一个在法律之下的人怎能说：\BibleQuote{我按内在的人喜悦天主的法律}；因为这喜于善，由此他不赞同恶，并非出于对刑罚的恐惧，而是出于对正义的爱（因为\Quote{喜悦}即意味着这一点），只能归因于恩宠。\par

% structure: p0047 source=h2 target=subsection reason=relative-heading-rank; base=h2

\subsection{第二十三章——什么是脱离这该死的肉身。}

因为当他说：\Quote{谁要从这死亡的身体中救我出来呢？}\BibleRef{罗 7:24}时，谁能否认当宗徒说这话时他仍在这死亡的身体中呢？当然，恶人不会从这身体中被救出来，同样的身体会归还给他们作为永罚。因此，从这死亡的身体中被救出来，就是从一切肉欲的软弱中得到医治，并领受身体，不是为了惩罚，而是为了光荣。与这段话充分一致的还有这句话：\Quote{我们这已获得圣神初果的人，自己也在心中叹息，等待着义子名分，我们身体的救赎。}因为我们所叹息的，正是那叹息，其中我们说：\Quote{我真不幸啊！谁要从这死亡的身体中救我出来呢？}此外，当他说：\Quote{因为我所做的，我不了解}时，这不就是指：\Quote{我不愿意，我不赞同，我不同意，我不去做}吗？否则，就与他上面所说的\Quote{法律叫人知罪}和\Quote{我若不是因法律，就不知罪}以及\Quote{罪借着那善的，在我身上产生了死亡，好成为罪}相矛盾。因为，他如何借法律知道了他所不知的罪？不知的罪如何显现？因此，说\Quote{我不了解}等于说\Quote{我不做}，因为我并没有以我的同意去犯那罪；正如主对恶人说的：\Quote{我不认识你们}\BibleRef{玛 7:23}，尽管毫无疑问，没有什么能隐藏于祂；又如经上所说：\Quote{那位没有认识罪的}\BibleRef{格后 5:21}，意思是指那位没有犯过罪的，因为祂没有认识祂所定罪的。\par

% structure: p0049 source=h2 target=subsection reason=relative-heading-rank; base=h2

\subsection{第二十四章——他得出结论说宗徒是以己身和那些在恩宠之下的人的身份说话。}

经过对这些以及这经文上下文中同类事情的仔细考虑，可以正确理解宗徒不仅以己身指代他自己，也指代那些同样处于恩宠之下的人，他们与他一样，尚未处于完美的平安中，在那平安中死亡将被胜利吞灭\BibleRef{格前 15:54}。关于此事，他后来又说：\Quote{如果基督在你们内，身体固然因罪而死，但神魂因义而成活。然而，如果那使耶稣从死者中复活的圣神住在你们内，那使基督从死者中复活的，也要借着那住在你们内的圣神，使你们有死的身体活过来。}\BibleRef{罗 8:10}因此，在我们有死的身体活过来之后，不仅不会同意犯罪，甚至连那未曾同意的肉体的情欲本身也不会存留。而在有死的肉身中没有这种对圣神的反抗，只可能在那位不是借肉身而来到人间的唯一者身上。至于宗徒们，因为他们是人，在这有死之生中带着一个败坏而沉累灵魂的身体\BibleRef{智 9:15}，因此就\Quote{总是被过度的情欲所玷污}，正如那人恶意断言的那样——我绝不敢这样说。但我要说，尽管他们免于同意堕落的贪欲，但他们仍因那用克制加以约束的肉身的贪欲而叹息，怀着如此谦卑和虔诚，以致他们宁愿没有它，而不是制服它。\par

% structure: p0051 source=h2 target=subsection reason=relative-heading-rank; base=h2

\subsection{第二十五章　第十二节——第六个诽谤：奥古斯丁断言甚至基督也并非无罪的。}

同样，关于他所添加的，即我说\Quote{基督甚至也并非无罪的，反而因肉身的必然性说了谎话，并被其他过犯所玷污}，他应当看清楚他是从谁那里听到这些，或在谁的书信中读到这些；因为，他或许没有理解，而用恶意的欺骗扭曲成了诽谤性的含义。\par

% structure: p0053 source=h2 target=subsection reason=relative-heading-rank; base=h2

\subsection{第二十六章　第十三节——第七个诽谤：奥古斯丁断言在圣洗中并非所有罪都被赦免。}

\Quote{他们也说，}他说，\Quote{圣洗并不给予完全的赦免，也没有除去罪行，而是将它们剃掉，以至于所有罪的根源都保留在邪恶的肉身中。}除了不信的人，谁能对着佩拉齐亚人断言这事呢？因此，我说圣洗赦免所有罪，除去罪责，而不是剃掉它们；而且\Quote{所有罪的根源}并不\Quote{保留在邪恶的肉身中，如同剃过的头发在头上一样，从那里罪可以重新长出被剪除。}因为正是我发明了那个比喻，为了他们能用来进行诽谤，仿佛我是这样想的和说的。\par

% structure: p0055 source=h2 target=subsection reason=relative-heading-rank; base=h2

\subsection{第二十七章——在何种意义上，情欲在重生者中被称作罪。}

但关于他们所说的私欲偏情，我相信他们或是受骗，或是骗人；因为连已受洗的人也必须以虔诚的心志与它争斗，无论他如何努力向前，并被圣神引导。然而，虽然这被称为罪，但它之被称为罪，并非因为它本身是罪，而是因为它是因罪而产生的，正如某人的笔迹被称为他的\Quote{手}，是因手写了它。而那些真正的罪，是那些按肉体的私欲或无知而非法行、说、思的事——这些事一旦发生，若不被宽恕，就使行的人有罪。这私欲偏情本身在洗礼中被如此去除，尽管它由所有出生的人所继承，却绝不伤害重生的人。然而，若他们按肉身生育子女，它又被传递下去；它又将伤害那些出生的人，除非他们以同样的形式在重生时获得赦免，并且它留在他们内，绝不阻碍将来的生命，因为它的罪债虽由生育而来，却已由重生除去；因此现在它不再是罪，而被称为罪，或许是因为它因罪而成其所是，或因它被犯罪的快乐所激动，虽然因正义之快乐的得胜，并不向其表示同意。也不是因着这个——它的罪债已在重生的洗涤中被除去——受洗的人在祈祷中说：\Quote{宽免我们的罪债，如同我们也宽免得罪我们的人。} \BibleRef{玛 6:12} 而是因着那犯罪的事，无论是对它表示同意，即当正义被那令人愉悦之事所胜，还是因无知而将恶当作善。这些事是通过行为、言谈或——最容易也最迅速的——思想而犯的。从这一切中，有哪个信徒敢夸口自己的心是纯洁的？或者谁夸口自己纯洁无罪？ \BibleRef{箴 20:9} 当然，祷文中的下文是为了私欲偏情而说的：\Quote{不要让我们陷于诱惑，但救我们脱离凶恶。}因为正如经上记载：\Quote{各人被诱惑，乃是被自己的私欲偏情所牵引、所诱惑；私欲偏情既怀孕，就生出罪来。} \BibleRef{雅 1:14}\par

% structure: p0057 source=h2 target=subsection reason=relative-heading-rank; base=h2

\subsection{第二十八章 [XIV]——许多人无罪行，无人无罪。}

私欲偏情的一切产物以及私欲偏情本身的旧罪债，都藉洗礼的洗涤被除去。而私欲偏情现在所生出的一切，若不是那些不仅被称为罪、甚至被称为罪行的产物，就藉每日祈祷的方法得以净化，即当我们说：\Quote{宽免我们的罪债，如同我们也宽免}，并藉着真诚的施舍。因为没有人会愚蠢到说主的那条诫命不适用于受洗的人：\Quote{宽免，也必蒙宽免；给予，也必蒙给予。} \BibleRef{路 6:37} 但若宗徒在说\Quote{若有人无罪行}之处——他说的是\Quote{若有人无罪行} \BibleRef{铎 1:6}——或是说\Quote{没有罪}之处——他说的是\Quote{没有罪行} \BibleRef{弟前 3:10}——那么无人能被正当地祝圣为教会中的牧职。因为许多受洗的信徒没有罪行，但我要说，在此生中无人无罪——无论贝拉基派人士何等自高自大，并因我这样说而狂怒破裂；这不是因为还有任何罪的残余在洗礼中未被赦免，而是因为我们留在此生软弱中的人，每日不断地犯下那些罪，而这些罪日日被赐给那些以信德祈祷、以仁慈工作的人。这是公教信德的正统道理，是圣神处处播种的——并非异端邪说的虚浮和骄傲自大。\par

% structure: p0059 source=h2 target=subsection reason=relative-heading-rank; base=h2

\subsection{第二十九章 [XV]——犹利安以他朋友的信德反对公教信徒的意见。首先，关于自由意志。}

因此，现在让我们看看——在他以为可以诽谤性地反对我所信的事，并捏造我不信的事之后——他本人如何陈述他自己的信德或贝拉基派的信德。他说：\Quote{与此相反，我们每日辩论，且不愿向违犯者让步，因为我们说自由意志是众人本性所固有的，且不能因亚当的罪而丧失；这论断为一切圣经的权威所证实。}若多少有必要说这话，你就不应反对天主的恩宠说这话——你不应向违犯者让步，而应纠正你的意见。但关于此事，我已尽力，且在我认为足够之处，在上文辩论过了。\par

% structure: p0061 source=h2 target=subsection reason=relative-heading-rank; base=h2

\subsection{第三十章——其次，论婚姻。}

他说：\Quote{我们说，那如今在全地举行的婚姻是天主所设立的，已婚者无罪，但淫乱者和奸淫者应受谴责。}这是真实而公教的道理；但你想由此引申出的，即从男女交合中出生的人不带来任何需由重生的洗涤除去的罪——这是虚假而异端的。\par

% structure: p0063 source=h2 target=subsection reason=relative-heading-rank; base=h2

\subsection{第三十一章——第三，论夫妻交合。}

他说：\Quote{我们说，性冲动——即那没有它便不可能有交合的力量本身——是天主所设立的。}对此我回答说：性冲动，并且用他的话来说，那没有它便不可能有交合的力量，原是由天主如此设立的，以致它没有任何可羞耻之处。因为受造者不宜羞于造物主的工程；而是因正义的惩罚，肢体的悖逆是对第一人悖逆的报应，他们因这悖逆而羞耻，并用无花果叶遮蔽那些先前并不羞耻的羞耻部位。\par

% structure: p0065 source=h2 target=subsection reason=relative-heading-rank; base=h2

\subsection{第三十二章 [XVI]——亚当和厄娃所穿的无花果叶裙。}

因为他们犯罪之后，并没有给自己穿外衣遮盖全身，而是用了围腰，\BibleRef{创 3:7}。有些不太仔细的译者将其译为\Quote{覆盖物}。这固然不错，但\Quote{覆盖物}是一个总称，可以指各种衣服和面纱。应当避免歧义，所以，既然希腊文称它们为\OriginalTerm{围腰}{\GreekText{περιζώματα}}，只遮盖身体羞耻部位，那么拉丁文也应当使用希腊词本身，因为现在习惯已经用它替代了拉丁词；或者，如同某些人那样，使用\Quote{围腰}一词；或者，如同另一些人更好地命名的那样，使用\Quote{角力围腰}。因为这个名称取自一个古罗马习俗：青年人在田野里赤身操练时，用腰带遮盖他们的羞耻部位；因此，直到今日，他们仍被称为\Emph{campestrati}，因为他们用腰带遮盖那些肢体。然而，如果犯罪所用的肢体在犯罪后应当被遮盖，那么人们其实不应该穿外衣，而应该遮盖手和口，因为他们是通过手取和口吃犯了罪。那么，当禁食被取用、诫命的违犯已经发生后，目光转向那些肢体是什么意思呢？那里感觉到什么未知的新奇事物，迫使它被注意？这正由眼睛的睁开所象征。因为\Person{亚当}给牲畜和飞鸟命名时，\Person{厄娃}看见那些树木美丽悦目时，他们的眼睛并不是闭着的；而是被打开——也就是说，专注——去思考；正如关于\Person{撒辣}的婢女\Person{哈加爾}所记载的：她开了眼睛，看见了一口井，\BibleRef{创 21:19}，虽然她先前肯定不是闭着眼睛的。因此，他们如此突然地为自己的赤身感到羞耻——他们每天都习惯于看见这些而不觉混乱——以至于现在再也无法忍受那些肢体裸露，而是立刻设法遮盖它们；难道他们没有——他在明处，她在暗处的冲动——察觉到那些肢体对于他们意志的选择不服从吗？而这些肢体本应像其他肢体一样受他们意志的指挥。他们应受此苦，因为他们自己也没有服从他们的\Person{主}。因此他们羞愧，因为他们如此没有向他们的造物主表示服事，以至于应受失去对那些用以生育儿女的肢体的掌管。\par

% structure: p0067 source=h2 target=subsection reason=relative-heading-rank; base=h2

\subsection{三十三、赤身的羞耻}

这种羞耻——这种必须脸红的——确实是每个人生来就有的，并且在某种程度上是由自然法则本身所规定的；以至于在这方面，即使是有道德的已婚者也会感到羞耻。没有人会走向如此邪恶和可耻的极端，因他知道\Person{天主}是自然的创造者和婚姻的制定者，而在任何人的眼前与妻子行房，或者不因那些冲动而脸红、不寻求隐秘——他不仅可以避开陌生人的目光，甚至还可以避开所有亲戚的目光。因此，让人性承认因自己的过错而临到它身上的恶，免得它被迫要么不因自己的冲动而脸红（这是最无耻的），要么因造物主的作为而脸红（这是最忘恩负义的）。然而，有道德的婚姻却为了生育儿女的好处而善用这种恶。但为了肉身的快感而同意情欲，这本身就是罪，尽管它可能因许可而被允许给已婚者。\par

% structure: p0069 source=h2 target=subsection reason=relative-heading-rank; base=h2

\subsection{34 [XVII]. 堕落之前乐园中能否有感官欲望？}

但是，你们白拉奇派的人啊，在坚持婚姻的尊贵和生育性的同时，请确定：如果没有人犯罪，你们愿意怎样认为那些人在乐园中的生活，并从以下四种情况中选择一种。因为毫无疑问，要么只要他们愿意，就随时行房；要么当不需要行房时，他们就克制情欲；要么情欲会在意志的召唤下兴起，正好在贞洁的谨慎预先察觉到需要行房的时候；要么根本没有情欲存在，如同其他每一肢体为自己的工作服务，生育器官也会毫无困难地服从那些愿意者的命令。从这四个假设中，选择你们喜欢的；但我认为你们会拒绝前两个，在这两个假设中，情欲要么被服从，要么被抵抗。因为第一个假设不符合如此伟大的德行，第二个假设不符合如此巨大的幸福。因为愿这种想法远离我们：如此巨大幸福的荣耀，要么因总是跟从预先的情欲而被最卑贱地奴役，要么因抵抗它而不得享受最丰盛的平安。我说，要抛弃这种想法：即那心灵要么因同意满足肉身的贪欲而满足——这种贪欲不是适时地为生育而兴起，而是带着无节制的兴奋——要么那安宁觉得有必要通过拒绝来约束它。\par

% structure: p0071 source=h2 target=subsection reason=relative-heading-rank; base=h2

\subsection{三十五、乐园中的欲望要么全然没有，要么服从意志的冲动}

但是，若你选择了另外两个选择中的任何一个，却没有必要以任何辩论与你争辩。因为即使你拒绝选择第四个选项（其中所有服从的肢体都享有最高的平安，没有任何情欲），既然你论证的紧迫性已经使你敌视它了；那么，毫无疑问，你会喜欢我放在第三位的那个选项：那肉体的情欲（其冲动达到最终快感，那快感令你非常愉悦）决不会在乐园中出现，除非随意志的命令，且是为生育所必要时。如果你乐意在乐园中安排这样的事，并且藉着这样一种既不会先于、也不会妨碍、更不会超越意志命令的肉体情欲，你认为孩子可以受孕诞生，我并无异议。因为就我对此事而言，我满足于：这样的肉体情欲如今不在人类中间，如同你承认在那福乐之地可能有的那样。因为现在它是什么，所有人的感觉确实承认，尽管是带着羞惭：因为它既以过度和纠缠的不安骚扰贞洁者——甚至当他们不愿意并以节制加以制止时——又常常从愿意者那里退却，并强加于不愿意者；这样，藉着它的不服，它证明自己无非是那最初的不服的惩罚。因此，合理地，当时第一人在遮盖他们的裸体时，以及现在任何自视为人的人（无论是否谦逊，甚至不谦逊之人）都为此感到羞愧——我们决不能说它是天主的作为，而是——那最初古老之罪的罚。然而，你——不是为了宗教的理性，而是为了激动的争辩；不是为着人类的谦逊，而是为你自己的狂妄，为使甚至肉体的情欲本身不被认为是败坏的，并且原罪不是由它衍生而来——正企图用你的论证将它绝对地招回乐园，如同它现在这样；并争辩说那情欲本可以在那里存在，它要么总是继之以可耻的同意，要么有时因可怜的拒绝而被抑制。然而，我对你乐意怎样想它并不十分在意。但是，凡藉着它而生的人，若不再生，无疑将被定罪；并且他必在\Person{魔鬼}的权下，除非\Person{基督}将他从中解救出来。\par

% structure: p0073 source=h2 target=subsection reason=relative-heading-rank; base=h2

\subsection{第三十六章\Person{尤利安}第四项反驳：人是天主的工程，且不受祂的权力强迫为善或为恶。}

\Quote{他说：我们主张，人是天主的工程，没有人被祂的权力强迫不愿意地要么行恶要么行善，而是人出于自己的意志行善或作恶；但在善工中，他常得天主恩宠的助佑，在恶中，他受\Person{魔鬼}的诱惑唆使。}对此，我回答说：人，就其为人而言，是天主的工程；但就其为罪人而言，他们在\Person{魔鬼}权下，除非被那成为天主与人之间的中保者从中拔出，祂成为中保，别无其他原因，只因为祂不能是一个来自罪人的罪人。并且，没有人被天主的权力强迫不愿意地行善或作恶，而是当天主舍弃一个人时，他理应走向恶；当天主助佑时，他毫无功绩地转向善。因为人不是出于不愿意而为善，而是藉天主的恩宠甚至得助到愿意的地步；因为经上记载并非徒然：\BibleQuote{因为是天主在你们内工作，使你们愿意，并使你们力行，为成就祂的善意。}\BibleRef{斐 2:13}以及\BibleQuote{意志是由天主预备的。}\BibleRef{箴 8:35}\par

% structure: p0075 source=h2 target=subsection reason=relative-heading-rank; base=h2

\subsection{第三十七章——善意之始是恩宠的恩赐。}

但是，你们认为人在善行中得到天主恩宠的助佑，以至于你认为在激发他的意志去从事那善行时，恩宠毫无作用；你们自己的话已充分说明了这一点。因为你们为何不说人是由天主恩宠的激励而行善，正如你们说他由魔鬼的暗示而行恶，而却说他在善行中始终由天主恩宠帮助？——好像他凭借自己的意志，毫无天主的恩宠，就承担了一件善行，然后在善行本身中得到天主性的帮助，即是说，为了他的善意的功绩；这样恩宠就作为应得的回报而被赐予——不是作为不配领受的恩赐而被赐予——于是恩宠就不再是恩宠了。见：\BibleRef{罗 11:6} 而这正是\Person{贝拉基}在\Place{巴勒斯坦}的审判中以欺诈的心所谴责的——即天主的恩宠，也就是，是照我们的功绩而赐予的。我恳求你们告诉我，\Person{保禄}（当时还是\Person{扫禄}）曾愿意过什么好事，难道不是大恶吗？当时他口吐杀气，在可怕的心智黑暗和疯狂中，去扫荡基督徒？见：\BibleRef{宗 9:1} 因为他有什么善意的功绩，使天主藉着一次奇妙的突然召唤，将他从那些恶事转变为善事呢？我该说什么呢？当他亲自呼喊道：\BibleQuote{不是出于我们所行的义德，而是照祂的仁慈，祂拯救了我们。}见：\BibleRef{铎 3:5} 那是我已经提到的主所说的话：\BibleQuote{没有人能到我这里来}——这被理解为\Quote{相信我}——除非由我父赐给了他。见：\BibleRef{若 6:66} 这是否是赐给已经愿意相信的人，为了善意的功绩？或者，意志本身，像扫禄的情况那样，是从上而被激发的，使他相信，即使他如此远离信仰甚至迫害信徒？因为主如何命令我们为迫害我们的人祈祷？我们这样祈祷：愿天主的恩宠因他们的善意而报答他们，而不是愿邪恶的意志本身被转变为善良的意志？正如我们相信，当时那些被他迫害的圣人为扫禄的祈祷并非徒然，为使他意志转变为那正在摧毁的信仰。确实，他的皈依由上而来，甚至通过一个明显的奇迹而显现。但在今天，有多少基督的仇敌，被天主隐秘的恩宠突然吸引到基督面前！如果我没有从福音中写下这个字，那人关于这事会怎样说我呢？因为即使现在，他不是在反对我，而是在反对那呼喊的主：\BibleQuote{没有人能到我这里来，除非那派我来的父吸引他！}见：\BibleRef{若 6:44} 因为祂不是说\Quote{除非祂引导他}，这样我们就能够以某种方式理解他的意志居先。因为谁是\Quote{被吸引}的，如果他已经是愿意的？然而，没有人来，除非他是愿意的。因此，他以奇妙的方式被吸引去愿意，由那知道如何在人内心工作的主。不是要那些不愿意的人相信——那是不可能的——而是要他们从不情愿变成情愿。\par

% structure: p0077 source=h2 target=subsection reason=relative-heading-rank; base=h2

\subsection{第三十八章【二十】——天主恩宠的德能被证明}

我们不是凭人的猜测来推断这事的真实性，而是凭神圣圣经最清晰的权威来辨明。在\BibleBook{编年纪}{下}中读到：\BibleQuote{那时，天主的手也感动犹大人，使他们同心合意，遵行王与众首领照着主的命所发布的命令。}见：\BibleRef{编下 30:12} 此外，主藉厄则克耳先知说：\BibleQuote{我要赐给你们一颗新心，在你们内注入一种新的精神；我要从你们的肉体内取出石心，给你们换上一颗血肉的心，好使你们遵行我的诫命，遵守我的典章，并实行它们。}见：\BibleRef{则 36:26} 艾斯德尔王后在祈祷中说的又是什么呢？她说：\BibleQuote{求赐我口舌伶俐，使我的言语在狮子面前有光彩，转变他的心去仇恨那攻击我们的人。}见：\BibleRef{艾 14:13} 她为何向天主如此祈祷呢，如果天主不在人的心中动工？但或许这妇人在这样祈祷时是愚昧的。那么，让我们看看，这祈求者的愿望是否白白地预先表达出来，而结果并没有像蒙垂听者那样发生。看，她来到国王面前。我们毋庸多言。因为她没有按自己的次序接近他，迫于她的巨大急需，正如经上所记载的：\Quote{他看她，像一只在暴怒冲动中的公牛。} \Quote{王后害怕了，她的脸色因昏厥而改变，她俯身在她前面的使女头上。天主转变了他，将他的愤怒转化为温和。}现在何必叙述其后之事？神圣圣经在那里证实，天主完成了她所求的，藉着在国王心中动工，使他产生祂所命令的意志，于是事情就按王后所求的成就了。天主早已垂听了她，使此事成就，祂在国王尚未听到恳求他的妇人的陈词之前，就以最隐秘而有效的力量改变了国王的心，将他的愤怒转化为温和——也就是从加害的意志转变成施恩的意志——正如宗徒的话：\BibleQuote{天主在你们内运作，使你们愿意。}写这些事的天主的人——不，是引导他们写下这些事的天主圣神自己——是否攻击了人的自由意志？绝无此事！而是祂在一切事上称赞了全能者最公义的审判和最慈悲的帮助。因为人知道在天主面前没有不义就够了。但祂如何分配这些恩赐，使一些人理应成为义怒之器，另一些人蒙恩成为仁慈之器——谁能识透主的心意，或谁曾作过祂的参谋？如果我们获得了恩宠的光荣，就不要将所领受的归于自己而忘恩负义。\BibleQuote{因为我们有什么不是领受的呢？}见：\BibleRef{格前 4:7}\par

% structure: p0079 source=h2 target=subsection reason=relative-heading-rank; base=h2

\subsection{第三十九章【二十一】——\Person{儒利安}关于旧约圣人们的第五项异议}

他说：\Quote{我们说，}他说，\Quote{旧约的圣人，在此地达成圆满的义德后，就进入永生——也就是说，因着对德行的爱，他们远离了一切罪恶；因为我们读到那些曾犯过罪的人，我们知道他们后来改过了。}无论你宣称古代义人拥有何种德行，唯有相信那位为赦免他们的罪而流血的中保才能拯救他们。因为他们自己的话是：\Quote{我信了，所以我说话。}由此宗徒保禄也说：\Quote{我们既然具有同样信德的精神，正如经上所载：\InnerQuote{我信了，所以我说}，我们也信，所以也说。}\BibleRef{格后 4:13} 什么是\Quote{同样的精神}，不就是这些说了此话的圣人所拥有的那一位圣神吗？宗徒伯多禄也说：\Quote{你们为什么愿意把重担放在外邦人的颈上，这个重担我们和我们的祖先都承受不了？但是，我们相信，藉着主耶稣基督的恩宠，我们得救，正如他们一样。}\BibleRef{宗 15:10} 你们这些恩宠的敌人不愿意相信古人也同样因耶稣基督的恩宠得救，却按照白拉奇划分时期，在他的著作中读到这一点：你们说，在法律以前，人凭自然得救；然后凭法律；最后凭基督，好像前两个时期（法律以前和法律之下）的人不需要基督的血一样，使这话落空：\BibleQuote{因为只有一个天主，也只有一个在天主与人之间的中保，就是人基督耶稣。}\BibleRef{弟前 2:5}\par

% structure: p0081 source=h2 target=subsection reason=relative-heading-rank; base=h2

\subsection{第四十章 [第二十二节]——第六项反对意见，论恩宠为众人之必要，及论婴儿圣洗。}

他们说：\Quote{我们承认基督的恩宠对所有人都是必要的，包括成人和婴儿；并且我们谴责那些说由两个已受洗者所生的婴儿不应该受洗的人。}我知道你说这样的话是出于什么意义——不是按照宗徒保禄，而是按照异端者白拉奇；也就是说，圣洗对婴儿是必要的，不是为赦罪，而仅仅是为天国；因为你们给予他们，即使未受洗，在天国之外也有救恩和永生的地方。你们也不考虑经上所记载的：\BibleQuote{信而受洗的必要得救；但不信的必被判罪。}\BibleRef{谷 16:16} 因此，在救主的教会中，婴儿藉着他人相信，正如他们藉着他人领受了在圣洗中赦免的那些罪一样。你们也不认为那些没有基督身体和血的人不能有生命，虽然祂亲自说过：\BibleQuote{你们若不吃我的肉，不喝我的血，在你们内便没有生命。}\BibleRef{若 6:34} 或者，如果你们被福音的话所迫，承认离开肉身的婴儿除非受洗，否则不能有生命或救恩，那么问问为什么那些未受洗的人，在判定人无罪的审判者面前，必须承受第二次死亡的审判？你们就会找到你们不想要的东西——原罪！\par

% structure: p0083 source=h2 target=subsection reason=relative-heading-rank; base=h2

\subsection{第四十一章 [第二十三节]——第七项反对意见，论圣洗的效力。}

他说：\Quote{我们谴责那些声称圣洗并没有除去一切罪恶的人，因为我们知道藉着这些奥迹，获得完全的洁净。}我们也这样说；但是你们不说婴儿也藉着同样的奥迹，从他们初生和可憎的传承的束缚中得以释放。因此，你们像其他异端者一样，应当从基督的教会中分离出来，教会自古以来就持守这一点。\par

% structure: p0085 source=h2 target=subsection reason=relative-heading-rank; base=h2

\subsection{第四十二章 [第二十四节]——他驳斥儒利安书信的结论。}

但现在他结束书信的方式是说道：\Quote{因此，不要让人迷惑你们，也不要让恶人否认他们有这样的想法。但如果他们说实话，要么给他们一个听审的机会，要么让那些现在不同意我的主教们，谴责我上面说过的他们与摩尼教徒共同持有的观点，如同我们谴责他们关于我们的言论一样，这样就能达成完全的和解；但如果他们不愿意，要知道他们是摩尼教徒，要远离他们的团体。}——这话与其说是应受谴责，不如说是应受轻视。因为我们中谁还会犹豫要咒骂摩尼教徒呢？他们说不从善天主生出人，没有设立婚姻，也没有赐给法律，这法律是藉着梅瑟传给希伯来人的。但反对白拉奇主义者，我们也不是无缘无故地咒骂他们，因为他们如此敌视天主的恩宠——这恩宠是藉着我们的主耶稣基督而来的——以至于说这恩宠不是白白的赐予，而是按照我们的功过赐予，这样恩宠就不再是恩宠了；\BibleRef{罗 11:6} 并且如此高抬自由意志（人正是藉此坠入深渊），以至于说因为善用自由意志，人就配得恩宠——然而没有人能善用自由意志，除非藉着恩宠，这恩宠不是按债务偿还，而是由天主的仁慈白白赐予的。并且他们如此争辩说婴儿已经得救了，以至于他们敢否认婴儿需要被救主拯救。并且，持有并散布这些可憎的教义，他们还不断地要求听审，而他们作为被定罪者，本应悔改。\par

\SourceRecord{Translated by Peter Holmes and Robert Ernest Wallis, and revised by Benjamin B. Warfield.；Nicene and Post-Nicene Fathers, First Series；Vol. 5.；Edited by Philip Schaff.；Buffalo, NY: Christian Literature Publishing Co.,；1887.；<http://www.newadvent.org/fathers/15091.htm>.}

% structure: p0001 source=h1 target=section reason=page-title

\section{\WorkTitle{驳斥佩拉奇亚主义的两封信（第二卷）}}

他着手审查佩拉奇亚主义者的第二封信，该信如同第一封一样充满对公教徒的诽谤——这封信是由他们以十八位主教的名义寄给\Place{得撒洛尼}的；首先，他通过比较异端著作彼此之间的异同，表明公教徒在厌恶佩拉奇亚主义者的教条时，绝没有陷入摩尼教徒的错误。他驳斥了关于罗马圣职人员在\Person{佐西默}后来谴责\Person{佩拉奇乌斯}和\Person{塞莱斯提乌斯}一事中犯有失职之罪的诽谤，表明佩拉奇亚主义的教条在罗马从未被认可，尽管有一段时间，由于\Person{佐西默}的宽大，\Person{塞莱斯提乌斯}得到了仁慈的对待，以期引导他改正错误。他表明，公教徒在恩宠的名义下，既不主张命运的教义，也不将偏袒归于天主；尽管他们确实说天主的恩宠不是按照人的功过赐予的，而且人最初的向善意愿是由天主激发的；因此，一个人若没有天主白白赐予的恩宠的推动，就根本不能开始从恶到善的改变。\par

% structure: p0003 source=h2 target=subsection reason=relative-heading-rank; base=h2

\subsection{第一章 引言：佩拉奇亚主义者指控公教徒为摩尼教徒。}

现在让我考虑第二封信，不仅是\Person{朱利安}一人所写，而是与他同属几位主教的共同信件，他们将其寄给了\Place{得撒洛尼}；让我在天主的帮助下尽我所能地回答它。以免我的这篇作品变得比题目本身的需要更长，何需反驳那些不包含他们教义隐伏之毒的东西，那些似乎只是为求得东方主教们的默许或帮助，或者是为了公教信仰而在他们所谓的反对摩尼教徒的亵渎之事的请愿呢？他们的目的无非是，在提出一种可怕的异端并自称是它的反对者时，以赞美本性为掩护，暗藏为恩宠的敌人。因为何时曾有人针对他们提出这些问题的争论？或者，有哪个公教徒不因他们谴责那些宗徒所预言将要离弃信仰、良心被烙伤、禁止结婚、禁食他们认为不洁净的食物、不认为万物由天主所造的人而感到不快呢？何时曾强迫他们否认天主的一切受造物都是好的，并且除了那非受造的天主自己以外，没有任何实体不是至上天主所造的？这些明显是公教真理的东西，并没有在他们的谴责和定罪中受到指责；因为不仅公教信仰极其痛恨摩尼教徒的愚蠢而有害的亵渎，而且所有非摩尼教徒的异端也痛恨它。因此，这些佩拉奇亚主义者也适当地诅咒摩尼教徒，并谴责他们的错误。但他们做了两件坏事情，为此他们自己也应受到诅咒：一是他们以摩尼教徒的名义诬告公教徒，二是他们自己也正在引入一种新错误的异端。因为他们并非因为没有染上摩尼教徒的疾病而信仰纯正。瘟疫的种类并不总是一个样——在身体上如此，在思想上也是如此。因此，正如身体的医生不会认为一个他宣布没有水肿的人就没有死亡的危险，如果他看到他患上另一种致命的疾病；同样，他们并不因为不是摩尼教徒就被承认真理，如果他们在另一种乖谬中癫狂的话。因此，我们与他们共同憎恶的东西是一回事，我们在他们身上憎恶的东西是另一回事。因为我们在厌恶他们时，我们也是厌恶他们所厌恶的东西；然而我们在他们身上憎恶的，正是使他们自身成为可憎的那些东西。\par

% structure: p0005 source=h2 target=subsection reason=relative-heading-rank; base=h2

\subsection{第二章 摩尼教与佩拉奇亚主义的异端彼此对立，同受天主教会的谴责。}

摩尼教徒说，善的天主并非所有本性的创造者；佩拉奇亚主义者则说，天主并不是各时代人类中所有人的净化者、救主和拯救者。公教会谴责两者，既维护天主创造万物以反对摩尼教徒，以免否认任何本性是由他所造；又维护所有时代的人性必须被视为败坏并需要寻求拯救以反对佩拉奇亚主义者。摩尼教徒斥责肉身的贪欲，认为它不是一个偶然的恶习，而是一种从永恒起就是邪恶的本性；佩拉奇亚主义者却称赞它，认为它根本不是恶习，甚至是自然的美善。公教信仰谴责两者，对摩尼教徒说：\Quote{这不是本性，而是恶习；}对佩拉奇亚主义者说：\Quote{这不是来自父，而是来自世界；}这样两者都承认它是一种需要医治的罪恶疾病——前者停止相信它是不可治愈的，后者停止赞美它。摩尼教徒否认善人的恶始于自由意志；佩拉奇亚主义者则说，即使坏人也有足够的自由意志去行善的命令。公教会谴责两者，对前者说：\BibleQuote{天主造人原很正直}\BibleRef{训 7:30}，对后者说：\BibleQuote{如果子使你们自由，你们就真自由了}\BibleRef{若 8:36}。摩尼教徒说，灵魂作为天主的一粒，因邪恶本性的混合而犯罪；佩拉奇亚主义者说，灵魂是正直的，虽然不是一粒，而是天主的受造物，并且即使在今生腐朽的生命中也没有任何罪。公教会谴责两者，对摩尼教徒说：\BibleQuote{你们或使树好，它的果子也好；或使树坏，它的果子也坏}\BibleRef{玛 12:33}，这话不会对不能改变自己本性的人说，除非因为罪不是本性，而是恶习；对佩拉奇亚主义者说：\BibleQuote{如果我们说我们没有罪，便是欺骗自己，真理也不在我们内}\BibleRef{若一 1:8}。在这些彼此对立的疾病中，摩尼教徒和佩拉奇亚主义者发生分歧，怀着不同的意愿但相似的虚荣，因不同的意见而分离，但因悖逆的心态而接近。\par

% structure: p0007 source=h2 target=subsection reason=relative-heading-rank; base=h2

\subsection{第三章 摩尼派与佩拉奇亚派在错谬上的联合与分离程度}

然而，他们确实同样反对基督的恩宠，同样使祂的圣洗无效，同样羞辱祂的肉身；但是，他们以不同的方式和不同的理由行这些事。因为摩尼派断言，天主的助佑是赐给善意本性的功德，而佩拉奇亚派则断言，是赐给善意意志的功德。前者说，天主因祂肢体们的劳力而欠此恩惠；后者说，天主因祂仆人们的德行而欠此恩惠。因此，在两种情况下，赏报都不是按恩宠，而是按债务来归算的。摩尼派以亵渎之心主张，重生的洗涤——即水本身——是多余的，毫无益处。但佩拉奇亚派则断言，在圣洗圣事中为除去罪过所说的，对婴儿毫无效力，因为他们没有罪；因此，在婴儿的圣洗中，就赦罪而言，摩尼派毁掉了有形元素，但佩拉奇亚派连那无形的圣事也毁掉了。摩尼派亵渎基督由童贞女诞生而羞辱祂的肉身；但佩拉奇亚派则使那些应得救赎之人的肉身与救赎者的肉身相等，从而羞辱了祂的肉身。因为基督固然不是在有罪的肉身中出生的，而是在有罪肉身的相似形中出生的，而其余人类的肉身都是在有罪中出生的。因此，摩尼派痛恨所有肉身，就从基督的肉身中夺走了明显的真理；而佩拉奇亚派则主张没有肉身是在有罪中出生的，就从基督的肉身中夺走了其独特和应有的尊严。\par

% structure: p0009 source=h2 target=subsection reason=relative-heading-rank; base=h2

\subsection{第四章 两种相反的错谬}

那么，让佩拉奇亚派不要再向公教徒提出他们自己所不是的指控，而要赶快改正他们自己的所是；并且不要希望因反对摩尼派可憎的错谬而被认为值得赞许，而要承认自己因没有除去自己的错谬而应受憎恨。因为两种错谬可能相互对立，尽管两者都应受谴责，因为它们同样与真理对立。因为如果佩拉奇亚派因憎恨摩尼派而应被爱，那么摩尼派也因憎恨佩拉奇亚派而应被爱。但我们的公教慈母断不可因某人憎恨他人而选择爱某些人，因为藉着主的警告和帮助，她应该避开两者，并渴望医治两者。\par

% structure: p0011 source=h2 target=subsection reason=relative-heading-rank; base=h2

\subsection{第五章 [III.] 佩拉奇亚派对罗马教会圣职人员的诽谤}

此外，他们指控罗马圣职界，写道：\Quote{他们因畏惧命令的驱使，竟不知羞耻地犯下背信的罪行；以致他们违背了先前的判决（在其中他们曾通过行动赞同公教教义），后来竟宣称人的本性是邪恶的。} 然而，白拉奇派却怀着一个虚假的希望，认为白拉奇或塞莱斯提乌斯那个新而可憎的教义能够得到某些罗马人的公教智识的接纳，当时那些狡猾的人物——虽然被邪恶的错误所扭曲，却并非可鄙之人，因为他们似乎更应当得到审慎的纠正，而非轻易的责罚——受到了比教会更严格的纪律所要求的更多的宽大处理。因为当如此众多且如此重要的教会文件在宗座与非洲主教之间往返传递时——而且，当在此事上的诉讼在那个宗座内完成，塞莱斯提乌斯在场并作出答辩时——可敬的教宗佐西默究竟有什么书信、什么法令，规定必须相信人生来没有任何原罪的玷污？他绝对从未说过这话——从未写过它。但因塞莱斯提乌斯曾在他的小册子中写了这些，仅作为他承认自己仍存疑惑且愿受教诲的事项之一，他这样一个敏锐之人改过的意愿（若他能被纠正，必将有益于多人）得到了认可，而非教义的错误。因此他的小册子被称为公教的，因为这也是公教的态度的一部分——若偶然在某事上人思想与真理要求不同，不是最精确地定义那些事，而是一旦被察觉和证明，就抛弃它们。因为宗徒说话时不是对异端者，而是对公教徒，他说：\BibleQuote{所以，我们凡是成全的人，都应怀有这种心意；若你们在什么事上另有心意，天主也要将这事启示给你们。} \BibleRef{斐 3:15} 人们认为这似乎在他身上发生了，当他回答说他同意荣福教宗依诺森的书信时，那些书信中关于此事的全部疑虑都被消除了。为了使这一点在他身上更充分、更明显，事情被推迟，直到来自非洲的书信到达——在那个省份，他的狡诈在某种程度上已更加为人所知。随后这些书信到达罗马，内容如下：对于较为迟钝和焦虑的人而言，他承认自己普遍同意依诺森主教的书信并不足够，而应当公开咒诅他在小册子中所作的有害言论；以免他若不这样做，许多见识更高的人反倒会相信，他的小册子中那些信仰的毒药已被公教宗座所认可，因为该宗座曾肯定那本小册子是公教的，而非相信它们因他回答同意教宗依诺森的书信而得到纠正。因此，当要求他到场，以便通过确定而清楚的回答使此人的狡诈或纠正明确显现、不为任何人所怀疑时，他却自行退避，拒绝接受审问。先前已经为他人益处而作出的延迟，若不能有益于那些极其乖张之人的固执和疯狂，也不会发生。但若（绝不可能如此）在罗马教会中关于塞莱斯提乌斯或白拉奇竟如此判断，以致他们那些教义（教宗依诺森亲自并连同他们一起谴责过）竟被宣布值得认可和维护，那么背信的标记就不得不因此打在罗马圣职界身上。然而如今，当最荣福的教宗依诺森的第一封书信（回复非洲主教的信）同样谴责了这些人试图向我们推荐的错误时；他的继任者，圣教宗佐西默，从未说过、从未写过应当持守这些人关于婴儿所想的教义；不仅如此，当塞莱斯提乌斯试图为自己开脱时，他甚至用重复的判决束缚塞莱斯提乌斯，要他同意上述宗座的书信——肯定地，在此期间关于塞莱斯提乌斯所做的一切较为宽大之事，只要保持了最古老而坚强的信德的稳定，就是最仁慈的纠正规劝，而非最有害的邪恶认可；而后来，由同一司祭职以重复的权威谴责塞莱斯提乌斯和白拉奇，则是对一种稍作间歇、最终必须施行的严厉的证明，而非对先前已知真理的否认或对真理的新认可。\par

% structure: p0013 source=h2 target=subsection reason=relative-heading-rank; base=h2

\subsection{第六章 [IV.]——关于塞莱斯提乌斯和佐西默所行之事}

然而，在这件事上，我们何须再多费唇舌？各地已有公开的记录和文件，其中记载了所有经过，可供人了解或回忆。谁看不出塞勒斯提乌斯是如何被你神圣前任的审问和塞勒斯提乌斯的回答所束缚？他当时宣称同意教宗依诺增爵的信函，并被一条最有益的锁链牢牢拴住，以致不敢再坚持婴儿的原罪在圣洗中不得除去。因为可敬的主教依诺增爵关于此事写给迦太基会议的信中有这样的话：\Quote{他原曾拥有自由意志，但轻率地使用其优势，堕入背信的深渊，被淹没而不得起身，永远被自己的自由所欺骗，本将一直受此毁灭的压迫，若非后来基督的来临以祂的恩宠解救了他，并藉着新重生的净化，以祂圣洗的洗涤洗净了所有过去的罪。}还有什么比这宗座判决更清楚、更显明的呢？塞勒斯提乌斯对此表示同意，当时你神圣前任对他说：\Quote{你谴责所有以你名义散布的言论吗？}他回答：\Quote{我依照你前任真福依诺增爵的判决谴责它们。}但在他名下所传的言论中，执事保利努斯曾向塞勒斯提乌斯提出反对，说他声称\Quote{亚当的罪只损害他自己，而非人类；新生婴儿所处的状态与亚当犯罪前相同。}因此，倘若他真心诚意地依照真福教宗依诺增爵的判决谴责保利努斯所提出的观点，那么此后他还能留下什么依据来争辩说婴儿身上没有从第一人从前过犯而来的罪，而这罪要在圣洗中藉新重生的净化得以洗净呢？然而，最终的结果表明他回答得虚假，因为他退出了审讯，以免被迫按非洲谕令明确提及并诅咒他就此问题在其论著中所写的话语本身。\par

% structure: p0015 source=h2 target=subsection reason=relative-heading-rank; base=h2

\subsection{第七章——他向塞勒斯提乌斯提出一个两难问题。}

关于同一案件，同一教宗回复努米底亚主教们的话又是什么呢？因为他收到了两个会议——迦太基会议和米莱维会议的来信。他难道不是极其清楚地谈论婴儿吗？因为他的话是这样：\Quote{因为你们的友爱所宣称的，即婴儿即使没有圣洗的恩宠也能获得永生的赏报，这是极其愚蠢的；除非他们吃人子的肉，喝他的血，他们便没有生命在自己内。\BibleRef{若 6:54}而那些人认为这是属于他们而不需重生，在我看来是想摧毁圣洗本身，因为他们宣称这些婴儿拥有我们所相信不经圣洗就无法赐予的东西。}对此，这个忘恩负义的人怎么说呢？当时宗座已因他的宣称而宽恕了他，仿佛他已被其最仁慈的宽厚所纠正。他对此怎么说？婴儿在生命终结后，即使生前未在基督内受洗，会是在永生中，还是不会？如果他说：\Quote{他们会在}，那么他如何回答说他已\Quote{依照真福依诺增爵的判决}谴责了以他名义所传的言论呢？看，真福教宗依诺增爵说，婴儿没有基督的圣洗和分享基督的体血就没有生命。如果他说：\Quote{他们不会}，那么，如果他们没有获得永生，岂不是理所当然地被定罪永死，如果他们没有原罪的话？\par

% structure: p0017 source=h2 target=subsection reason=relative-heading-rank; base=h2

\subsection{第八章——关于婴儿的公教信仰。}

那些胆敢写下邪恶不敬之言并将它们寄给东方主教的人，对这一切有何话说？塞勒斯提乌斯被认为已同意可敬的依诺增爵的信函；前述教长的信函被宣读，他写道，未受洗的婴儿不能有生命。那么，谁能否认，既然他们没有生命，随之而来他们就有死亡？那么，婴儿何以遭受如此悲惨的惩罚，如果没有原罪？那么，那些背弃信仰、反对恩宠的人，在若西默主教手下，为何又指控罗马圣职人员背信，仿佛他们在后来谴责塞勒斯提乌斯和贝拉基时，与先前在依诺增爵手下时曾持有不同的看法？因为，既然通过可敬的依诺增爵关于婴儿除非在基督内受洗否则永死的信函，公教信仰的古道彰显无遗，那么，偏离那判决的人，才会成为罗马教会的背信者；幸亏天主保佑，这事并未发生，而是在一再谴责塞勒斯提乌斯和贝拉基时，那判决本身被坚持到底。愿他们明白自己正陷于他们所指控别人的处境，并愿他们今后从背信信仰中痊愈。因为公教信仰并不说人的本性是恶的——即人最初被造主造为人的那一部分；也不是说天主在那本性中创造的东西——即当祂从人产生人时——是恶的；而是说人从第一人的那罪所继承的东西是恶的。\par

% structure: p0019 source=h2 target=subsection reason=relative-heading-rank; base=h2

\subsection{第九章【第五卷】——他回应贝拉基派的诽谤。}

如今我们必须检视他们在信中所指控我们、并简要提及的那些事。对此，我的答复如下：我们并不是说，因亚当之罪，自由意志从人的本性中消失了；而是说，在受制于魔鬼的人身上，自由意志仅能用于犯罪；而对于良善圣洁的生活，除非人的意志本身借天主的恩宠得以释放，并在行动、言语、思想的每一善行上获得助佑，否则自由意志便无能为力。我们说，除上主天主外，无人是出生者的创造者，婚姻不是由魔鬼设立的，而是由天主亲自设立的；然而，因传生之过，所有人都生于罪下，因此所有人都处于魔鬼的权下，直到他们在基督内重生。我们也并非以恩宠之名维护命运，因为我们说天主的恩宠不以任何人的功绩为先导。不过，若有人愿意称全能天主的旨意为\Quote{命运}，我们固然回避不敬的新奇词语，却也无意在言辞上争论。\par

% structure: p0021 source=h2 target=subsection reason=relative-heading-rank; base=h2

\subsection{第十章——为何佩拉纠派假借恩宠之名，诬告天主教徒主张命运。}

但当我更为仔细地思考，他们究竟有何理由认为应当指控我们以恩宠之名主张命运时，我首先查看了他们随后所说的话。因为他们认为应当这样指控我们：\Quote{他们以恩宠之名，}他们说，\Quote{如此主张命运，竟说除非天主将行善的愿望——且是不完美的善——注入不情愿、抗拒的人心中，否则他既不能远离恶，也不能持守善。}之后稍隔一段，在他们陈述自己所坚持的观点时，我留意了他们就此所说的话。\Quote{我们承认，}他们说，\Quote{圣洗对各年龄段都是必要的，此外恩宠还协助每个人的良好意愿；但恩宠并不将对德行的爱注入不情愿的人心中，因为天主不偏待人。}从他们这些话中，我看出他们之所以认为或希望别人认为我们以恩宠之名主张命运，是因为我们说天主的恩宠不是因我们的功绩赐予，而是按祂自己最仁慈的旨意，如祂所说：\Quote{我要恩待谁，就恩待谁；我要怜悯谁，就怜悯谁。}接着自然又加上：\Quote{由此可见，不在于人愿意，也不在于人奔跑，而在于天主施予怜悯。}\BibleRef{罗 9:16}在此，任何人若认为或说宗徒是命运的维护者，同样是愚昧的。但在此，这些人充分暴露了自己；因为他们诽谤我们，说我们以恩宠之名主张命运，因为我们说天主的恩宠不是因我们的功绩赐予，这无疑表明他们自己说恩宠是因我们的功绩赐予；如此，他们的盲目无法掩盖并隐瞒他们所相信和所想的，尽管当这一观点被指控时，佩拉纠在巴勒斯坦的主教审判中，以狡猾的恐惧谴责了它。因为确实，从他自己的弟子塞莱斯提乌斯的话中，他被指控他自己也惯于说天主的恩宠是因人的功绩赐予的。他对此表示厌恶（或假装厌恶），至少立即从口头上予以绝罚；但正如他后来的著作所表明的，以及他的追随者的主张所显明的，他将它藏匿于诡诈的心底，直到日后胆量得以在书信中发表，而当时因恐惧而否认的诡计所隐藏的东西。然而，佩拉纠派的主教们并不畏惧，至少不羞于将其书信寄给天主教的东方主教们，在信中指控我们是命运的维护者，因为我们甚至不说恩宠是按我们的功绩赐予的；尽管佩拉纠因畏惧东方主教们，不敢这样说，因而被迫谴责了它。\par

% structure: p0023 source=h2 target=subsection reason=relative-heading-rank; base=h2

\subsection{第十一章[VI]——将命运的指控掷回给对手。}

然而，骄傲之子、天主恩宠的仇敌、新兴的白拉奇异端者啊，难道果真如此：凡说人的一切善功皆由天主的恩宠所先导，且天主的恩宠不是因功绩而赐予，以免恩宠不再成为恩宠，若不是白白赐予，而是当作配得者的债务而偿还的，这样的人在你们看来，就是在主张\OriginalTerm{命运}{fate}？难道你们自己不也说——无论你们的意图如何——圣洗为所有年龄的人都是必要的？你们不也在同一封信中，将关于圣洗的意见和关于恩宠的意见并列写下了吗？为何圣洗（它赐给婴孩）没有通过这种并列本身提醒你们应当如何看待恩宠？因为你们的原话是：\Quote{我们承认圣洗为所有年龄的人是必要的，而且恩宠也帮助每个人的善志；但它并不将爱德之德注入一个不情愿的人，因为天主并不\OriginalTerm{人格偏袒}{acceptance of persons}。}在你们的这些话中，我暂时不对你们关于恩宠所说的作任何评论。但请就圣洗说明原因，为何你们说它为所有年龄的人是必要的；请说明为何它对婴孩是必要的。确实，因为它赋予他们某种善；而这一善既非微小也非平常，而是意义重大。因为尽管你们否认婴孩沾染原罪（它在圣洗中被赦免），但你们并不否认在那\Quote{重生之洗}中，他们从人子被收纳为天主子；你们甚至宣讲这一点。那么告诉我们，那些在基督内受圣洗并离世的婴孩，无论他们是谁，是如何领受如此崇高的恩赐的？又是凭何种先行的功绩？如果你们说，他们因父母的虔诚而配得此恩，那么就会有人反问你们：为何有时虔诚之人的子女被剥夺了这恩惠，而恶人的子女却获得了它？因为有时从宗教人士所生的后代，在年幼时、刚出母胎，在未及接受重生之洗之前就被死亡抢先夺去，而基督仇敌所生的婴孩却因基督徒的怜悯而在基督内受圣洗——已受圣洗的母亲哀悼自己未受圣洗的孩子，而贞洁的童女将不贞之母所遗弃的陌生婴孩领来受圣洗。在此，父母的功绩显然不存在，而按你们自己的承认，婴孩本身的功绩也不存在。因为我们知道，你们并不相信人的灵魂在住进此肉身之前曾在某处生活过，并行了善或恶的事，以致配得在肉身中如此不同的际遇。那么，什么原因使这个婴孩获得了圣洗，而那个婴孩被拒绝了？难道他们因没有功绩就有了\OriginalTerm{命运}{fate}？还是在这事上，天主有\OriginalTerm{人格偏袒}{acceptance of persons}？因为你们两者都说过——先是命运，后是人格偏袒——既然两者都必须被驳斥，便可留下你们想要引入来对抗恩宠的功绩。那么，请回答关于婴孩功绩的问题：为何有些婴孩在受圣洗后离世，有些未受圣洗离世，而他们既不因父母的功绩拥有如此卓越的恩赐（即从人子成为天主子）也不因之缺失，既无父母的功绩，亦无自己的功绩？你们当然沉默，并且发现自己恰恰处于你们所归咎于我们的相同境地。因为如果没有功绩，你们便说随之而来的是命运，并因此希望人的功绩被理解为存在于天主的恩宠中，以免被迫承认命运；看，你们反倒在婴孩圣洗中更主张了命运，因为你们承认在婴孩中没有功绩。但如果，关于要受圣洗的婴孩，你们否认有任何功绩先行，却又不承认有命运，那么当我们说天主的恩宠是白白赐予的，以免它不再是恩宠，并且不是当作先行的功绩所应得的债务而偿还时，你们为何尖叫说我们是主张\OriginalTerm{命运}{fate}的人？——你们没有看出，在罪人的成义中，既然没有功绩（因为这是天主的恩宠），所以它不是命运（因为这是天主的恩宠），并且它不是人格偏袒（因为这是天主的恩宠）吗？\par

% structure: p0025 source=h2 target=subsection reason=relative-heading-rank; base=h2

\subsection{第十二章 论\Quote{命运}一名之下所指为何。}

因为主张\OriginalTerm{命运}{fate}的人断言，不仅行动和事件，而且我们的意志本身，都取决于人受孕或出生时的星辰位置；这些位置他们称为\OriginalTerm{星座}{constellations}。但天主的恩宠超乎一切星辰、一切天穹，更超乎所有天使。简言之，主张\OriginalTerm{命运}{fate}的人将人的善恶行为及祸福都归于命运；但天主在人的不幸中，以应得的报应追讨他们的功绩，而在幸运中，则以慈悲的旨意白白赐予不配得的恩宠；两者皆不依星辰的时间性排列，而是依祂的严厉与良善的永恒而高深的计划。我们因此看到，两者均不属于\OriginalTerm{命运}{fate}。在此，如果你们回答说，天主的这种仁慈（祂不追随功绩，而是以白白的慷慨赐予不配得的恩惠）倒应称为\OriginalTerm{命运}{fate}，然而宗徒却称之为恩宠，说：\BibleQuote{你们得救是由于恩宠，藉着信德，而且这并不是出于你们自己，而是天主的恩赐；不是出于功行，免得有人自夸，}\BibleRef{弗 2:8}——难道你们不思考、不察觉，主张\OriginalTerm{命运}{fate}的并非我们（以恩宠之名），反倒是你们将以恩宠为名的神圣恩宠称为\OriginalTerm{命运}{fate}吗？\par

% structure: p0027 source=h2 target=subsection reason=relative-heading-rank; base=h2

\subsection{第十三章 [VII.]——他驳斥关于人格偏袒的诽谤。}

而且，我们正确地称它为\Quote{偏待人}，当审判者忽略他所审理的案件的功过，而偏袒一方反对另一方，因为他在那人的身上找到某种值得尊敬或怜悯的东西。但如果有人有两个债务人，他选择免除一人的债务，而要求另一人偿还，他给予他所愿意的人，并不欺骗任何人；这也不应称为\Quote{偏待人}，因为没有不义。偏待人可能在那些见识浅薄的人看来是另一回事，例如葡萄园的主人给只工作了一小时的工人与那些承担了整天重担和炎热的工人同样的工资，使工作时间如此不同的工人得到相等的报酬。但那些向家主抱怨这件事似乎有偏待人的工人们，家主是如何回答的呢？他说：\Quote{朋友，我没有亏负你。你岂不是同我议定了一个银币吗？拿你的走吧；但我愿意给这最后来的，如同给你一样。难道我不可做我所愿意的事吗？难道因为我善良，你就眼红吗？}这里，确实就是全部的公义：\Quote{我愿意这样。对你说，}他说，\Quote{我给了你应得的；我白给了他；我并没有从你那里拿走什么来给他；我没有减少或拒绝我亏欠你的。}\Quote{难道我不能做我所愿意的事吗？难道因为我善良，你就眼红吗？}因此，正如这里没有偏待人，因为一人白白地得到尊重，而另一人并没有被亏负他应得的；同样，按照天主的旨意，一人被召叫，另一人不被召叫，被召叫的那人得到了一种无偿的恩惠，这恩惠的起始就是召叫本身；没有被召叫的那人则得到恶的报应，因为所有人都是有罪的，这是由于罪恶藉着一个人进入了世界。在关于工人的比喻中，确实，只工作了一小时的工人和那些工作时间长十二倍的工人都得到了一个银币——尽管按照人的推理，无论多么虚妄，后者按其劳动量本应得到十二个银币——他们在恩惠上被平等对待，而不是一些人得救，另一些人被定罪；因为那些工作更久的人，是从家主那里得到一切的：他们被召叫而来，并被养活而无缺乏。但那里说：\BibleQuote{所以，祂愿意怜悯谁，就怜悯谁；愿意使谁硬心，就使谁硬心} \BibleRef{罗 9:18} \BibleQuote{祂造了一个器皿为尊贵，又造了一个器皿为卑贱} \BibleRef{罗 9:21}，这实在是无功德而白给的，因为那不被给予的人也是出于同一团泥；但恶是应得地、作为债而被报应的，因为在毁灭的团块中，恶并不被不义地报应在恶人身上。而对于那被报应的人，这恶是恶，因为那是他的惩罚；而对于那施行报应者，这是善，因为那是祂的权利。在两个同样有罪的债务人身上，如果一人的债务被免除，另一人的债务被追讨，那两人同等亏欠的债务，并没有偏待人的问题。\par

% structure: p0029 source=h2 target=subsection reason=relative-heading-rank; base=h2

\subsection{第十四章——他用一个例子说明他的论点。}

但为了使我说的话通过例子展示得清楚，我们假设某个娼妓生下了一对双胞胎，被遗弃后由别人收养。其中一个没有领洗就死了；另一个领了洗。在这种情况下，我们能说什么\Quote{命运}或\Quote{运气}呢？它们在这里是根本不存在的。有什么\Quote{偏待人}呢？因为在天主那里是没有偏待人的，即使在这些情况下可能存在任何这样的东西，显然他们没有任何东西可以使这一个比那一个更受偏爱，也没有他们自己的功过——无论是善的，使其中一个配得领洗；或是恶的，使另一个配得没有领洗而死。在他们父母身上有任何功过吗？父亲是个奸淫者，母亲是个娼妓？但无论这些功过如何，在如此不同情况下死去的人身上，肯定没有任何不同的功过，而是两者都是共同的。因此，如果既不是命运，因为没有星辰使他们不同；也不是运气，因为没有什么偶然的事件产生这些事；也不是人物或功过的不同造成了这一点；那么，就领了洗的孩子而言，还剩下什么，除非是天主的恩宠，白白地赐给那些为了尊贵而造的器皿；但就未领洗的孩子而言，是天主的义怒，按照团块本身的功过报应给那些为了卑贱而造的器皿？但是，在领了洗的那个孩子身上，我们强迫你们承认天主的恩宠，并使你们相信没有它自己的功过在前；至于那个没有领洗而死去的孩子，为什么那件圣事——你们也承认对所有年龄都是必需的——会缺乏，以及这样在他身上可能受到了什么惩罚，这是你们这些不愿意承认有原罪的人所应留意的。\par

% structure: p0031 source=h2 target=subsection reason=relative-heading-rank; base=h2

\subsection{第十五章——宗徒以不解决的方式回应了这个问题。}

既然在这两个双胞胎的事例中，我们无疑面对同一个案例，那个问题——为什么一个以这种方式而死，另一个以另一种方式而死——的难点，被宗徒以仿佛不解的方式解决了；因为，当他提出关于两个双胞胎类似的事情时，看到那话（不是出于行为，因为他们还没有行善或行恶，而是出于那召叫者的旨意），\BibleQuote{年长的要服事年幼的}，\BibleRef{罗 9:11}以及\BibleQuote{我爱了雅各伯，我恨了厄撒乌}，\BibleRef{罗 9:11}并且他将这深奥之事的可畏感延长，甚至到了说\BibleQuote{所以祂愿意怜悯谁，就怜悯谁；愿意使谁硬心，就使谁硬心}的地步；\BibleRef{罗 9:18}他立刻明白了困难所在，便将一个反驳者的话放在自己面前，他要用权威来制止这反驳者。因为他说道：\Quote{那么，你对我说：\InnerQuote{祂为什么还责怪人呢？有谁能抗拒祂的旨意呢？}}而他对说这话的人回答说：\BibleQuote{人啊！你是谁，竟敢向天主抗辩？受造的怎能对造它的说：\InnerQuote{你为什么这样造了我？}难道陶工不能从同一团泥中，造出一个器皿为尊贵，另一个为卑贱吗？}\BibleRef{罗 9:19}接着，他继续揭示这个伟大而隐藏的奥秘，按他认为适合向人启示的程度，说：\BibleQuote{但如果天主，愿意显示祂的义怒并彰显祂的大能，在极大的忍耐中容忍了那些为毁灭而准备的义怒之器，正是为了在祂所预备为光荣的怜悯之器上，彰显祂光荣的丰富。}\BibleRef{罗 9:22}这不仅是天主恩宠的援助，更是天主恩宠的证明——援助，即是在怜悯之器上；而证明，则是在义怒之器上；因为在这些器上，祂显示自己的忿怒并彰显自己的能力，因为祂的美善如此强大，以至于祂甚至善用邪恶；而在那些器上，祂在怜悯之器上彰显祂光荣的丰富，因为惩罚者的正义从义怒之器所要求的，拯救者的恩宠则给怜悯之器豁免了。那白白赐予一些人的仁慈也不会显明，除非天主对同样有罪且来自同一团泥的另一些人，显示了他们真正应得的惩罚，并以正义的审判定了他们的罪。\BibleQuote{是谁使你不同呢？}\BibleRef{格前 4:7}同一宗徒对那仿佛自夸于自身及自身优点的人说。\BibleQuote{是谁使你不同}于义怒之器呢？当然，是从那因一人的过犯而将众人带往灭亡的丧亡团块中分别出来？\BibleQuote{是谁使你不同呢？}而当他仿佛回答：\Quote{是我的信德使我不同——我的旨意，我的功绩}——他说：\BibleQuote{你所有的一切，哪一样不是领受的呢？既然领受了，为什么还要夸口，好像没有领受一样呢？}——这就是说，好像那使你不同的是属于你自己的。因此，使你不同的是那一位，祂赐给你那使你成为不同的事物，藉免除所应得的惩罚，赐予不当得的恩宠。祂使彼此不同——当黑暗笼罩深渊时，祂说：\BibleQuote{要有光；就有了光，并分开}——亦即，使彼此相异——\BibleQuote{光和黑暗}。\BibleRef{创 1:2}因为当只有黑暗的时候，祂没有找到可以使之不同的事物；而是藉创造光明，祂使之不同；因此，可以对那被成义的恶人说：\BibleQuote{你们从前曾是黑暗，但现在你们在主内是光明。}\BibleRef{弗 5:8}因此，那夸耀的，必须夸耀，不是在自己，而是在主内。祂使彼此不同——对那些尚未出生、尚未行善或行恶的人（为要祂拣选的旨意得以存立，不是出于行为，而是出于那召叫者）说了：年长的要服事年幼的，并后来藉先知的口赞扬同一旨意，说了：\BibleQuote{我爱了雅各伯，但恨了厄撒乌。}\BibleRef{拉 1:2}因为他说了\Quote{拣选}，在这件事上，天主不是发现由别人所造好的、祂可以拣选的东西，而是祂自己所造好的、祂可以找到的东西；正如关于以色列的遗民所记载的：\BibleQuote{因恩宠的拣选，有了遗民；但如果是出于恩宠，就不再出于行为，否则恩宠不再是恩宠了。}\BibleRef{罗 11:5}因此，你们确实愚蠢，当真理宣布说：\BibleQuote{这话不是出于行为，而是出于那召叫者的旨意}，你们却说\Person{雅各伯}是因着天主预知他将来要做的事而被爱，这样就违背了宗徒所说\BibleQuote{不是出于行为}的话，好像他不能说\Quote{不是出于现在的行为，而是出于未来的行为}似的。但他说：\BibleQuote{不是出于行为}，为要颂扬恩宠；\BibleQuote{但如果是出于恩宠，就不在于行为；否则，恩宠就不再是恩宠了。}\par

% structure: p0033 source=h2 target=subsection reason=relative-heading-rank; base=h2

\subsection{第十六章——白拉奇主义者被双胞胎婴儿死亡的情况所驳斥，一个在圣洗恩宠之后，另一个没有。}

但为了除去你们黑暗中的所有藏身之处，我已向你们提出了这样的双胞胎案例：他们未得到父母功劳的帮助，且都在婴儿期初便死去，一个领受了圣洗，另一个未领受圣洗；免得你们说天主预知了他们未来的行为，如同你们论及\Person{雅各伯}和\Person{厄撒乌}时所说的，与宗徒相反。因为祂怎能预知那些在婴儿期就将死去的婴儿中不会存在的事物呢？祂的预知是不会错的。又或者，对于那些被带离此生的人——使恶念不能改变他们的悟性，诡计不能欺骗他们的灵魂——如果连那尚未做出、说出或想出的罪也像已犯下一样受惩罚，这对他们有何益处呢？因为，如果某些人因那些他们既不能从父母那里遗传（如你们所说），也不能因犯下或甚至思虑而自己招致的罪愆而被定罪，是最荒谬、愚蠢和无意义的，那么，你们面前那未领受圣洗的双胞胎兄弟（相对于领受了圣洗者）便无声地问你们：他为何在幸福上与其兄弟不同？为何他受到那样的不幸，以至于当他的兄弟被收养为天主的子女时，他自己却未能领受你们承认对每个年龄段都必要的圣事？如果正如没有命运或宿命，天主也不看情面，那么同样，没有功劳就没有恩宠的恩赐，也没有原罪。对这个哑童，你们完全屈服你们的舌头和声音；对这个一言不发的证人——你们无话可说！\par

% structure: p0035 source=h2 target=subsection reason=relative-heading-rank; base=h2

\subsection{第十七章 [VIII]——连对不完全之善的渴望也是恩宠的恩赐，否则恩宠便会按功劳赐予。}

现在，让我们尽己所能来看看这件事的本质，即他们主张必须先存在于人内，以便人被看作为配得恩宠的帮助，而恩宠不是像白白的赐予那样给予，而是作为应得的回报；于是恩宠便不再是恩宠了。然而，让我们看看这是什么。他们说：\Quote{以恩宠之名，他们如此断言命运，以至于说，除非天主已将善的渴望——且是不完全之善——灌输给不情愿和抗拒的人，否则他既不能远离邪恶，也不能趋向善。}我已经指出他们关于命运和恩宠所说的空洞之言。现在我应该考虑的问题是：天主是否将善的渴望灌输给不情愿和抗拒的人，使他不再不情愿、不再抗拒，而是同意善并愿意善。因为那些人坚持认为，人心中善的渴望始于人自身；而且，这开始的功劳还伴随着完成的恩宠——至少，如果他们连这一点也承认的话。因为佩拉吉乌斯说，如果有恩宠的帮助，善会被\Quote{更容易}地成就。通过这一附加语——即加上\Quote{更容易}——他肯定表明他的观点是，即使缺少恩宠的帮助，善仍能通过自由意志完成，尽管更为困难。但是，让我为当前的对手们规定一下他们在此事上应有的想法，而不提及这异端本身的作者。让我们允许他们凭其自由意志，甚至摆脱佩拉吉乌斯本人，而更留意我们在我们正答复的这封信中所写的话。\par

% structure: p0037 source=h2 target=subsection reason=relative-heading-rank; base=h2

\subsection{第十八章——善的渴望是天主的恩赐。}

因为他们曾认为，我们可以被诘难说，我们声称\Quote{天主把意愿灌输给不愿和抗拒的人}，不是关于任何极大的善，而是\Quote{甚至是不完美的善}。那么，也许他们自己在某种意义上至少为恩宠留了一个位置，认为人可以在没有恩宠的情况下有善的意愿，但只是不完美的善；而对于完美的善，他很难有恩宠的意愿，但若没有恩宠，则根本不可能有。的确，即便是这样，他们也在说天主的恩宠是按我们的功过赐予的。这原是伯拉纠在东方教会会议中，因害怕被定罪而谴责过的。因为如果善的意愿是在没有天主恩宠的情况下由我们自己开始，那么功过本身就已经开始了——恩宠的帮助像是欠债一样来到这功过之上；如此，天主的恩宠就不是白白施予，而是按我们的功过赐予的。但为了给将来的伯拉纠一个答复，主并没有说\Quote{没有我，你们很难做什么}，而是说\BibleQuote{没有我，你们什么也不能做}。\BibleRef{若 15:5}而且，为了也给这些将来的异端者一个答案，在同一句福音言说中，他没有说\Quote{没有我，你们什么也不能\Emph{完成}}，而是说\Quote{什么也不能\Emph{做}}。因为如果他曾说\Quote{完成}，他们可能会说，天主的帮助不是为了开始善（这是出于我们自己），而是为了完成它。但让他们也听听宗徒的话。因为当主说\BibleQuote{没有我，你们什么也不能做}时，在这一句话中，他包含了开始和结束。宗徒确实，像是主言说的阐释者，非常清楚地分别了这两者，当他说：\BibleQuote{因为那在你们内开始了美好工作的，必会完成它，直到耶稣基督的日子。}\BibleRef{斐 1:6}但在圣经中，同一位宗徒的著作里，我们找到更多关于我们正在讨论的问题。因为我们现在谈论的是善的意愿，如果他们愿意承认这是由我们自己开始并由天主完成的，那么让他们看看他们能怎样回答宗徒，当他说道：\BibleQuote{并不是我们凭自己配思考什么事，而是我们的配得出于天主。}\BibleRef{格后 3:5}他说\BibleQuote{思想什么事}—他当然指的是\Quote{思想任何善事}；但思想比意愿更小吗？因为我们意愿的一切我们都思想，但我们并非思想的一切都意愿；因为我们有时也思想我们不意愿的。既然思想比意愿更小——因为一个人可以思想他尚未意愿的善，并通过进展后来意愿他先前无意愿时所思想的——那么，我们怎么对于那较小的事，即思想某件善事，凭自己并不足够，而是我们的配得出于天主；而对于那更大的事，即对某件善事的意愿，没有神圣的帮助，凭自由意志就足够了呢？因为宗徒在这里说的不是\Quote{并不是我们凭自己足够思想那完美的事}，而是说\BibleQuote{思想任何事}，与它相对的\Quote{什么}是\Quote{没有}。这正是主所说的意思：\BibleQuote{没有我，你们什么也不能做}。\par

% structure: p0039 source=h2 target=subsection reason=relative-heading-rank; base=h2

\subsection{第19章【IX】——他解释那些\Person{伯拉纠派}滥用的圣经章节。}

但确然，关于所写的：\BibleQuote{人心筹算自己的道路，舌头的应对出于主。}\BibleRef{箴 16:1}他们被不完全的理解误导，以致认为预备人心——即开始善——属于人而不需天主恩宠的帮助。愿这远离于应许的子女！仿佛当他们听到主说\BibleQuote{没有我，你们什么也不能做}\BibleRef{若 15:5}，他们会反驳主说：\Quote{看，没有你，我们也能预备人心}；或者当他们听到\Person{保禄宗徒}说\BibleQuote{并不是我们凭自己配思考什么事，而是我们的配得出于天主}\BibleRef{格后 3:5}，仿佛他们也会反驳他说：\Quote{看，我们凭自己足以预备自己的心，从而也思想一些善事；因为谁能在没有善意的情况下预备自己的心向善呢？}愿任何人不要这样理解这经文，除了骄傲的自由意志维护者和公教信仰的背弃者！因此，既然经上记载：\BibleQuote{人心筹算自己的道路，舌头的应对出于主。}那就是说，人预备自己的心，然而不是没有天主的帮助，天主如此触动人心，以致人预备心。但在舌头的应对中——即神圣的舌头对已预备的心所作的回答——人没有份；而是全部出于主天主。\par

% structure: p0041 source=h2 target=subsection reason=relative-heading-rank; base=h2

\subsection{第20章——天主的行动甚至在人的行为中也是必要的。}

因为正如经上所说：\Quote{预备人心是人分内的事，舌头的回答却来自主；}同样也说道：\Quote{你要张开口，我要将它充满。}因为，若非藉祂的帮助——没有祂我们什么也不能做——我们不能开口，但我们仍是藉祂的协助并藉我们自己的行动力开口，而主则无需我们行动便将它充满。因为预备人心和张开口，不就是预备意志吗？然而在同一部圣经中也读到：\Quote{意志是由主所预备的，}\EditorialNote{箴言第八章}以及\Quote{你要开启我的嘴唇，我的口要宣扬你的赞美。}所以天主在经文中劝诫我们预备我们的意志。读经说\InnerQuote{预备人心是人分内的事}；然而，为使一个人能做到这事，天主帮助他，因为意志是由主所预备的。还有：\InnerQuote{张开口。}祂这样以命令的方式说，以致没有人能够做到，除非藉祂的协助，而对祂曾说过：\Quote{你要开启我的嘴唇。}这些人中有谁如此愚蠢，竟主张口是一回事，嘴唇是另一回事？并且以惊人的琐碎说，人自己张口，而天主开启人的嘴唇？然而天主甚至从那种荒谬中约束他们，当祂对祂的仆人梅瑟说：\Quote{我要开启你的口，并教导你应当说什么。}\BibleRef{出 4:12}因此，在祂说\Quote{你要张开口，我要将它充满}这句话时，似乎其中一样属于人，另一样属于天主。但在这里，当祂说\Quote{我要开启你的口，并教导你}时，两者都属于天主。这是为什么呢？除非因为在一种情况下，祂与人合作，人作为行动者；在另一种情况下，祂独自完成？\par

% structure: p0043 source=h2 target=subsection reason=relative-heading-rank; base=h2

\subsection{第二十一章——人不行任何天主不使他去行的善事。}

因此，天主在人身上行许多人不行的善事；但人不行任何天主不使人去行的善事。因此，人若没有从主而来，就不会有对善的渴望，除非那是善；但如果是善，我们只有从那至善且不能共有的善者那里获得它。因为对善的渴望不就是爱吗？宗徒若望毫无歧义地说：\BibleQuote{爱出自天主。}\BibleRef{若一 4:7}它的开始不是出于我们自己，它的完成也不是出于天主；但若爱出自天主，我们便从天主那里领受整个的爱。愿天主完全阻止这种愚妄，即我们认为自己是祂恩惠中的首位，而祂是末位——因为\Quote{祂的仁慈必先于我。}祂正是那被忠实而真诚地歌颂的：\Quote{因为你以甘饴的福气临到他。}这里，还有什么比我们所谈论的对善的渴望更恰当地理解呢？因为善开始被渴慕，正是当它开始变得甘甜之时。但是当善行是出于对惩罚的恐惧，而非出于对正义的爱时，善还未被善行。那在行为中看似做成的，若一个人能免受惩罚而逃避，他宁愿不做，那么在心里也并非真正做成。因此，\Quote{甘饴的福气}就是天主的恩宠，藉此在我们内产生：祂所命令我们的，使我们喜悦，我们渴慕它——也就是说，我们爱它。如果天主不先于我们，不仅不能完成，甚至不能由我们开始。因为，若没有祂我们实际上不能做什么，我们就既不能开始，也不能完成——因为开始，经上说\Quote{祂的仁慈必先于我；}完成，经上说\Quote{祂的仁慈必跟随我。}\par

% structure: p0045 source=h2 target=subsection reason=relative-heading-rank; base=h2

\subsection{第二十二章 [X.]——蒙召者是按谁的旨意被选？}

那么，为什么在接下来的内容中，当他们提到自己认为的观点时，他们承认说：\Quote{恩宠也协助每人的善志，但它并不将美德的热望注入一颗不情愿的心。}因为他们这样说，仿佛人凭自己，没有天主的协助，就有一个善志和对美德的渴望；而这先有的功劳，堪受随后而来的天主恩宠的协助。因为他们或许以为宗徒如此说：\BibleQuote{我们知道，天主使一切事协助那些爱祂的人，即那些按旨意蒙召的人，} \BibleRef{罗 8:28}从而期望人们理解是按人的旨意，这旨意作为一种善功，蒙召的天主的仁慈会随之而来；却不知道经上所说：\Quote{按旨意蒙召的人，}乃是要理解为天主的旨意，而非人的旨意，藉此祂在创世以前预知并预定他们成为与祂儿子肖像相似的人，就拣选了他们。因为并非所有蒙召的都是按旨意蒙召的，既然\BibleQuote{许多人蒙召，少数人蒙选。}\BibleRef{玛 20:16}因此，那些在创世以前蒙拣选的，就是按旨意蒙召的人。关于天主的这个旨意，我前面已经提及关于双胞胎厄撒乌和雅各伯所说的：\BibleQuote{为使天主的旨意按照拣选得以坚定，不是出于行为，而是出于那召叫的；经上说：年长的要服事年幼的。}\BibleRef{罗 9:11}天主的这个旨意也在另一处提及，当致书弟茂德时，他说：\BibleQuote{你要按照天主的能力为福音劳苦，祂拯救了我们，以圣召召叫了我们，不是按我们的行为，而是按祂的旨意和恩宠，这恩宠在永恒年代以前已在基督耶稣内赐给了我们，如今却因我们的救主耶稣基督的显现而显示出来。}\BibleRef{弟后 1:8}那么，这就是天主的旨意，关于此旨意经上说：\Quote{祂使一切事协助那些按旨意蒙召的人。}但随后的恩宠确实协助人的善志，然而，若恩宠不先存在，善志本身也不会存在。人的渴望，即被称为善的渴望，虽然在开始存在时受恩宠协助，却并非没有恩宠就开始，而是由祂所激发，宗徒论及祂说：\BibleQuote{但感谢天主，祂为你们把同样的渴望放在弟铎心里。}\BibleRef{格后 8:16}如果天主赐予渴望，使每个人能为他人拥有它，那么谁还会赐予，使人能为自己拥有它呢？\par

% structure: p0047 source=h2 target=subsection reason=relative-heading-rank; base=h2

\subsection{第二十三章——凡天主所命令的，没有一样不是天主所赐予的。}

既然这些事如此，我看见主在圣经中并没有命令人任何事，为考验人的自由意志，而这件事不是因祂的良善而开始的，或者不是为证明恩宠的助佑而求得的；人也不会因信德的开端而从恶变为善，除非天主的白白且无代价的仁慈在他内成就这事。其中一人回想自己的思想，正如我们在圣咏中所读到的，说：\BibleQuote{难道天主忘记施恩了吗？难道祂在愤怒中收起了祂的仁慈吗？于是我说：我现在开始了；这个改变是出于至高者的右手。}因此，当他说\BibleQuote{我现在开始了}的时候，他并没有说\Quote{这个改变是出于我的意志}，而是说\BibleQuote{出于至高者的右手}。所以，因此，要这样思想天主的恩宠：从他良好改变的开始，直到完成的终末，凡夸耀的，应因主而夸耀；因为，正如没有主没有人能完成善事，同样，没有主也没有人能开始善事。但让这成为本书的结束，以便读者的注意力得到更新和加强，为接下来的内容。\par

\SourceRecord{Translated by Peter Holmes and Robert Ernest Wallis, and revised by Benjamin B. Warfield.；Nicene and Post-Nicene Fathers, First Series；Vol. 5.；Edited by Philip Schaff.；Buffalo, NY: Christian Literature Publishing Co.,；1887.；<http://www.newadvent.org/fathers/15092.htm>.}

% structure: p0001 source=h1 target=section reason=page-title

\section{驳斥佩拉纠派的两封信（第三卷）}

\Emph{奥古斯丁继续驳斥佩拉纠派在同一封送往得撒洛尼的信中诽谤性地提出的其他事项；并针对他们的异端，阐释真正天主教徒关于法律用途的说法；他们关于圣洗圣事的效力和德能的教导；关于旧约和新约之间的差异；关于先知和宗徒的义德和成全；关于基督身上的罪之称谓，当祂被说成在罪恶肉体的相似中定了罪，或成了罪；最后，他们关于在未来生命中完成诫命的宣认。}\par

% structure: p0003 source=h2 target=subsection reason=relative-heading-rank; base=h2

\subsection{第一章 — 陈述}

还有一些他们诽谤性地反对我们的事情；他们还没有开始阐述他们自己的想法。但为了避免这些文字的冗长令人反感，我将他们反对的事项分为两卷——前一卷已完成，即这部完整作品的第二卷，我在这里开始另一卷，并将其作为第三卷与第一卷和第二卷合并。\par

% structure: p0005 source=h2 target=subsection reason=relative-heading-rank; base=h2

\subsection{第二章 — 佩拉纠派对旧约法律用途的歪曲}

他们声称\Quote{我们说旧约法律被赐下，不是为了它可以成义顺从的人，而是为了它可以成为更大罪的原因。}当然，他们不理解我们关于法律的说法；因为我们说的是宗徒所说的话，而他们不理解宗徒。因为谁能说那些顺从法律的人没有被成义呢？因为除非他们被成义，他们就不能顺从？但我们说，借着法律，天主所愿意做的事被听到；借着恩宠，法律被遵守。\Quote{因为不是法律的听众，}宗徒说，\Quote{在天主前是义人，而是法律的实行者将被成义。}\BibleRef{罗 2:13}因此，法律使人成为义德的听闻者，恩宠使人成为实行者。\Quote{因为法律所不能做到的，}同一宗徒说，\Quote{因它因肉体而软弱，天主派遣了自己的儿子在罪恶肉体的相似中，并为罪在肉体中定了罪，为要使法律的义德成就在我们身上，我们不是照着肉体，而是照着圣神行事。}\BibleRef{罗 8:3}这就是我们所说的——让他们祈祷，但愿有一天他们能理解，而不是争辩以致永远不理解。因为法律不可能由肉体，即由肉性的骄傲所成全，在这种骄傲中，那些傲慢的人，他们不知道天主的义德——即来自天主的义德，使人成为义人——并且渴望建立自己的义德——就好像凭自己的意愿，不受上天的帮助，法律就能成全一样——没有顺服于天主的义德。\BibleRef{罗 10:3}因此，法律的义德成就在那些不照着肉体——即不照着人，不知道天主的义德并渴望建立自己的义德——而照着圣神行事的人身上。但谁照着圣神行事呢？除非被天主圣神引导的人。\Quote{因为凡被天主圣神引导的，都是天主的子女。}\BibleRef{罗 8:14}因此\Quote{文字叫人死，圣神却叫人活。}\BibleRef{格后 3:6}而且文字并不因为能叫人死就是恶的；它只是定了罪人的过犯。他说：\Quote{因为法律是圣的，诫命是圣的、公义的、好的。那么，}他说，\Quote{那好的成了我的死亡吗？绝对不是；而是罪，为要显出是罪，借着那好的在我内产生了死亡，以致借着诫命，罪就成为极其罪恶的。}\BibleRef{罗 7:12}这就是\Quote{文字叫人死}的意思。\Quote{因为死亡的刺是罪，而罪的力量是法律；}\BibleRef{格前 15:56}因为借着禁令，它增加了罪的欲望，从而杀死一个人，除非恩宠来帮助使他活过来。\par

% structure: p0007 source=h2 target=subsection reason=relative-heading-rank; base=h2

\subsection{第三章 — 天主教教义的圣经证实}

这就是我们所说的；这是他们反对我们的事，即我们说\Quote{法律被赐下是为了成为更大罪的原因。}他们没有听到宗徒说：\Quote{因为法律惹起愤怒；因为哪里没有法律，哪里就没有过犯；\BibleRef{罗 4:15}}和：\Quote{法律是为过犯而添加的，直到那应许的后裔来到；\BibleRef{迦 3:19}}和：\Quote{如果曾有一种能赐生命的法律被赐下，义德就真的出自法律了；但圣经把所有人都圈在罪中，为要使那因耶稣基督的信德而得的恩许，能赐给信者。}因此，从西奈山传下来的旧约，法律在那里被赐下，生出为奴的儿子，那就是阿加尔。他说：\Quote{我们却不是婢女的子女，而是自由妇女的子女。}因此，那些接受文字法律的人不是自由妇女的子女，他们不仅显出是罪人，而且更是犯法者；但那些接受恩宠之圣神的人，就是借着她，那圣洁、公义、良善的法律本身得以成全的人。这就是我们所说的：让他们留心，不要争辩；让他们寻求光照，不要诬告。\par

% structure: p0009 source=h2 target=subsection reason=relative-heading-rank; base=h2

\subsection{第四章 — 关于圣洗圣事效果的歪曲}

\Quote{他们断言，}他们说，\Quote{而且，圣洗并不使人成为新人——即并不赐予完全的罪赦；但他们主张，人一部分成为天主的子女，另一部分仍属于今世，即属于魔鬼。}他们欺骗；他们设陷阱；他们搪塞；我们并不这样说。因为我们是说：所有魔鬼的子女也都是今世之子；但并非所有今世之子也都是魔鬼的子女。我们绝不敢说，圣祖\Person{亚巴郎}、\Person{依撒格}、\Person{雅各伯}，以及这类的其他人在婚配中生子时是魔鬼的子女，也不敢说那些直到如今和将来仍继续生子的信徒是魔鬼的子女。而且我们不能反驳主所说的，\BibleQuote{今世之子也娶也嫁。}\BibleRef{路 20:34}因此，有些人是今世之子，却并非魔鬼的子女。因为尽管魔鬼是一切罪恶的创始者和根源，但并非每一罪都使人成为魔鬼的子女；因为天主的子女也犯罪，因为若他们说没有罪，就是自欺，真理也不在他们内。\BibleRef{若一 1:8}但他们犯罪，是由于他们仍属于今世的那种状态；但因着那使他们成为天主的子女的恩宠，他们必定不犯罪，因为凡由天主生的，就不犯罪。\BibleRef{若一 3:9}但不信使人成为魔鬼的子女；不信被特别称为罪，好像它是唯一的罪，如果没有说明罪的性质。正如当提到\Quote{宗徒}时，若不指明是哪位宗徒，则只理解为\Person{保禄}；因他以多封书信而更为人所知，且他比众人更劳苦。为此，主论及圣神时所说的，\BibleQuote{祂要指证世界关于罪，}\BibleRef{若 16:8}祂的意思是叫我们理解不信；因为祂在解释时说过这话，\BibleQuote{关于罪，因为他们不信我，}\BibleRef{若 16:9}以及当祂说，\BibleQuote{我若没有来对他们讲话，他们就没有罪。}\BibleRef{若 15:22}因为祂的意思不是说他们先前没有罪，而是想指出他们那种不信，即当祂在他们面前对他们讲话时，他们仍不信祂；因为他们属于那一位，关于他宗徒说，\BibleQuote{按着空中掌权者的首领，就是现今在不信之子身上运行的。}\BibleRef{弗 2:2}因此，凡没有信德的人都是魔鬼的子女，因为他们内在之人中没有理由可以使其因人性软弱、无知或任何邪恶意志所犯的罪获得赦免。但那些是天主的子女的人，如果他们\BibleQuote{说自己没有罪，就是自欺，真理不在他们内；但立刻}（正如接续所说的）\BibleQuote{当他们承认自己的罪时}（这是魔鬼的子女不做的，或不是按照天主子女特有的信德去做的），\BibleQuote{祂是信实的、公义的，必要赦免他们的罪，并洗净他们的一切不义。}\BibleRef{若一 1:8}为使我们的说法更易理解，让我们亲自听\Person{耶稣}的话，祂显然是对天主的子女说话，祂说：\BibleQuote{你们纵然是恶的，尚且知道把好东西给你们的儿女，何况你们的天父，岂不更要把好东西赐给求祂的人！}\BibleRef{玛 7:11}因为如果这些人不是天主的子女，祂就不会对他们说，\Quote{你们的天父。}然而祂说他们是恶的，并知道把好东西给他们的儿女。那么，他们在身为天主的子女这一点上是恶的吗？绝无此事！他们是恶的，是因为他们仍属于今世，尽管如今借着圣神的凭据成为了天主的子女。\par

% structure: p0011 source=h2 target=subsection reason=relative-heading-rank; base=h2

\subsection{第五章—圣洗除去一切罪过，但不立即治愈所有软弱。}

因此，\OriginalTerm{圣洗}{Baptism}确实洗去一切罪恶——绝对一切罪恶，无论是行为、言语或思想，无论是原罪还是后加，无论是因无知所犯或明知而纵容；但它并不除去\OriginalTerm{软弱}{weakness}，这软弱是重生之人打好仗时所抵抗的，但作为人，当他暂时陷于任何过犯时，他所同意的；因前者而感恩欢呼，因后者而叹息祈祷。因前者说：\BibleQuote{我当如何报答主赐给我的一切恩惠？}因后者说：\BibleQuote{宽免我们的罪债。}\BibleRef{玛 6:12}因前者说：\BibleQuote{上主，我的力量，我必要爱慕祢。}因后者说：\BibleQuote{上主，求祢怜恤我，因为我软弱。}因前者说：\BibleQuote{我的眼睛常仰望主，因为祂必从网罗中拔出我的脚。}因后者说：\BibleQuote{我的眼睛因激怒而苦恼。}圣经中充满了无数这样的章节，\OriginalTerm{天主的子女}{children of God}因\OriginalTerm{信德}{faith}，时而为天主的恩惠欢欣，时而为自己的邪恶哀叹，只要他们在这生活的软弱中仍是\OriginalTerm{世俗的子女}{children of this world}；然而，天主将他们与\OriginalTerm{魔鬼的子女}{children of the devil}区别开来，不仅是藉着\OriginalTerm{重生之洗}{laver of regeneration}，而且还藉着那藉\OriginalTerm{爱德}{love}运行的\OriginalTerm{信德}{faith}的\OriginalTerm{正义}{righteousness}，因为义人因信德而生活。但这软弱，我们与之奋战，时而失败，时而进步，直至身体死亡，且极重要地能在我们内战胜什么，将被另一次\OriginalTerm{重生}{regeneration}所消灭，对此主说：\BibleQuote{在重生时，当人子坐在祂光荣的宝座上，你们也要坐在十二个宝座上，等等。}\BibleRef{玛 19:28}无疑地，在这段话中，主将末次复活称为重生，而\Person{保禄}宗徒也将之称为\OriginalTerm{义子}{adoption}和\OriginalTerm{救赎}{redemption}，他说：\BibleQuote{连我们这些具有圣神初果的，也在自己心中叹息，等候义子，就是我们身体的救赎。}\BibleRef{罗 8:23}我们不是已经藉着圣洗重生、得为义子并被救赎了吗？然而，还有一个重生、一个义子、一个救赎，我们现今应当耐心等候它终将来临，以便那时我们不再是世俗的子女。因此，凡从圣洗中除去我们藉其唯一领受之物的，是败坏信德；但如今就已将那藉圣洗确将赐予但我们尚待来日获得的归于它的，是割断\OriginalTerm{望德}{hope}。因为若有人问我，我们是否因圣洗而得救，我无法否认，因为宗徒说：\BibleQuote{祂藉着重生之洗和圣神的更新拯救了我们。}\BibleRef{铎 3:5}但如果他问，是否藉着同一洗濯，祂已经绝对地、完全地拯救了我们，我答说：并非如此。因为同一位宗徒也说：\BibleQuote{我们得救是在乎希望；但看得见的希望就不是希望了：谁还盼望所见之物呢？如果我们盼望那未看见的，我们就耐心等候。}\BibleRef{罗 8:24}因此，人的\OriginalTerm{救恩}{salvation}在圣洗中完成，因为无论他从父母继承的何种罪恶，或他自身在受洗前所犯的，都被赦免；但将来的救恩将是如此，以致他完全不能犯罪。\par

% structure: p0013 source=h2 target=subsection reason=relative-heading-rank; base=h2

\subsection{关于旧约和古圣人的诽谤}

如果事情如此，那么他们随后向我们提出的那些异议就可被驳倒。因为哪个\OriginalTerm{大公的}{catholic}会说他们指控我们所说的，即\Quote{圣神在旧约中不是美德的帮助者}，除非我们如此理解\Quote{旧约}，正如宗徒称之为\BibleQuote{出自西乃山，生子为奴}？但因其中预象了\OriginalTerm{新约}{New Testament}，当时那些按时代秩序领悟此理的天主的人，的确是\OriginalTerm{旧约}{Old Testament}的\OriginalTerm{管家}{stewards}和承载者，但被显示为\OriginalTerm{新约}{New Testament}的\OriginalTerm{继承人}{heirs}。难道我们要否认那说\BibleQuote{天主，求祢在我内造一个洁净的心，使我内心重新有一个正直的灵}的人属于新约？或是那说\BibleQuote{祂把脚放在磐石上，稳定了我的脚步，并在我口中赐给一首新歌，就是赞美我们天主的诗篇}的人？或是那位在西乃山旧约之前的信德之父，宗徒论到他说：\BibleQuote{弟兄们——我按人的常理说——即使是人的遗嘱，既经批准，就没有人废弃或增添。应许原是向亚伯拉罕和他的后裔说的。不说'向许多后裔'，而是单指向一位，就是基督。我是说，}他说，\BibleQuote{天主所批准的遗嘱，那四百三十年以后的法律，不能使之失效，以致应许落空。因为如果继承出于法律，就不出于应许；但天主曾藉应许将产业赐给亚伯拉罕。}\par

% structure: p0015 source=h2 target=subsection reason=relative-heading-rank; base=h2

\subsection{新约比旧约更古老；但它是后来才启示的。}

在此，如果我们问：他所称为由天主所确认、而四百三十年以后制定的法律未能使之失效的此一约，应当理解为新约还是旧约？谁能犹豫而不回答：\Quote{新约，但隐藏于预言性的影子里，直到那应揭示于基督的时刻到来}？因为如果我们说是旧约，那么那来自西奈山、生子为奴的将是什么呢？因为四百三十年以后制定了法律，他断言这法律不能使亚巴郎应许的盟约失效；他愿意这亚巴郎所立的约更属于我们，他愿意我们作为自由妇人的子女，而非婢女的子女，是凭应许而非凭法律的继承人，正如他所说的：\BibleQuote{因为若继承是出于法律，便不是出于应许；但天主却藉应许赐给了亚巴郎。}\BibleRef{迦 3:18} 因此，因为法律是四百三十年以后制定的，它介入使得过犯增多；\BibleRef{罗 5:20} 既然因罪，人骄傲自恃于自己的义被显明为过犯，而何处罪恶增多，恩宠便更加增多，\BibleRef{罗 5:20} 这是借着如今在法律上跌倒、投靠天主仁慈的谦卑之人的信德。因此，当他说：\BibleQuote{因为若继承是出于法律，便不再是出于应许；但天主却藉应许赐给了亚巴郎。}\BibleRef{迦 3:18} 好像有人会对他说：\Quote{那么为什么后来要制定法律呢？}他便接着说道：\BibleQuote{那么法律是什么呢？}\BibleRef{迦 3:19} 对这问话他立即答道：\BibleQuote{它是因为过犯而添设的，直到那应许所指向的子孙来到。}\BibleRef{迦 3:19} 他后来又这样说道：\BibleQuote{因为如果那些属于法律的人成为继承人，信德就落空了，应许也归于无效。因为法律带来义怒：哪里没有法律，哪里就没有过犯。}\BibleRef{罗 4:14} 他在前一段证据中所说的：\BibleQuote{因为若继承是出于法律，便不再是出于应许；但天主却藉应许赐给了亚巴郎}，他在后一段证据中所说的是：\BibleQuote{因为如果那些属于法律的人成为继承人，信德就落空了；应许也归于无效。}这充分表明，天主藉应许赐给亚巴郎的，是属于我们的信德（这信德当然属于新约）。他在前一段证据中所说的：\BibleQuote{那么法律是什么呢？}并答道：\BibleQuote{它是为了过犯而添设的}，他在后一段证据中立即接着说：\BibleQuote{因为法律带来义怒：哪里没有法律，哪里就没有过犯。}\par

% structure: p0017 source=h2 target=subsection reason=relative-heading-rank; base=h2

\subsection{第八章——亚巴郎之前和之后的所有义人都是应许和恩宠的子女。}

那么，无论是亚巴郎，或是他之前或之后的义人，直到亲自领受那来自西奈山、生子为奴之约的梅瑟，以及他以后的先知们，以及直到若翰洗者的天主圣人们，他们都是按依撒格——自由妇人的儿子——而成为应许和恩宠的子女——不是出于法律，而是出于应许，是承继天主并与基督同为承继者。我们绝不能否认，义人诺厄和更早时期的义人，以及从那时直到亚巴郎时期凡能成为义人者，无论是明显还是隐藏的，都属于那在上的耶路撒冷，她是我们的母亲，尽管他们在时间上早于撒辣，而撒辣本身孕育了那自由母亲的预言和预像。那么，在亚巴郎之后——他曾得到应许，要成为万民之父——所有中悦天主的人，岂不更应被尊为应许的子女！因为从亚巴郎和跟随他的义人开始，所找到的不是更真实的血统，而是更明确的预言。\par

% structure: p0019 source=h2 target=subsection reason=relative-heading-rank; base=h2

\subsection{第九章——谁是旧盟约的子女。}

但那些人属于旧约，\BibleQuote{它来自西奈山，生子为奴}，即是哈加尔。他们领受了神圣、公正和良善的法律，却以为文字足以赐予他们生命；他们不寻求天主的仁慈，好能成为法律的实行者，反而对天主的正义无知，并企图建立自己的正义，不顺服于天主的正义。属于这一类的，有在旷野中抱怨天主、制造偶像的群众；也有在应许之地本身与异神行淫的群众。然而，这些群众即使在旧约本身也受到严厉的斥责。此外，当时凡只追求天主在那里所应许的现世福乐，不明白这些应许在新约下所预示的，并以获取这些应许的渴望和失去它们的恐惧来遵守天主的诫命——他们确实没有遵守，只是自以为遵守。因为在她们里面没有由爱德运行的信德，只有属世的贪欲和肉体的恐惧。那如此履行诫命的人，无疑是不情愿地履行，内心并不实行；因为如果能不受到惩罚而忽略这些他所渴望和畏惧之事，他宁愿完全不履行。因此，在他内心的意志中，他是有罪的；而颁布命令的天主正是在这里察看。那些是地上的耶路撒冷的子女，关于她，宗徒说：\BibleQuote{她与自己的子女同为奴隶}\BibleRef{迦 4:25}，并且她属于旧约，\BibleQuote{它来自西奈山，生子为奴，即是哈加尔}。钉死主并继续留在同样不信中的，就是这类的人。因此，如今在犹太人的一大群人中仍有他们的子女，尽管新约如预言所说，已经藉基督的血显明并坚立；福音从祂受洗并开始教导的那河，直到地极，已被传扬。而根据他们所诵读的预言，这些犹太人分散于普天下，好使连他们的著作也不缺少对基督真理的见证。\par

% structure: p0021 source=h2 target=subsection reason=relative-heading-rank; base=h2

\subsection{第十章——旧法律也是由天主所赐。}

正是为此缘故，天主制定了旧约，因为天主乐意将天上的应许隐藏在世上的应许之中，仿佛作为奖赏，直到时期圆满；并赐给一个渴求世上福乐、因而心硬的民族一部法律——这部法律虽是属神的，却刻在石版上。因为，除了旧约诸书中的圣事——那些仅为象征意义而颁定的（虽然在其中，由于它们应被属灵地理解，法律恰当地被称为属神的）——其余关于虔敬与善行的事项，绝对不能通过某种解释而归于某种象征意义，而必须按其字面实行。毫无疑问，没有人会怀疑，天主的这部法律不仅当时对那个民族是必要的，而且如今对我们正确安排生活也是必要的。因为，如果基督为我们卸去了那许多规条的重轭，以致我们不再按肉体受割礼，不再用牲畜献祭，不在每周第七日安息（甚至不作必要的工作），以及其他此类事情；而是按属灵意义遵守它们，在除去预表的阴影之后，警醒于它们所预示之事的真光中；那么，我们岂能说，经上记载\BibleQuote{凡捡到他人遗失的任何物品，应归还失主}\BibleRef{肋 6:3}这条诫命与我们无关呢？还有其它许多类似的事，使人学习虔敬正直地生活？尤其是十诫本身，它包含在那两块石版上（除去对安息日的肉体遵守——它象征属灵的圣化和安息）？谁能说基督徒不应以宗教的服从事奉独一天主，不应拜偶像，不应妄称上主的名，应孝敬父母，不应犯奸淫、杀人、偷盗、作假证，不应贪恋别人的妻子或任何属于别人的东西？谁如此不虔敬，竟说他因是基督徒，不处于法律之下，而处于恩宠之下，就不遵守这些法律诫命？\par

% structure: p0023 source=h2 target=subsection reason=relative-heading-rank; base=h2

\subsection{第十一章——旧约与新约子女的区别。}

但显然有这重大区别：那些处于法律之下、被文字杀死的人，他们行这些事或出于获得地上幸福的欲望，或出于害怕失去它；因此他们并非真正行这些事，因为肉身的欲望——借此罪反被交易或增加——并未因另一种欲望而得到医治。这些人属于旧约，它生子为奴；因为肉身的恐惧和欲望使他们成为仆人，福音的信德、望德和爱德并未使他们成为子女。但那些处于恩宠之下、被圣神赋予生命的人，他们以藉着爱德而工作的信德行这些事，怀着对善事的希望——不是肉身的而是属神的，不是地上的而是天上的，不是暂时的而是永恒的；他们尤其相信中保，藉着祂毫不怀疑恩宠之神被施予他们，使他们能善行这些事，并在犯罪时获得赦免。这些人属于新约，是恩许的子女，由天主圣父和自由的母亲重生。所有古时的义人，连旧约的僕人梅瑟本人，也是这类人——他是新约的继承人；因为凭着我们生活的同一信德，他们生活了，相信基督的降生成人、受难和复活是未来之事，而我们相信这些已经实现——直到洗者若翰本人，他仿佛是旧约分施的某种界限；他预示中保将要来临，不是以未来的阴影或预表暗示，也不是以预言宣告，而是用手指指着祂说：\BibleQuote{请看天主的羔羊；请看除免世罪者。}\BibleRef{若 1:29}仿佛在说：许多义人渴望看见祂，从人类之初就相信祂将要来临，关于祂的应许曾对亚巴郎说过，梅瑟写过，法律和先知都为祂作证：\BibleQuote{请看天主的羔羊，除免世罪者。}从若翰开始，此后所有关于基督的事都开始成为过去或现在，而在此之前的所有义人则相信、希望、渴望它们是未来之事。因此，信德是相同的，不仅对于那些虽未被称为基督徒但实际上是基督徒的人，而且对于那些不仅已是而且也被称为基督徒的人；两者都因圣神而有同一恩宠。为此宗徒说：\BibleQuote{我们既然具有同一信德之神，正如经上所载：\InnerQuote{我信了，所以我说了}，我们也信，所以也说话。}\BibleRef{格后 4:13}\par

% structure: p0025 source=h2 target=subsection reason=relative-heading-rank; base=h2

\subsection{第十二章——旧约是一回事，旧约书是另一回事。}

因此，因着已经流行的言语习惯，一方面，法律和众先知——他们预言直到若翰——被称为\Quote{旧约}；虽然这更确定地被称为\Quote{旧盟约}而非\Quote{旧约}；但这一名称又被权威以另一种方式使用，无论是明言还是暗示。因为宗徒是明确的，当他说：\BibleQuote{直到今日，几时诵读梅瑟时，还有同样的帕子盖在旧约的诵读上；因为它是未被揭开的，因为它在基督内已失效了。}\BibleRef{格后 3:14}因为如此无疑地，旧约是指梅瑟的职事。此外，他说：\BibleQuote{我们应以圣神的新，而不以文字的旧来事奉。}\BibleRef{罗 7:6}以文字的名称表明同一个约。在另一处也说：\BibleQuote{祂也使我们成为新约的合格仆役；不是文字，而是圣神：因为文字叫人死，圣神却叫人活。}\BibleRef{格后 3:6}在这里，藉着提到新约，他当然意味着前者被理解为旧约。但更明显的是，虽然他没有说旧或新，但他区分了两个约和亚巴郎的两个儿子——一个是婢女的，一个是自由妇的，正如我上面提到的。还有什么比他所说的更明确：\BibleQuote{请告诉我，你们愿意受法律管制的人，没有听过法律吗？因为经上记载：亚巴郎有两个儿子，一个是婢女所生，一个是自由妇人所生。那婢女所生的是按血统生的，那自由妇人所生的是因恩许生的。这些都是寓意：这些是两约，一个出自西乃山，生子为奴，即哈加尔。西乃山在阿剌伯，它与现今的耶路撒冷相关，因为耶路撒冷与她的儿女同为奴隶。但那在上面的耶路撒冷是自由的，她是我们的母亲？}还有什么比这更清楚、更确定、更远离一切晦暗和歧义，对于恩许的子女而言？稍后又说：\BibleQuote{弟兄们，我们像依撒格一样，是恩许的子女。}\BibleRef{迦 4:28}又稍后说：\BibleQuote{弟兄们，我们不是婢女的子女，而是自由妇的子女。}\BibleRef{迦 4:31}并藉着基督使我们得自由。因此，让我们选择，称古时的义人为婢女的儿女还是自由妇的儿女。我们决不可说他们是婢女的儿女；因此，如果他们是自由妇的儿女，他们就属于圣神的新约；作为赋予生命者，宗徒将祂与叫人死的文字相对。有何理由说他们不属于新约的恩宠呢？我们正是藉着他们的言语和面容，指控并反驳这些极端狂妄和忘恩负义的同一恩宠的仇敌。\par

% structure: p0027 source=h2 target=subsection reason=relative-heading-rank; base=h2

\subsection{第十三章——为什么一个盟约称为旧约，另一个称为新约。}

但有人会说：\Quote{那在四百三十年后由梅瑟所传的怎么被称为旧的，而早在多年以前给亚巴郎所传的却被称为新的呢？}凡在此事上感到困扰的人，不要争辩，而要诚恳地先明白：当那从较早时期传的被称为\InnerQuote{旧}，而这从较晚时期传的被称为\InnerQuote{新}时，是从启示的角度来看它们的名称，而非从设立的角度。因为旧约是通过梅瑟启示的，梅瑟传了神圣、公义而良善的法律，这法律的目的不是废除罪，而是使人认识罪——借此使骄傲的人被定罪，他们想要建立自己的义，好像不需要天主的帮助，并因文字的罪而逃向恩宠之神，不是靠自己的义成义，而是靠天主的义——即天主所赐予他们的义。因为正如同一宗徒所说：\InnerQuote{藉法律，惟有罪的认识。但是如今，天主的义在法律之外已显示出来，法律和先知都为此作证。}\BibleRef{罗 3:20}因为法律本身，因着在其中无人能成义，就为天主的义提供了见证。因为这一点是明显的：在法律下没有人能在天主面前成义，因为\InnerQuote{义人因信德而生活。}\BibleRef{迦 3:11}因此，虽然法律不义人，当他被定罪为犯法时，却将人送往成义的天主那里，从而为天主的义作证。再者，先知们预言基督，也见证了天主的义：\InnerQuote{基督由天主成了我们的智慧、正义、圣化者和救赎者，为应验经上所记载的：凡要夸耀的，应因主而夸耀。}\BibleRef{格前 1:30}但那法律从一开始就被隐藏了，当时本性自身定了恶人的罪，因为他们对别人做了自己不愿别人对自己做的事。但新约在基督内的启示是在祂降生成人时显明的，其中显现了天主的义——即来自天主的、赐予人的义。因此他说：\InnerQuote{但是如今，天主的义在法律之外已显示出来。}\BibleRef{罗 3:21}这就是为什么前者被称为旧约，因为它是在较早时期启示的；后者被称为新约，因为它是在较晚时期启示的。总之，旧约属于旧人，人必须从此开始；新约属于新人，人应从旧状态过渡到新状态。所以在旧约中有地上的许诺，在新约中有天上的许诺；因为这属于天主的仁慈，使人不可认为任何属地的幸福——无论何种——能由任何人获得，除非来自万物的创造主。但如果为求得地上的幸福而敬拜天主，那是奴隶的敬拜，属于婢女的子女；但如果为求天主本身，以在永生中天主成为万物中的万有，则是自由的服事，属于自由妇女的子女——她是我们在天上的永恒母亲——她起初似乎是荒胎的，因为她没有显现的子女；但现在我们看到关于她的预言：\InnerQuote{不怀孕、不生养的女子，欢呼吧！未曾经过产痛的女人，高声歌唱吧！因为被遗弃者的子女比有夫之妇的子女更多。}\BibleRef{依 54:1}——即比那在某种形式上与法律结合、同其子女一起为奴的耶路撒冷更多。因此，在旧约时期，我们说圣神在那些当时就按照依撒格作了承诺之子女的人心中，不仅是辅助者——这些人认为这足以支持他们的意见——而且是德能的赐予者；而他们否认这一点，反而将德能归于他们的自由意志，这与那些诚心向天主呼号的圣祖们相反：\InnerQuote{上主，我的力量，我爱慕祢。}\par

% structure: p0029 source=h2 target=subsection reason=relative-heading-rank; base=h2

\subsection{第十四章 [V.]——关于先知与宗徒的正义的诬蔑。}

他们又说，\Quote{我们并非将所有宗徒或先知都定为完全圣洁，而是说他们与更坏的人相比不那么邪恶；天主赐予证义的所谓义德即是如此，以致正如先知称索多玛在犹太人面前被称义一样，我们也说圣人们在与罪犯相比时行了某种善。} 愿我们绝不如此说；但要么他们不能理解，要么不愿观察，要么为歪曲而假装不知我们所说的。因此，让他们听，或是他们自己，或是他们正试图欺骗的那些无经验无学问的人。我们的信仰——即大公信仰——区分义人与不义之人，不是凭行为之法，而是凭信仰之法，因为义人因信德而生活。由此区分产生的结果是：一个人若一生不杀人、不偷盗、不作伪证、不贪他人之物、孝敬父母、贞洁到断绝一切肉身交合（甚至夫妻交合）、慷慨施舍、忍耐屈辱；不仅不剥夺他人之物，甚至被人夺去也不追讨；或甚至变卖所有财产分给穷人，自己一无所有——这样的品行看似可赞，但若对天主无真确的大公信仰，仍必由此生归入定罪。另一人却因由爱德行事的正确信德而有善工：在婚姻的诚实中保持节欲，虽不像前一人那样完全戒绝，而是偿付和回报肉身的债务，不仅为生子也为享乐而行房，但只与妻子，这是宗徒容许已婚者作为可宽恕的；不能那么耐心忍受屈辱，而是被报复的欲望激起怒气，但为了能说\BibleQuote{如同我们宽免我们的欠债者}，在被请求时宽恕；拥有私产，从中施舍一些，但不如前一人那么慷慨；不夺取他人之物，却为争取自己的权利而争讼（虽依教会审判而非民事审判）：这人在道德上似乎远不如前者，但因他对天主有正确信德（他凭此生活，并在一切过犯中自责，在一切善工中赞美天主，将羞愧归于自己，将光荣归于天主，并从祂领受罪的赦免和对善行的爱），他将由此生得救，进入与基督一同为王的行列。因此，若不是因信德？虽无善工不能救人（因为它不是无效的信德，而是由爱德运行的），但即使罪也因它被解除，因为义人因信德生活；没有它，即使看似善工也变为罪：\BibleQuote{凡不出于信德的都是罪。} \BibleRef{罗 14:23} 由于这巨大差异，虽然毫无疑问持守贞洁优于夫妻贞洁，但一位甚至再婚的妇女，若她是公教徒，优于一位异端的发愿贞女；并非因前者在天国中更好，而是因后者根本不在那里。我们描述为道德更佳的前一人，若有真信德，则超越后者，尽管二人都在天国；但若缺失信德，则他被后者超越，以致他自己根本不在那里。\par

% structure: p0031 source=h2 target=subsection reason=relative-heading-rank; base=h2

\subsection{第十五章——宗徒与先知的成全}

既然所有义人，无论较古者还是宗徒都凭对基督耶稣我们主的正确信德而生活；且他们的信德伴有如此圣洁的道德，以致虽然在今世不能有来世那样的全德，但来自人性软弱的一切罪孽都能因他们信德本身的虔诚而不断消除：由此产生的结果是，与天主将定罪的恶人相比，必须称这些人为\Quote{义人}，因为藉着虔诚的信德，他们与那些恶人截然相反，以致宗徒呼喊：\BibleQuote{信者与不信者有什么干系？} \BibleRef{格后 6:14} 但显然，伯拉纠派这些现代异端，自以为是对圣人们虔诚的敬爱者和赞美者，因他们不敢说圣人们有不完美的德行；然而那位蒙选的器皿承认了这一点，他思考自己仍在何种状态，以及那败坏的身体拖累灵魂，说：\BibleQuote{这不是说我已经达到，或已经成为成全的；弟兄们，我不以为自己已经得着了。} \BibleRef{斐 3:12} 而稍后，那曾否认自己是成全的人说：\BibleQuote{所以我们凡是成全的人，都当存这样的心。} \BibleRef{斐 3:15} 为表明按今生的尺度，有一定的成全，并且即使有人知道自己尚未成全，这也应归于那成全。因为古时圣洁的司铎中有谁更成全、更卓越？而天主命令他们首先为自己的罪献祭。新约子民中有谁比宗徒更圣洁？然而主命令他们在祈祷中说：\BibleQuote{宽免我们的罪债。} 因此，所有背负这可朽坏的肉身、在这今生软弱中叹息的虔诚人，只有一个希望：\BibleQuote{我们有耶稣基督作中保，祂是义者，为我们的罪作补赎。} \BibleRef{若一 2:1}\par

% structure: p0033 source=h2 target=subsection reason=relative-heading-rank; base=h2

\subsection{第十六章【VI】——关于基督内之罪的歪曲}

他们没有一个正义的辩护人，那些（即使这是唯一的区别）与义人完全且广泛地区分开来的人。我们绝不可以说，正如他们自己诽谤性地断言的那样，这位正义的辩护人\Quote{因肉身的必要而说了谎言}；但我们说，祂在罪过的肉身上，以罪过的肉身相似，论及罪，祂定罪了罪。他们或许不明白这一点，被诬蔑的欲望所蒙蔽，并不知道在圣经中罪这个名字有多少种用法，就把罪归给了基督。因此，我们断言基督既没有罪——无论灵魂还是身体都没有；并且，祂取了罪恶肉身的相似，论及罪，祂定罪了罪。这个由宗徒在某种程度上隐晦地提出的断言，有两种解释方式：要么事物的相似通常以它们所相似的事物之名来称呼，这样宗徒可能有意将这罪恶肉身的相似称为\Quote{罪}；要么按法律，为罪的祭献被称为\Quote{罪}，所有这些都是基督肉身的预像，基督的肉身是真正的、唯一的为罪的祭献——不仅是为那些在圣洗中全部被洗净的罪，也为那些后来因这一生的软弱而潜入的罪，为此普世教会每日在祈祷中向天主呼求：\Quote{宽免我们的罪债}，而它们因那独一无二的为罪的祭献而被宽免，宗徒按法律说话，毫不犹豫地称那祭献为\Quote{罪}。此外，还有他更清楚的一段话，没有任何双重歧义：\Quote{我们替基督求你们，与天主和好吧。祂替我们成了罪，祂本不认识罪，好叫我们在祂内成为天主的正义。}\BibleRef{格后 5:20}因为我上面提到的段落\Quote{论及罪，祂定罪了罪}，因为没有说\Quote{论及祂的罪}，任何人都可以理解为祂说祂在犹太人的罪上定罪了罪；因为在他们钉死祂的罪上，祂流出了祂的血以赦免罪。但这一段，天主使基督自己成了\Quote{罪}，祂原不认识罪，在我看来，没有比理解为基督成了为罪的祭献，并因此被称为\Quote{罪}，更合适的了。\par

% structure: p0035 source=h2 target=subsection reason=relative-heading-rank; base=h2

\subsection{第十七章——他们关于在来世遵守诫命的诽谤。}

谁能容忍他们指责我们说：\Quote{我们声称复活后我们将有这样的进步，以至于人们能够开始遵守他们在此世不愿遵守的天主诫命}？因为我们说在那里将完全没有罪，没有任何与罪之私欲的斗争；难道他们自己敢否认这一点吗？智慧和对天主的认识，那时在我们内得以圆满；在主内有如此喜乐，以致那是完满而真实的平安——除了那如此远离真理以至于连得到真理都不能的人，谁会否认呢？但这些事不会存在于诫命中，而是存在于那些应在此世遵守的诫命的赏报中；的确，忽视这些诫命不会导致那里的赏报。但在此世，天主的恩宠赐予遵守祂诫命的意愿；如果在这些诫命中有什么没有完全遵守，祂因我们在祈祷中所说的\Quote{愿祢的旨意奉行}以及\Quote{宽免我们的罪债}而予以宽恕。在此世，诫命是我们不得犯罪；在那里，赏报是我们不能犯罪。在此世，诫命是我们不顺从罪的私欲；在那里，赏报是我们没有罪的私欲。在此世，诫命是\Quote{愚昧的人，你们要明白；糊涂的人，你们何时才能明智？}在那里，赏报是完满的智慧和圆满的知识。宗徒说：\BibleQuote{因为我们现今藉着镜子观看，模糊不清，到那时就要面对面了；现今我所认识的，只是局部的，那时我就要完全认识，如同我被完全认识一样。}\BibleRef{格前 13:12}在此世，诫命是\Quote{你们要因上主我们的助佑而欢跃}以及\Quote{义人们，你们要因主而喜乐}；在那里，赏报是以完满而不可言喻的喜乐而喜乐。最后，在诫命中记载：\BibleQuote{饥渴慕义的人是有福的}；但在赏报中：\BibleQuote{因为他们要得到饱饫。}\BibleRef{玛 5:6}请问，他们从哪里得到饱饫，除非他们饥渴慕义的那义？那么，谁如此厌恶——不仅远离神圣的感知，也远离人的感知——以至于说在一个人身上，当他饥渴慕义时，能有他满得饱饫时那样的义？但当我们饥渴慕义时，如果对基督的信德在我们心中警醒，我们所信的是饥渴慕义什么，岂不正是基督？\BibleQuote{因为祂由天主成了我们的智慧、正义、圣化与救赎，好使经上记载的：凡夸耀的，应因主而夸耀。}\BibleRef{格前 1:30}因为我们只是相信祂，而不看见祂，所以我们饥渴慕义。因为只要我们还在身体内，我们就远离主；因为我们藉信德行走，而非藉见到。但当我们看见祂，并确实达到见到时，我们将以不可言喻的喜乐而喜乐；那时我们将充满正义，因为现在我们以虔敬的渴望对祂说：\Quote{当祢的荣耀显现时，我将得到满足。}\par

% structure: p0037 source=h2 target=subsection reason=relative-heading-rank; base=h2

\subsection{第十八章——正义的圆满与完全的安全感，连保禄在此生也不曾有。}

然而，那种骄傲——我不敢说是无耻，简直是疯狂——尚未等同于天主的众天使，却自以为已经拥有与天主的众天使同等的义德；它没有考虑到那样一位伟大而圣洁的人，他无疑曾渴望渴求义德的完满，当他因启示的伟大而不愿受抬举时；然而，为了不让他受抬举，他未被留给自己选择和意志，反而接受了\Quote{一根刺在肉里，撒殚的使者来击打我；为此我三次求主使它离开我，主对我说：我的恩宠为你够了，因为能力在软弱中得以成全。}\BibleRef{格后 12:7}什么能力？岂不就是那不被抬举的能力？谁怀疑这属于义德？天主的众天使确实拥有这种义德的完满，因为他们常常见到圣父的面，从而也见到整个圣三，因为他们通过圣子在圣神中看见。但没有什么比这启示更崇高，然而在那些欢乐者的凝视中，没有一位天使需要一个撒殚的使者来击打他，以免如此伟大的启示使他受抬举。宗徒保禄确实尚未拥有那种美德之完满，也尚未等同于天主的众天使；但他身上有抬举自己的软弱，这软弱必须被撒殚的使者在击打中遏制，以免他因启示的伟大而受抬举。因此，尽管最初的抬举使撒殚堕落，但那位至高的医生，深知如何善用甚至邪恶之事，却从撒殚的使者中，针对抬举的祸害，施加了一种有益虽痛苦的治疗，正如甚至用蛇来制作解毒剂以对抗蛇毒一样。那么，\Quote{我的恩宠为你够了}是什么意思？岂不是说，免得你因屈服而败在撒殚使者的击打下？\Quote{能力在软弱中得以成全}又是什么意思？岂不是说，在这软弱的处所中，美德得以完满，以致在软弱本身的临在中，抬举得以被抑制？这软弱确实将在未来的不朽中被治愈。因为那健康若仍需治疗，甚至来自撒殚使者的击打，又怎能称为完美呢？\par

% structure: p0039 source=h2 target=subsection reason=relative-heading-rank; base=h2

\subsection{第十九章——论此生人的义德在何种意义上被称为完美。}

由此得出，现今在义人身上的美德被称为完美，达到这一点：其完美包括对不足本身的正知和谦卑的承认。就这软弱而言，那小小的义德按其尺度是完美的，当它甚至理解自己缺乏什么时。因此，宗徒称自己既完美又不完美——不完美，在于他想到了为那义德所缺的，他仍饥渴其满全；完美，则在于他不羞于承认自己的不完美，并勉力向前以求达到。正如我们可以说行路者是完美的，其行程已良好地推进，但他的意图并未实现，除非他实际到达。因此，当他说\Quote{按照法律中的义德，我是无可指摘的}时，他立刻补充说：\Quote{凡先前对我有益的事，我因基督都看作损失。不但如此，我也将一切看作损失，因以认识我主基督耶稣为至宝；我因祂，将一切都损失了，看作粪土，为赢得基督，并能在祂内，不持有我自己在法律中的义德，而是由于信基督的义德，即天主基于信德的义德。}\BibleRef{斐 3:6}请看！宗徒当然没有说谎：\Quote{按照法律中的义德，他是无可指摘的}；然而那些对他有益的事，他为了基督而抛弃，看作损失、伤害、粪土。不但是这些，还有他之前提到的所有其他事物；不是因为任何一种知识，而是如他自己所说：\Quote{以认识我主基督耶稣为至宝}，这无疑他现今仍在信德中拥有，但尚未在眼见中。因为那时基督的知识将成为至宝，当祂被启示出来，所信的被看见时。由此，他在另一处这样说：\Quote{因为你们已经死了，你们的生命与基督一同藏在天主内；当基督——你们的生命显现时，你们也要与祂一同在光荣中出现。}\BibleRef{哥 3:3}同样，主自己说：\Quote{爱我的人，必蒙我父爱他，我也要爱他，并将我自己显示给他。}\BibleRef{若 14:21}因此，若望福音者说：\Quote{亲爱的，现在我们已是天主的子女，将来如何，还未显明；但我们知道，当祂显现时，我们要相似祂，因为我们要看见祂实在怎样。}\BibleRef{若一 3:2}那时基督的知识将成为至宝。因为现今，它仿佛隐藏于信德中；但在眼见中，它尚未显为至宝。\par

% structure: p0041 source=h2 target=subsection reason=relative-heading-rank; base=h2

\subsection{第二十章——保禄为何轻看法律中的义德。}

因此，蒙福的保禄抛弃了他以前那些基于自己正义的成就，视之为\Quote{损失}和\Quote{粪土}，为\Quote{赢得基督，并在他内被找到，并非拥有自己的正义，即法律而来的正义}。为何他自己的正义是法律而来的呢？因为那法律是天主的法律。除了玛尔基雍、摩尼以及诸如此类的害虫，谁会否认这一点呢？既然那是天主的法律，他却称那法律而来的正义是\Quote{他自己的}，并且他不愿拥有它，反而把它当作\Quote{粪土}抛弃。为什么这样呢？正是因为我在上文已经指出的：那些不认识天主的正义，而企图建立自己的正义，且不服从天主的正义的人，是处于法律之下。\BibleRef{罗 10:3}因为他们以为凭自己意志的力量就能履行法律的命令，并因骄傲而自满，不转向助佑恩宠。因此，文字叫他们死\BibleRef{格后 3:6}——或者公开地，因为他们明知自己行不出法律所命令的；或者自以为行得出，实际上却没有以出于天主的属神之爱来做。这样，他们要么明显是恶人，要么是自欺的义人——要么在公开的不义中被明断，要么在虚伪的正义中愚蠢地自高。因此，令人惊奇但真实的是，律法的正义不是由法律本身或法律中的正义所成全的，而是由恩宠圣神中的正义所成全的。因为律法的正义成全在那些不随从肉体、只随从圣神的人身上，正如经上所记。但按法律中的正义，宗徒说自己是在肉体上、而非在圣神上无可指摘；并且他说法律而来的正义是他自己的，而非天主的。所以必须明白，律法的正义不是按法律中的或法律来的正义，即按人的正义而被成全的，而是按恩宠圣神中的正义，即按天主的正义——人从天主得来的正义——而被成全的。这可以更清楚简洁地表述如下：律法的正义成全时，不是法律命令，人仿佛凭自己的力量服从；而是圣神助佑，人的自由意志（却因天主的恩宠而得以自由）去实行。因此，律法的正义在于命令什么是中悦天主的、禁止什么是不中悦天主的；而法律中的正义在于服从文字，此外不寻求天主为圣善生活提供任何帮助。因为当他说\Quote{并非拥有自己的正义，即法律而来的正义，而是由于基督的信德而来的正义}时，他加上\Quote{即从天主来的正义}。因此，那正是天主的正义；骄傲的人不认识这正义，就企图建立自己的正义；它之所以被称为天主的正义，不是因为天主借它成为正义的，而是因为人从天主得到了它。\par

% structure: p0043 source=h2 target=subsection reason=relative-heading-rank; base=h2

\subsection{第二十一章——正义在此生绝不能达至圆满。}

如今，按这天主的正义，即我们从天主得来的正义，信德借爱德行事。但信德行事，为使人以某种方式达到他如今虽未看见却相信的那一位；当他看见那一位时，那时在信德中模糊如同镜中的，终将成为面对面的看见；那时连爱德本身也将达至圆满。因为愚蠢至极的是说：天主在未看见时被爱，与看见时被爱是同样的程度。此外，如果在此生中（正如任何虔诚之人所不怀疑的），我们越爱天主就越发正义，那么谁能怀疑，当天主之爱圆满时，虔敬而真实的正义也将达至圆满？那时，法律将被完全遵守，一无所缺；按宗徒所说，法律的圆满就是爱。因此，当他说\Quote{并非拥有自己的正义，即法律而来的正义，而是由于耶稣基督的信德而来的正义，即从天主来的、基于信德的正义}之后，他接着说\Quote{为认识他和他的复活的大能，以及与他同受苦难}。\BibleRef{斐 3:9}这一切在宗徒身上尚不圆满和完美；但他仿佛在道路上奔跑，奔向它们的圆满和完美。因为他在另一处说\Quote{我如今只认识一部分，到那时我将完全认识，如同我被完全认识一样}，他岂已完美地认识了基督？\BibleRef{格前 13:12}他岂已完美地认识了他复活的大能？因为他还需要在将来肉体复活时更充分地体验。他岂已完美地认识了与他同受苦难？因为他尚未为他经历死亡的痛苦。最后，他加上说\Quote{为使我能达到从死者中的复活}。\BibleRef{斐 3:11}然后他说\Quote{这不是说我已经获得，或者已经达至圆满}。那么，他承认自己尚未获得的、尚未达至圆满的，岂不正是那从天主来的正义？这正是他所渴望的，他不愿拥有自己的、法律而来的正义。因为他当时正谈论此事，并且为了抵抗天主的恩宠的仇敌（基督曾为此被钉十字架）而说了这些话；这些仇敌的族类中也有这些人。\par

% structure: p0045 source=h2 target=subsection reason=relative-heading-rank; base=h2

\subsection{第二十二章——人性正义与完美的本质。}

因为从他开始说这些话的地方，他如此开始：\BibleQuote{你们要提防狗，提防邪恶的工人，提防妄自行割的人。因为我们才是真割损，我们是在圣神内事奉天主}——或如一些抄本所记\BibleQuote{事奉天主圣神}，或\BibleQuote{事奉天主之神}——\BibleQuote{并在基督耶稣内夸耀，而不信赖血肉。}\BibleRef{斐 3:2} 这里显然他是针对犹太人说的，他们按肉体遵守法律，力图建立自己的义德，被文字杀死，而不被圣神赋予生命，并自夸，而宗徒和所有恩许的儿女却在基督内夸耀。然后他加上：\BibleQuote{虽然我也可以信赖血肉。若别人自以为可以信赖血肉，我更可以。}\BibleRef{斐 3:4} 并列举所有按肉体有光荣的事，最后说：\BibleQuote{就法律上的义德而言，是无可指摘的。} 当他说他视这一切为损失、为废物、为粪土，为获得基督，他加上了我正在讨论的这段话：\BibleQuote{并且在他内被找到，不是有我自己因法律而来的义德，而是有因基督的信德而来的义德，即那出于天主的义德。} 他承认他尚未获得这义德的圆满，这义德只会在对基督的卓越认识中实现，为此他说一切都是损失；因此他承认自己尚未达到完美。\BibleQuote{但我向前直跑，} 他说，\BibleQuote{为获得那为基督耶稣所获得的。}\BibleRef{斐 3:12} \BibleQuote{为获得那为基督耶稣所获得的}与\BibleQuote{为认识，如同我被认识}大致相同。\BibleQuote{弟兄们，} \BibleQuote{我并不以为我已经获得了；我只专注一件事：忘记背后，努力面前，向着标杆直跑，为得天主在基督耶稣内从高天召我来得的奖赏。}\BibleRef{斐 3:13} 这句话的语序是：\BibleQuote{我只专注一件事。} 这\Quote{一件事}正是主对玛尔大所告诫的，他说：\BibleQuote{玛尔大，玛尔大！你为了许多事操心忙碌，其实需要的只有一件。}\BibleRef{路 10:41} 宗徒渴望获得这似乎在路上的一件，他说他向前直跑，为得从高天召来的奖赏。因为谁会在想要获得他所宣称正在追求的那件时迟延呢？那时他将有与圣天使同等的义德，当然，不会有撒殚的使者来击打他们，以免他们因启示的伟大而高傲。然后，他告诫那些自以为已拥有那义德的圆满的人：\BibleQuote{所以，我们凡是成熟的人，应当持这种态度。}\BibleRef{斐 3:15} 似乎他说：如果按这短暂生命的容量，我们在某种程度上是成熟的，让我们明白这也属于那成熟：我们感知到我们在那将显于基督的天使般的义德上尚未完美。\BibleQuote{若你们在什么事上，} 他说，\BibleQuote{有别的想法，天主也要将这事启示给你们。}\BibleRef{斐 3:15} 如何启示呢？除非给那些在信德之路上行走和前进的人，直到那漫游结束，他们来到实际的看见。为此他接着说：\BibleQuote{然而，我们已到达什么程度，就该照着什么程度行走。}\BibleRef{斐 3:15} 然后他总结应当提防那些本章开头所论及的人：\BibleQuote{弟兄们，你们要一同效法我，并留意那些照着我们的榜样而行的人。因为有许多人，我多次告诉你们，现今也含泪告诉你们，他们的结局是丧亡，}\BibleRef{斐 3:16} 等等。这些人正是他开头所说的：\BibleQuote{你们要提防狗，提防邪恶的工人}及其后的话。因此，所有敌视基督十字架的人，都是力图建立自己的义德，即由法律而来的义德，——其中只有文字命令，而圣神不成就——不顺服于天主的法律。因为如果那些属于法律的成了继承者，信德就成了空虚。\BibleQuote{如果义德是由于法律，那么基督就白死了；这样，十字架的绊脚石就除去了。} 因此，那些主张义德由于法律的人就是基督十字架的仇敌，因为法律的作用是命令，而不是帮助。但天主的恩宠，藉着主耶稣基督，在圣神内，帮助我们的软弱。\par

% structure: p0047 source=h2 target=subsection reason=relative-heading-rank; base=h2

\subsection{第二十三章——没有基督恩宠的信德就没有真正的义德。}

因此，按法律内的义德生活，却无基督恩宠的信德——如宗徒所说，他无可指摘地生活——的人，必须被视为没有真正的义德；这并非因为法律不真实不圣洁，而是因为愿意服从那命令的条文，而没有使生命复活的圣神，仿佛凭自由意志的力量，这并不是真正的义德。但义人因信德而生活的义德，由于人从恩宠的圣神由天主获得，才是真正的义德。尽管在某些义人身上，按此生的能力，这义德并非不当地被称为完美，但比起天使平等的伟大义德，仍属微小。尚未拥有这义德的人，一方面，就他已有的而言，他说自己是完美的；就他仍欠缺的而言，他说自己是不完美的。但显然，那较低的义德构成功绩，而较高的义德成为赏报。因此，人不追求前者，便不能达到后者。所以，在人的复活之后，否认将有义德的圆满，并认为那生命的身体中的义德将如同这死亡的身体中所能有的义德，是极其愚昧的。但千真万确的是，人们不会在那里开始履行他们在此处不愿服从的天主诫命。因为将有最完美义德的圆满，但不是人追求所命令的、逐步努力达到圆满的圆满；而是在眨眼之间，如同死者的复活本身一样，因为完美义德的伟大将作为赏报赐给在此处服从了诫命的人，而不会作为有待完成的事命令给他们。但我这样说：他们履行了诫命，是要我们记念，这些诫命本身就包含那祈祷，圣洁的恩许儿女每日真诚地以它祈祷说：\Quote{愿祢的旨意承行}\BibleRef{玛 6:10}，并说\Quote{宽免我们的罪债}\BibleRef{玛 6:12}。\par

% structure: p0049 source=h2 target=subsection reason=relative-heading-rank; base=h2

\subsection{第24章 [VIII.]——伯拉纠异端的三个主要主张}

因此，当伯拉纠派被这些及类似的真理证言和话语所迫，不否认原罪；不说我们赖以成义的天主恩宠不是白白赐予的，而是按我们的功绩；不说在会死的人中，无论多么圣洁和行善，甚至经过重生的洗涤后，直到结束此生，赦罪对他并非必要——因此，当他们被迫不作出这三个断言，并借此使相信他们的人离弃救主的恩宠，将自高自大的人诱入骄傲，投入魔鬼的审判时：他们引入其他问题的迷雾，在这些迷雾中，他们的不敬——在更单纯的人眼前，无论是由于他们较迟钝或在圣经上较不熟练——得以隐藏。这些是受造物的赞美、婚姻的赞美、法律的赞美、自由意志的赞美、圣人的赞美的模糊问题；好像我们中的任何人习惯于贬低这些事物，而不是以应有的赞美来宣扬一切以光荣造物主和救主。但即使受造物也不愿如此被赞美，以至于不愿被医治。婚姻越是受赞美，就越少归因于它那可羞耻的肉体贪欲，这贪欲不是出于父，而是出于世界；并且婚姻在人类中发现它，而非制造它；因为，此外，它实际上存在于许多人中而没有婚姻，并且如果没有人犯罪，婚姻本身可能没有它。法律，圣洁、公正而良善，既不是恩宠本身，没有恩宠也不能凭它做任何正确的事；因为法律不是为赐予生命而赐下的，而是因过犯而添上的，为要叫所有被罪定罪的人都被圈住，使那因信耶稣基督的恩许赐给相信的人。\BibleRef{迦 3:22} 被俘虏的自由意志，除了犯罪，毫无用处；对于义德，除非被神圣地解放和援助，否则毫无用处。同样，所有的圣人，无论是从古时的亚伯尔直到若翰洗者，还是从宗徒们自己直到此时，并从此直至世界终结，都应在主内被赞美，而非在自己内被赞美。因为即使那些早期圣人的声音是：\Quote{我的灵魂要因主而受赞美}。后来圣人的声音是：\Quote{因天主的恩宠，我成为今日的我}\BibleRef{格前 15:10}。对所有圣人来说，都属于：\Quote{那夸耀的，应在主内夸耀。}所有圣人的共同宣认是：\Quote{如果我们说我们没有罪，就是欺骗自己，真理不在我们内}\BibleRef{若一 1:8}。\par

% structure: p0051 source=h2 target=subsection reason=relative-heading-rank; base=h2

\subsection{第25章 [IX.]——他表明天主教的观点是介于摩尼教徒和伯拉纠派之间的中间立场，并反驳两者}

但既然，在这些我所举出的五个特定点上，他们寻求藏身之处，并从中编织诽谤，但他们被圣经所抛弃和定罪，他们以为可以用摩尼教徒这个可恨的名字来吓唬那些他们所能影响的人，以免他们的耳朵在反对他们最歪曲的教导时顺从真理；因为毫无疑问，摩尼教徒亵渎地谴责这五个教条中的前三个，说无论是人类受造物、婚姻还是法律，都不是由至高无上的真天主所制定的。但他们不接受真理所说的话，即罪恶起源于自由意志，一切罪恶，无论是天使的还是人的，都来自自由意志；因为他们宁愿相信，在他们背离天主时，邪恶的本性一直是邪恶的，并与天主同永远。此外，他们还尽其所能地诅咒圣祖和先知。这就是现代异端者的想法，以为通过提出摩尼教徒的名字，就能逃避真理的力量。但他们逃避不了；因为真理追赶他们，并同时推翻了摩尼教徒和白拉奇主义者。因为当一个人出生时，就他作为人而言，他有某种善，这谴责了摩尼教徒，赞美了造物主；但就他带有原罪而言，这谴责了白拉奇主义者，并坚持救主是必需的。因为即使那个本性被称为\Emph{可治愈的}，它也驳斥了两种教导；因为一方面，如果它是健康的，就不需要医治，这反对白拉奇主义者；另一方面，如果其中的恶是永恒不变的，它根本不能被治愈，这反对摩尼教徒。此外，对于婚姻，我们赞美它为天主所订立，但不说肉体的私欲应该归因于它，这既反对白拉奇主义者（他们赞美这种私欲本身），又反对摩尼教徒（他们把它归因于一个外来的邪恶本性，而实际上它是对我们本性的一种偶然之恶，不应通过分离于天主来分离，而应通过天主的仁慈来治愈）。此外，我们说法律是圣洁、公义、良善的，不是为了称义恶人，而是为了定骄傲之人的罪，为了过犯的缘故——这既反对摩尼教徒（因为按照宗徒，法律被赞美），又反对白拉奇主义者（因为按照宗徒，没有人因法律成义）；因此，为了使那些被文字所杀的人（即被命令行善的法律因过犯而定罪的人）活过来，恩宠之神白白地带来帮助。此外，我们说意志在作恶上是自由的，但行善必须由天主的恩宠使之自由，这反对白拉奇主义者；但说它起源于先前不是恶的东西，这反对摩尼教徒。此外，我们以应得的赞美在天主内尊崇圣祖和先知，这反对摩尼教徒；但我们说，即使对他们，无论他们多么正义和取悦天主，主的救赎也是必需的，这反对白拉奇主义者。因此，大公信仰发现两者（也像其他异端者一样）都与它对立，并通过圣经的权威和真理的光照定罪两者。\par

% structure: p0053 source=h2 target=subsection reason=relative-heading-rank; base=h2

\subsection{第二十六章 [X.]——白拉奇主义者仍通过引入关于灵魂起源的不必要问题来寻找藏身之处。}

白拉奇主义者确实为他们藏身之处的乌云增添了一个关于灵魂起源的不必要问题，目的是通过其他问题的晦涩来扰乱明显的事物，从而建立一个藏身之处。因为他们说\Quote{我们保卫灵魂的连续传播与罪恶的连续传播。}我不知道他们在哪里或何时读过这个，无论是在维护大公信仰之人的演讲或著作中，因为尽管我找到了一些天主教作者关于这个主题的论述，但真理的辩护尚未针对那些人进行，也没有任何回答他们的焦虑。但我要说，按照圣经，原罪是如此明显，而它在婴儿身上通过重生的洗礼被消除，这是由大公信仰的如此古老和权威所证实，由教会如此清晰一致的见证所闻名，以至于任何人关于灵魂起源的探究或肯定，如果与此相反，就不可能是真的。所以，无论谁，关于灵魂或任何其他晦涩问题，建立任何可能摧毁这个真实、最牢固、最著名的真理的体系，无论他是教会的儿子还是敌人，都必须被纠正或避开。但让这本书就此结束，以便后续的内容有另一个开端。\par

\SourceRecord{Translated by Peter Holmes and Robert Ernest Wallis, and revised by Benjamin B. Warfield.；Nicene and Post-Nicene Fathers, First Series；Vol. 5.；Edited by Philip Schaff.；Buffalo, NY: Christian Literature Publishing Co.,；1887.；<http://www.newadvent.org/fathers/15093.htm>.}

% structure: p0001 source=h1 target=section reason=page-title

\section{\WorkTitle{驳斥伯拉纠派两函（卷四）}}

在前面的各卷中，\Person{奥古斯丁}已驳斥了针对公教徒的诬蔑，现在他进而揭露伯拉纠派第二函剩余部分中隐藏的诡计，即其教义的五个要点——即：赞美受造物、赞美婚姻、赞美法律、赞美自由意志、赞美圣人。伯拉纠派恶毒地夸耀，他们在这几点上的分歧不仅与摩尼教徒有关，更与公教徒有关。因此，这五点可归结为他们的三重谬误：即前两点是否认原罪；中间两点是主张恩宠依功绩而赐予；第五点是声称圣人今生未犯罪。\Person{奥古斯丁}指出，摩尼教和伯拉纠派这两种异端都与公教信德相悖，同样令人憎恶。公教信德宣称：第一，由善的天主所创造的本性是善的，但它需要救主，因为原罪从第一人的过犯传给了众人；第二，婚姻是善的，确实由天主设立，但与婚姻相伴的私欲却是恶的，是由罪带来的；第三，天主的法律是善的，但它只能显明罪，而不能除去罪；第四，自由意志确实是人本性固有的，但现在它却被奴役，除非由恩宠释放，否则无力行义；第五，无论旧约还是新约的圣人，确实享有真实的义德，但这义德并非完美的，也不至于使他们完全无罪。最后，他引用了\Person{西彼连}和\Person{盎博罗削}为公教信德所作的见证，有的关于原罪，有的关于恩宠的帮助，还有的关于现世义德的不完美。\par

% structure: p0003 source=h2 target=subsection reason=relative-heading-rank; base=h2

\subsection{\Emph{第一章 [I.]——伯拉纠派的五种遁词。}}

在我已考虑并答复了那些问题之后，他们重复了那封我已驳斥的信中所含的同样内容，但方式不同。先前，他们是作为反对我们、似乎是我们错误地思想的东西提出来的；但后来，作为解释自己的思想，以相反的一面提出了同样的东西，并加上了两件他们之前没有提及的事——即\Quote{他们说圣洗对各年龄都是必要的}，以及\Quote{死亡经由亚当传给我们，而不是罪}。这些事本身也必须在适当的地方予以考虑。因此，由于在我刚刚完成的前一卷书中，我提到他们提出了五个方面的障碍，其中隐藏着他们敌视天主恩宠和公教信德的教义——即赞美受造物、赞美婚姻、赞美法律、赞美自由意志、赞美圣人——我认为更适宜对所有他们坚持的、并且反对我们的那些观点进行一个总体的区分，并指出其中哪些属于那五点中的哪一点，从而使我的回答通过这种区分更加清晰和简洁。\par

% structure: p0005 source=h2 target=subsection reason=relative-heading-rank; base=h2

\subsection{\Emph{第二章 [II.]——赞美受造物。}}

他们在以下陈述中完成了对受造物的赞美，就其所涉及的人类的当前问题而言：\Quote{天主是一切出生者的创造者，人的子女是天主的工程；一切罪孽并非来自本性，而是来自意志。}他们在这赞美中联系了\Quote{他们说圣洗对各年龄都是必要的，以便}，即\Quote{受造物本身可以被收纳为天主的子女；并不是因为它从父母那里获得了任何必须在重生之水中被净化的东西。}他们又在赞美中加上\Quote{他们说主基督就其婴孩时期而言，未被任何罪的污点沾染}；因为他们断言，祂的肉身由于与所有婴孩共享的本性，而非由于祂自己卓越和独特的恩宠，而完全没有罪的玷污。也属于这一点的是，他们引入了\Quote{灵魂的来源}这个问题，从而试图使所有婴孩的灵魂都与基督的灵魂平等，主张他们同样也没有沾染任何罪的污点。为此，他们还声称\Quote{从亚当传给其余人类的，除了死亡之外，没有别的恶}，他们又说\Quote{死亡并不总是恶的，例如对殉道者而言，它反而为了赏报；而且，并不能因为各种人中都复活的身体的解体，就称死亡为善或恶，而是源于人类自由的功过差异。}他们在这封信中写了这些关于赞美受造物的内容。\par

他们确实按照圣经赞美婚姻，\Quote{因为主在福音中说：\InnerQuote{从起初创造人的，造了他们一男一女，并说：你们要生育繁殖，充满大地。}}虽然这在福音的那段经文中没有记载，但在法律中是有记载的。他们还加上\Quote{所以，天主所结合的，人不可拆散。}\BibleRef{玛 19:4} 我们也承认这些是福音的话。\par

在赞美法律时，他们说：\Quote{旧的法律，按照宗徒的话，是圣洁、公义、良善的；对于那些遵守其诫命、藉着信德正义生活的人，如先知、圣祖和所有圣人，能够赐予永生。}\par

在赞美自由意志时，他们说：\Quote{自由意志并没有丧失，因为主藉先知说：\InnerQuote{如果你们愿意听从我，就必享用地上的美物；如果你们不愿意，也不听从，刀剑必吞噬你们。}\BibleRef{依 1:19} 同样，恩宠辅助任何人的善意，但并不将德行的渴望注入不情愿的心中，因为天主不偏待人。}\par

他们在赞美圣徒时隐藏自己，说：\Quote{圣洗完美地更新人，因为宗徒见证说，藉着水的洗涤，教会从外邦人中成为圣洁无瑕的；\BibleRef{弗 5:26} 圣神也在古代帮助虔诚的灵魂，正如先知对天主说：\InnerQuote{愿祢善神引导我走上正途；} 而且所有的先知、宗徒和圣徒，无论新约还是旧约，天主为之作证的都是义人，不是与恶人相比，而是按德行的标准；而且在将来，善行和恶行都有报应。但没有人能在那时履行此地可能轻忽的诫命，因为宗徒说：\InnerQuote{我们必须出现在基督的审判台前，好使每人按他所行的，或善或恶，领受与身体相关的事。}} \BibleRef{格后 5:10}\par

在所有这一切中，无论他们说什么赞美受造物和婚姻的话，都试图将我们带回这一点：即没有原罪；无论赞美法律和自由意志，都指向这一点：恩宠不帮助无功劳的人，因此恩宠不再是恩宠；无论赞美圣徒，都指向这一点：圣人的现世生命似乎没有罪，他们不必祈求天主赦免他们的债务。\par

% structure: p0012 source=h2 target=subsection reason=relative-heading-rank; base=h2

\subsection{第三章——公教徒赞美本性、婚姻、法律、自由意志和圣人，其方式既谴责白拉奇派也谴责摩尼派。}

让每一个以大公的思想对这些不敬神和该受惩罚的教义感到战栗的人，在这三部分划分中，避开这五重错误的潜伏之处和陷阱，并如此谨慎地分辨，既要远离摩尼派，又不倾向白拉奇派；同样，既要与白拉奇派分离，又不与摩尼派联合；或者，如果他已被其中之一所束缚，就应当如此挣脱两者的手，而不冲向另一者。因为它们似乎是彼此对立的；因为摩尼派通过诋毁这五点而显明自己，而白拉奇派则通过赞美它们来隐藏自己。因此，凡按大公信仰的准则在由肉体和灵魂的善受造物所生的人中显耀造物主（因为这是摩尼派所不愿承认的），同时承认因第一个人的罪所传递给他们的败坏，连婴儿也需要一位救主（因为这是白拉奇派所不愿承认的）——这样的人，无论他是谁，都谴责并避开两者。他如此区分可耻情欲的恶与婚姻的祝福，既不像摩尼派那样指责我们出生的根源，也不像白拉奇派那样赞美我们紊乱的根源。他如此持守法律是通过梅瑟由圣善、公义、良善的天主所赐予的（这是摩尼与宗徒相反所否认的），同时说它显示罪却不除去罪，命令义却不给予义（这又是白拉奇与宗徒相反所否认的）。他如此肯定自由意志，说天使和人的恶不是始于某种我不知道的永远邪恶的本性（这根本不是本性），而是始于意志本身，这推翻了摩尼派的异端，然而即使如此，被俘虏的意志也不能凭自身呼吸到健康的自由，除非依靠天主的恩宠，这推翻了白拉奇派的异端。他如此赞美天主内的天主圣人们，不仅是在基督降生成人之后及其后的圣人们，甚至包括古代的那些人——摩尼派敢于亵渎他们——却相信圣人们关于自己的忏悔，而非白拉奇派的谎言。因为圣人的话是：\Quote{如果我们说我们没有罪，就是欺骗自己，真理就不在我们内。} \BibleRef{若一 1:8}\par

% structure: p0014 source=h2 target=subsection reason=relative-heading-rank; base=h2

\subsection{第四章——白拉奇派与摩尼派论对受造物的赞美。}

鉴于这些事实，对那些十字架基督的仇敌、神圣恩宠的反对者——新异端者来说，如果他们死于自己的另一种瘟疫，那么他们看似脱离了摩尼教的错误，这又有什么益处呢？他们赞美受造物，说\Quote{善的天主是那些出生者的创造者，万物藉祂而造成，人的儿女是祂的工程}，而摩尼教却说他们是黑暗王子的工程；在这两者之间，或两者之中，天主在婴孩中的受造物正在灭亡，这于他们有何益呢？因为两者都拒绝让这受造物藉基督的体血得救——一方是因为他们摧毁那体血本身，仿佛祂根本没有从人或在人身上取过这些；另一方是因为他们断言婴孩中没有任何邪恶，使他们需要藉这体血的圣事得救。在这两者之间，婴孩中的人性受造物，有着良好的起源，却有着败坏的传承：它为它的善事承认一位最卓越的创造者，为它的恶事寻求一位最仁慈的救赎者；它以摩尼教为它恩惠的诽谤者，以伯拉纠派为它邪恶的否认者，两者均为迫害者。虽然婴孩没有发言的能力，却以它无声的目光和隐藏的软弱向两者的不虔敬虚妄说话：对一方说\Quote{相信我是由那创造美好事物的祂所创造的}；对另一方说\Quote{容许我被那创造我的祂所医治}。摩尼教说\Quote{这婴孩中除了要得救的善灵之外，什么都没有；其余的，}不属于善的天主，而属于黑暗王子，\Quote{应予摒弃}。伯拉纠派说\Quote{当然，这婴孩中没有什么需要拯救的，因为我们已经证明整个是安全的}。两者都说谎；但是，现在仅仅指控肉体的人，比那赞美整个人而显出残忍的人更可容忍。但摩尼教因亵渎天主——整个人的创造者——而不帮助人的灵魂；伯拉纠派因否认原罪而不让神圣恩宠来帮助人的婴孩。因此，是公教信仰使天主施怜悯，因为它谴责两种有害的教义，从而来帮助婴孩得救。它对摩尼教说：\Quote{请听宗徒呼喊：\BibleQuote{难道你们不知道你们的身体是住在你们内的圣神的宫殿吗？} \BibleRef{格前 6:19} 并相信善的天主是身体的创造者，因为圣神的宫殿不能是黑暗王子的工程。}它对伯拉纠派说：\Quote{你们所看的这个婴孩\InnerQuote{是在罪孽中怀胎的，母亲在罪中孕育了它。}为什么，仿佛在辩护它没有任何恶事时，你们不允许它藉怜悯得救？没有一个人是不染污秽的，即使是地上的生命只有一天的婴孩。容许这些可怜者获得罪的赦免，藉着那一位，祂无论在小事上还是大事上，都不可能有什么罪。}\par

% structure: p0016 source=h2 target=subsection reason=relative-heading-rank; base=h2

\subsection{第5章。——在伯拉纠派意见中的特殊优势是什么？}

那么，他们声称\Quote{所有的罪不是源于本性，而是源于意志}，并凭这判断的真理抵抗摩尼教——他们说邪恶的本性是罪的原因——这于他们有何益呢？由于他们不愿承认原罪，虽然原罪本身也源于第一人的意志，他们却使婴孩带着罪愆离开身体。他们宣称\Quote{圣洗为所有年龄都是必要的}，而摩尼教说它为每个年龄都是多余的，而他们说在婴孩中，就罪的赦免而言，这是虚假的，这于他们有何益呢？他们坚持\Quote{基督的肉身}（摩尼教争辩说要么根本不是肉身，要么是幻像的肉身）不仅是真实的肉身，而且\Quote{灵魂本身没有被任何罪污所玷污}，同时其他婴孩被他们（伯拉纠派）放在与祂的婴孩期同等的位置上，带着相等的纯洁，以至于那肉身与这些婴孩比较起来似乎不保持自己的圣洁，而这些婴孩从那肉身得不到任何救恩，这于他们有何益呢？\par

% structure: p0018 source=h2 target=subsection reason=relative-heading-rank; base=h2

\subsection{第六章。——不仅死亡，而且罪也借着亚当进入了我们。}

在这一点上，确实，当他们说\Quote{死亡由亚当传给了我们，而非罪}时，他们并没有将摩尼教徒当作对手：因为摩尼教徒也否认原罪来自第一人——起初身体和灵魂都是纯洁正直的，后来因自由意志而堕落——随后罪与死亡一起传给并继续传给所有人；但他们说，肉体从起初就是邪恶的，是由一个恶神并与一个恶神一起创造的；而一个善的灵魂——即天主的一部分——由于它先前被束缚在饮食中的玷污的功绩，进入了人，并因此通过交媾被束缚在肉体的锁链中。这样，摩尼教徒与伯拉纠派一致认为，不是第一人的罪责传给了人类——既不是藉着肉体（他们说肉体从未善过），也不是藉着灵魂（他们断言灵魂带着它自己在肉体之前被污染的玷污的功绩进入人的肉体）。但伯拉纠派怎么说\Quote{只有死亡藉着亚当传给了我们}呢？因为如果我们死是因为他死了，但他死了是因为他犯了罪，那么他们说惩罚没有罪责就传了下来，而无辜的婴孩承受了不义的惩罚，因为他们得到了死亡却没有死亡的功过。公教信仰知道天主与人之间唯一的中保，基督耶稣之人，祂屈尊为我们承受了死亡——即罪的惩罚——却没有罪。正如祂独自成为人子，好使我们藉着祂成为天主的子女，同样祂独自代替我们承担了没有应得惩罚的死亡，好使我们藉着祂获得不应得的恩宠。因为对于我们，没有善是应得的，同样对于祂，没有恶是应得的。因此，向那些祂即将赐予不应得生命的人显扬祂的爱，祂愿意为他们承受不应得的死亡。伯拉纠派试图使中保的这一特殊特权归于无效，这样主旨中就不再有什么特殊了——如果亚当因自己的罪责承受了应得的死亡，而婴孩从他那里没有带任何罪责，却要承受不应得的死亡。因为死亡确实将许多善事赐予善人，有人甚至恰当地论证了\Quote{论死亡的益处}；然而，由此除了天主的仁慈还能宣告什么呢？既然罪的惩罚被转化为有益的用途？\par

% structure: p0020 source=h2 target=subsection reason=relative-heading-rank; base=h2

\subsection{第七章。—— \Quote{在祂内众人皆犯罪}是什么意思？}

但这些人如此说话，意图将人从宗徒的话语中扭转至他们自己的思想。因为宗徒说：\Quote{藉着一个人，罪进入了世界，死亡藉着罪也进入了，这样死亡就传到了所有人身上，}\BibleRef{罗 5:12} 他们却要那里被理解的不是\Quote{罪}传了过去，而是\Quote{死亡}。那么，接下来所说的\Quote{为此，众人都犯了罪}是什么意思？因为宗徒或者是在说，在他所说的\Quote{藉着一个人，罪进入了世界}的那个\Quote{一个人}中，众人都犯了罪，或者是在那个\Quote{罪}中，或者肯定是在\Quote{死亡}中。因为他说的是\Emph{在其中}（使用阴性代词形式）而不是\Emph{在其中}（使用阳性代词）众人都犯了罪，但这并不应困扰我们，因为在希腊语中，\Quote{死亡}是阳性。那么，让他们选择他们愿意的——或者是在那个\Quote{人}中众人都犯了罪，正如所说，因为他犯罪时，众人都在他内；或者是在那个\Quote{罪}中众人都犯了罪，因为那是所有人在普遍中所做的事，所有出生的人都将继承；或者剩下的一种说法：在那个\Quote{死亡}中众人都犯了罪。但我并不清楚这如何能被理解。因为众人都是在罪中死去，而不是在死亡中犯罪；因为当罪先发生时，死亡跟随其后，而非死亡先发生，罪跟随其后。因为罪的刺就是死亡——也就是说，刺的打击造成了死亡，而不是死亡用来打击的刺。正如毒药，若被喝下，就被称为死亡的杯，因为由那杯造成了死亡，并非因为杯是由死亡造成，或由死亡给予。但如果\Quote{罪}不能由宗徒的那些话被理解为\Quote{在其中众人犯了罪}，因为在希腊语中，这封书信是从那里翻译过来的，\Quote{罪}是以阴性表达的，那么剩下的理解就是众人是在第一个\Quote{人}中犯的罪，因为当他犯罪时，众人都在他内；并且罪是由生育从他那里衍生出来，除非由重生，否则不得赦免。因为圣依拉略也是这样理解所写的\Quote{在其中众人犯了罪}；他说：\Quote{在其中}，也就是在亚当内，\Quote{众人都犯了罪。}然后他加上：\Quote{显然，众人在亚当内犯了罪，如同在整体中；因为他自己被罪败坏，他所生的一切人都生而处于罪之下。}当他写下这些时，依拉略毫无歧义地指出了我们应当如何理解\Quote{在其中众人犯了罪}这句话。\par

% structure: p0022 source=h2 target=subsection reason=relative-heading-rank; base=h2

\subsection{第八章。—— 死亡因罪临到众人。}

但是，同一宗徒说我们藉基督与天主和好，是为了什么原因呢？除非因为我们成了祂的仇敌。而这是什么？不就是罪吗？因此先知也说：\Quote{你们的罪使你们与天主隔绝。}\BibleRef{依 59:2} 因着这隔绝，中保被派遣，为除去世界的罪，我们因罪而与天主分离成了仇敌，好使我们得以和好，从仇敌成为儿女。宗徒必定是在谈论此事；因此他插入了\Quote{罪藉着一人进入了世界}这句话。因为他先前说：\Quote{但是天主向我们证明了自己的爱，当我们还是罪人的时候，基督为我们死了。何况现在，我们既因祂的血而成义，更要藉着祂从义怒中得救。因为如果我们作仇敌的时候，藉着祂圣子的死，得与天主和好，那么，和好之后，我们更要因祂的生命得救了。不但如此，我们还藉着我们的主耶稣基督，藉着祂我们现已获得了和好，在天主内欢跃。}然后他接着说：\Quote{因此，就如罪藉着一人进入了世界，死亡藉着罪也进入了世界，这样死亡就传到了众人，因为众人都犯了罪。}白拉奇派为何要回避这事？如果和好对所有人都是必须的，那么因着使我们成为仇敌的罪，人人都犯了罪，好使我们需要和好。这和好是在重生之洗中，在基督的肉和血中，没有这些，连婴孩也不能在自己内有生命。因为就如有一人因罪而死亡，也有一人因义德而带来生命；因为\Quote{如在亚当内众人都死了，照样在基督内众人都要复活；}\BibleRef{格前 15:22} 并且\Quote{就如因一人的罪，众人都被定罪；同样，因一人的义行，众人都获得生命的成义。}\BibleRef{罗 5:18} 有谁听了这些话，还以如此顽固的邪恶不敬，胆敢争辩说死亡藉着亚当传给了我们，却没有罪呢？除非他们是天主恩宠的反对者和基督十字架的仇敌——他们若固执于此，结局就是丧亡。但是，为了对付他们蛇一般的诡诈——他们试图用这种诡诈败坏朴实的心灵，藉着赞美受造物，把他们从信德的纯朴中引诱出来——这样说就够了。\par

% structure: p0024 source=h2 target=subsection reason=relative-heading-rank; base=h2

\subsection{第九章 论婚姻的赞美}

此外，关于对婚姻的赞美，他们反对摩尼教徒——摩尼教徒将婚姻归于黑暗之君，而非真实和良善的天主——这些人抵抗真虔诚的话，说：\Quote{主在福音中说：\InnerQuote{起初，造了他们男和女，又说：你们要生育繁殖，充满大地。所以，天主所结合的，人不可拆散}}? \BibleRef{玛 19:4} 这对他们有什么益处呢？藉着真理引诱人走向谬误？因为他们说这话，是为了使人们认为婴孩出生时毫无瑕疵，因此不需要藉着基督与天主和好，因为他们没有原罪——这原罪使得和好对于所有人都是必须的，藉着那一位没有罪而来到世界的，正如所有人的仇恨是由那一位引起的，因他的罪进入了世界。天主教徒为了人性得救而相信这一点，并不有损于对婚姻的赞美；因为对婚姻的赞美是两性正当的结合，而不是对恶习的邪恶辩护。因此，这些人藉着对婚姻的赞美，试图将人们从摩尼教徒那里吸引到自己这边来，他们不过是想要改变病症，而不是医治它。\par

% structure: p0026 source=h2 target=subsection reason=relative-heading-rank; base=h2

\subsection{第十章 论法律的赞美}

再者，在赞美法律时，他们虽然反对摩尼教徒，说了真话，但这于他们有何益处呢？他们想引导人们脱离那种观点，转而接受他们针对公教徒所持的错误立场。因为他们说：\Quote{我们承认，就连旧法律，依照宗徒所言，也是圣洁、公义、良善的，并且它能使遵守其诫命、凭信德正义生活的人获得永生，就如先知、族长和诸圣一样。} 他们用这些精心措辞的话赞美法律，以反对恩宠；诚然，那法律虽是公义、圣洁、良善的，却并不能使所有那些天主的人获得永生，惟有在基督内的信德才能。因为这信德藉着爱德行事，不凭那使人死的文字，而凭那使人活的圣神，天主的恩宠——那法律如同启蒙师一样，藉着威慑使人不敢犯法，引领人归向它，以便那法律本身不能赐予的，能赐予人。因为针对他们所说\Quote{法律能使遵守其诫命的先知、族长和诸圣获得永生}的话，宗徒回答说：\BibleQuote{如果成义是藉着法律，那么基督就白白地死了。} \BibleRef{迦 2:21} \BibleQuote{如果嗣业是藉着法律，便不再是出于恩许。} \BibleRef{迦 3:18} \BibleQuote{如果那些属于法律的人成为嗣业者，信德便落了空，恩许也就失效了。} \BibleRef{罗 4:14} \BibleQuote{但没有人能藉着法律在天主面前成义，这是明显的，因为义人因信德而生活。} \BibleRef{迦 3:11} \BibleQuote{法律并不出于信德，而是：\InnerQuote{行这些事的人，必因这些事而生活。}} \BibleRef{迦 3:12} 宗徒从法律引用的这段证言，就现世生命而言，人们因害怕丧失现世生命，通常行法律的事，而非信德的事；因为违反法律的人，按同一法律的规定，应由百姓处死。或者，若必须从更高的意义来理解，即\Quote{行这些事的人，必因这些事而生活}是针对永生写的，那么法律的权能是如此表达的：人自身软弱无力，不能实行法律所命令的，便应寻求天主恩宠的帮助，这恩宠出于信德，因为连信德本身也是因天主的仁慈赐予的。因为信德是这样获得的，正如天主按各自的分量赐给了每人信德的度量。\BibleRef{罗 12:3} 因为人若不从自身拥有信德，而是领受了刚毅、爱德和节制的神，那么那位外邦人的导师亲自说：\BibleQuote{因为我们没有领受怯懦的神，而是领受了刚毅、爱德和节制的神。} \BibleRef{弟后 1:7} ——无疑，也领受了信德之神，关于此神他说：\BibleQuote{既然我们拥有同一信德之神。} \BibleRef{格后 4:13} 所以，法律确实说：\Quote{行这些事的人，必因这些事而生活。} 但要行这些事，并因这些事而生活，需要的不是颁行此事的法律，而是获得此事的信德。然而，这信德，为能配得领受这些事，它本身也是白白赐予的。\par

% structure: p0028 source=h2 target=subsection reason=relative-heading-rank; base=h2

\subsection{第十一章——白拉奇派人士说法律本身是天主的恩宠}

但恩宠的仇敌在暗中设下更隐蔽的陷阱，以更猛烈地反对同一恩宠，其方式莫过于赞美那无疑值得赞美的法律。因为他们藉着各种说话方式和所有辩论中变换的词语，希望人们理解法律就是\Quote{恩宠}——即是说，我们从主天主那里得到知识的帮助，从而知道必须做的事，而不是得到爱的感动，使我们认识后能以圣善的爱去实行，这爱才是真正的恩宠。因为没有爱的法律知识使人自高自大，并不建树，正如同一宗徒最明确地所说：\BibleQuote{知识使人自高自大，惟有爱德才能建树。} \BibleRef{格前 8:1} 这话与\BibleQuote{文字使人死，神使人活} \BibleRef{格后 3:6} 相似。因为\Quote{知识使人自高自大}对应\Quote{文字使人死}；而\Quote{爱德建树}对应\Quote{神使人活}；因为\BibleQuote{天主的爱，藉着所赐与我们的圣神，已倾注在我们心中了。} \BibleRef{罗 5:5} 因此，法律的知识使人成为骄傲的犯法者；但藉着爱德的恩赐，他乐于成为法律的实行者。那么，我们并不是藉着信德废了法律，而是坚固法律，\BibleRef{罗 3:31} 因为法律藉着威吓引向信德。这样，法律确实激起愤怒，好让天主的仁慈将恩宠赐给那受惊而悔改、并藉着我们的主耶稣基督——天主的智慧——实践法律正义的罪人；关于这智慧经上记载：\BibleQuote{她的舌上带着法律和仁慈。} \BibleRef{箴 3:16} ——法律用以威吓，仁慈用以帮助——法律藉着祂的仆人，仁慈藉着祂自己——法律好比\Person{厄里叟}打发去复活寡妇儿子的棍杖，却未能救活他，\BibleQuote{因为如果颁行了能赐生命的法律，那么成义诚然是出于法律了，} \BibleRef{迦 3:21} 但仁慈好比\Person{厄里叟}本人，他预像基督，藉着复活死人来象征新约的伟大圣事。\par

% structure: p0030 source=h2 target=subsection reason=relative-heading-rank; base=h2

\subsection{第十二章 [VI.]——论赞美自由意志}

此外，他们为反对\OriginalTerm{摩尼教徒}{Manicheans}而称赞\OriginalTerm{自由意志}{free will}，引用先知的话说：\BibleQuote{如果你们愿意并听从我，就必享用这地的美物；但如果你们不愿意，也不听从我，刀剑必吞灭你们}\BibleRef{依 1:19}。这对他们有何益处呢？事实上，他们与其说是在反对摩尼教徒，不如说是在赞颂自由意志以反对\OriginalTerm{天主教徒}{Catholics}。因为他们希望所说的\BibleQuote{如果你们愿意并听从我}被理解为：在先行意志中已有了随后恩宠的功劳；这样，恩宠就不再是恩宠，因为它不再被白白赐予，而是当作债务偿还。但如果他们这样理解\BibleQuote{如果你们愿意}这句话，以至于承认那预备甚至这个善意本身的，就是那位关于他说过\BibleQuote{意志由主预备}的那一位\BibleRef{箴 8:35}，那么他们就会以天主教徒的身份使用这个见证，不仅会战胜摩尼教的古老异端，也不会建立\OriginalTerm{白拉奇主义者}{Pelagians}的新异端。\par

% structure: p0032 source=h2 target=subsection reason=relative-heading-rank; base=h2

\subsection{\Emph{第十三章——天主的旨意是恩宠的效果。}}

他们称赞同一个自由意志，说\Quote{恩宠帮助每个人的善意}，这对他们有何益处呢？如果他们不将功劳归于善意（现在，他们按照欠债而非按恩宠的工资付给是应得的），而是理解和承认，连那随后被帮助的善意本身，若没有恩宠先行的预备，也不能存在于人之中，那么这话就会毫无疑虑地被当作公教精神接受。因为如果没有天主的先行的仁慈，人如何能有善意呢？因为那正是\BibleRef{箴 8:35}所说的\BibleQuote{意志由主预备}。但当他们说了这些话：\Quote{恩宠也帮助每个人的善意}，并随即加上\Quote{但并不将德性的爱注入抵抗的心}时，若是由那些我们已知其意的人所说，这话或可理解。因为，对抵抗的心，首先是靠天主恩宠本身获得对神圣召唤的听从，然后，在这颗不再抵抗的心里，德性的渴望被点燃。然而，在任何人按天主而行的一切事上，天主的仁慈先于他。他们不愿这样，因为他们选择做白拉奇主义者，而非天主教徒。因为骄傲的不敬虔很喜欢这一点：即一个被迫承认是由主赐予的，似乎不是白白给予他，而是偿还给他；以致于沉沦之子（而非恩许之子）自以为他们自己使自己成为善的，而天主则给他们这些已由自己成为善的人，按他们的工作偿还应得的酬报。\par

% structure: p0034 source=h2 target=subsection reason=relative-heading-rank; base=h2

\subsection{\Emph{第十四章——圣经为恩宠所作的见证。}}

因着这骄傲，他们心里的耳朵被塞住了，以致听不见：\BibleQuote{因为你有什么不是领受的呢？}\BibleRef{格前 4:7}；听不见：\BibleQuote{离了我，你们什么也不能做}\BibleRef{若 15:5}；听不见：\BibleQuote{爱是从天主来的}\BibleRef{若一 4:7}；听不见：\BibleQuote{天主按量分配了信德}\BibleRef{罗 12:3}；听不见：\BibleQuote{风随意向哪里吹}\BibleRef{若 3:8}，以及\BibleQuote{凡受天主圣神引导的，都是天主的子女}\BibleRef{罗 8:14}；听不见：\BibleQuote{除非蒙我父恩赐，谁也不能到我这里来}\BibleRef{若 6:65}；听不见\Person{厄斯德拉}所写的：\BibleQuote{我们祖先的上主是可称颂的，因为祂将这意愿放在王心里，要装饰耶路撒冷上主的殿宇}\BibleRef{厄上 8:25}；听不见主通过\Person{耶肋米亚}所说的话：\BibleQuote{我必把我的敬畏放在他们心里，使他们不离开我；我必眷顾他们，使他们得到好处}\BibleRef{耶 32:40}；尤其是通过\Person{厄则克耳}先知所说的那段话，天主在其中清楚地表明：祂并非因人的任何功德而使人成义、使人服从祂的命令，而是以善报恶，并为自己的缘故而非为他们而做。因为祂说：\BibleQuote{吾主上主这样说：以色列家族，我这样做，不是为了你们，而是为了我的圣名，就是你们所到的列邦中，被亵渎的名。我要在列邦中显扬我的大名，就是你们在他们中间被亵渎的名。当我在你们身上，显示了我的圣洁，在他们眼前时，列邦就要承认我是上主——吾主上主说。我要把你们从列邦中领出，从各国聚集你们，领你们回到你们的地域。我要在你们身上洒清水，洁净你们，洗净你们的一切不洁和各种偶像。我还要给你们一颗新心，在你们五内注入一种新的精神；从你们肉体内去掉那颗硬心，给你们一颗易感的心。我要把我的神放在你们五内，使你们遵行我的法令，遵守我的诫命并付诸实行。}\BibleRef{则 36:22-27}过了不多几句，同一位先知又说：\BibleQuote{吾主上主这样说：我这样做并不是为了你们，以色列家族，使你们知道，你们应惭愧，因你们的行径而羞愧。吾主上主这样说：在我洗净你们的一切不洁之日，我又使城市有人居住，废墟得以重建，过去荒芜之地，如今在过路人眼前变成果园；荒凉、废置和毁坏的城市，如今都成了巩固的居所。四周剩余的异民，将知道我是上主，是我重建了废墟，培植了荒芜之地；我上主说了，也实行了。吾主上主这样说：我还要为此而允许以色列家族的请求，使他们多得人口，像羊群一样增多；人会像耶路撒冷在庆节日的羊群一样，那些荒废的城市也要充满人；这样他们就知道我是上主。}\BibleRef{则 36:32-38}\par

% structure: p0036 source=h2 target=subsection reason=relative-heading-rank; base=h2

\subsection{第十五章——从这些圣经经文证明恩宠是白白且有效的。}

腐烂的皮肤还有什么可令它自高自大，并能在夸耀时藐视夸耀于主？\BibleRef{格前 1:31} 它还有什么剩下，当它无论说了什么它所作的事，以致那先前从人而来的功绩起源于人之后，天主随后才作那人所应得的事——它将要得到回答，它将要受到抗议，它将要被反驳：\Quote{我这样做；但为了我圣名的缘故，不是为你们的缘故，我这样做——上主天主说}？\BibleRef{则 36:22} 没有什么比这更能推翻白拉奇派，当他们说天主的恩宠是按我们的功绩赐予的。其实，白拉奇本人也谴责了这一点，若不是因为纠正它，也是因为害怕东方审判官。没有什么比这更能推翻那些说\Quote{我们这样做，好使我们堪当那些天主以此实行的事}之人的骄傲。回答你的不是白拉奇，而是主自己：\Quote{我这样做，不是为了你们，而是为了我圣名的缘故。}\BibleRef{则 36:22} 你能从不善的心里行出什么善来？但为使你有一颗善心，祂说：\Quote{我要赐给你们一颗新心，在你们内放一新精神。} 你能说，我们要先遵行祂的正义，遵守祂的判断，并且这样做，好使我们配得祂所赐的恩宠吗？但你们这些恶人，能行什么善，你们怎能行那些善事，除非你们自己成为善人？然而，除了那说\Quote{我要眷顾他们，使他们为善}的祂，谁能使人成为善人呢？还有谁说\Quote{我要把我的精神放在你们内，使你们遵行我的正义，遵守我的判断，并且实行它们}？难道你们还不醒悟吗？你们还没有听见\Quote{我要使你们行走，我要使你们遵守，}最后，\Quote{我要使你们实行}吗？什么！你们还在自高自大吗？我们确实行走，不错；我们遵守；我们实行；但祂使我们行走、遵守、实行。这就是天主使我们为善的恩宠；这是祂预先行施的仁慈。那荒废、荒凉和掘起的地方，却还要被建造、耕种和加固，它们有什么功德？这些事是因为它们的荒废、荒凉、连根拔起的功绩吗？决不是。因为这类事物是恶的功绩，而那些恩赐是善的。因此，善的事物是为恶的事物赐下的——所以是白白的；不是债，因此是恩宠。主说：\Quote{我，}\Quote{上主。} 这样的话难道不约束你吗？人的骄傲啊，当你说，我做这些事好从主那里配得被建造和种植。你没有听见吗？\Quote{我这样做不是因你们的缘故；我，上主，建造了毁坏的城市，我种植了荒凉的土地；我，上主，说了，我做了，但不是为了你们，而是为了我圣名的缘故。} 谁使人数增多如羊，如圣羊，如耶路撒冷的羊？谁使这些荒凉的城市满有羊一样的人，除非是那继续说话的主：\Quote{他们要知道我是上主}？但祂用什么羊来充满城市，正如祂所应许的？那些祂找到的，还是那些祂创造的？让我们询问圣咏；看，它回答；让我们听：\Quote{请你们前来，让我们俯伏朝拜上主，在上主面前哭泣，因为祂是造我们的天主；祂是我们的天主，我们是祂牧场的子民，是祂手下的羊群。} 祂因此创造羊，用它们充满荒凉的城市。有什么稀奇，当对那一只羊，即教会——其成员都是人的羊——说\Quote{因为我是创造你的上主}。你向我夸耀什么自由意志，除非你成为羊，否则它不能自由地行正义？那么，祂既使人成为祂的羊，祂就释放人的意志去服从孝爱。\par

% structure: p0038 source=h2 target=subsection reason=relative-heading-rank; base=h2

\subsection{第十六章——为什么天主使一些成为羊，另一些不。}

然而，为何天主使这些人成为羊，而不使那些人成为羊呢？在祂那里，难道有偏袒人吗？这正是蒙福的宗徒对那些过于好奇而不合宜地提出此问题的人所作的回答：\BibleQuote{人啊！你是谁，竟敢向天主抗辩？受造的岂能对造它的说：你为什么这样造了我？}\BibleRef{罗 9:20}这正属于那深奥的事，连同一宗徒在试图探究时也多少感到畏惧，并喊道：\BibleQuote{啊，天主的智慧和知识的富饶，是多么深奥！祂的判断是多么不可测量，祂的道路是多么不可追踪！谁曾知道主的心意？或者谁曾作过祂的参谋？谁曾先给了祂，而应受祂的偿还呢？因为万物都出于祂、借着祂、归于祂；愿光荣归于祂，至于无穷之世。}愿那些在恩宠之前、甚至反对恩宠而辩护功绩的人，不要胆敢探究那不可深究的问题，他们想先给天主，好使自己再得偿还——当然，他们想先凭自由意志给出一些东西，好使恩宠作为报酬再赐予他们；愿他们明智地理解或忠信地相信，就连他们自以为先给出的，也是从祂那里领受的，因为万物都出于祂、借着祂、归于祂。但为何这一个人领受，而那一个人不领受？两人都不配领受，无论谁领受，都是不配地领受。——愿他们衡量自己的力量，不要探究过于强大之事。让他们知道天主并无不义就够了。因为当宗徒在雅各伯身上找不到任何功绩使他优于他的孪生兄弟时，他说：\BibleQuote{那么，我们该说什么呢？难道天主有不义吗？绝对不可！因为祂对梅瑟说：\InnerQuote{我要恩待谁，就恩待谁；我要怜悯谁，就怜悯谁。}因此，不在乎那愿意的，也不在乎那奔跑的，而在乎那显示仁慈的天主。}因此，即使这深奥的问题仍未解决，让祂自由的怜悯成为我们感恩的对象；然而，同一宗徒已如此解决了这个问题，他说：\BibleQuote{如果天主愿意显示祂的愤怒并彰显祂的权能，就以极大的忍耐包容那些准备丧亡的忿怒之器；并为了显明祂光荣的财富，赐予那些祂为光荣所预备的慈悲之器。}\BibleRef{罗 9:22}当然，除非是应得的，忿怒不会被偿还，以免天主有不义；但怜悯，即使赐予时并非应得，也不会使天主有不义。因此，让慈悲之器明白怜悯是如何白白地赐予他们，因为对忿怒之器——他们与慈悲之器有共同的毁灭之因和尺度——所偿还的忿怒是正义且应得的。这对那些因自由意志而试图摧毁恩宠之慷慨的人，已足够了。\par

% structure: p0040 source=h2 target=subsection reason=relative-heading-rank; base=h2

\subsection{论圣人的赞美}

事实上，在赞美圣人时，他们不会以那法利塞人的虚荣（仿佛满溢于自足和圆满）来驱使我们，而是以那税吏的热诚\BibleRef{路 18:10}，使我们饥渴慕义；但他们——反对否认圣洗的摩尼教徒——说\Quote{人在圣洗中得以完全更新}，并引用宗徒的见证：\BibleQuote{他作证说，借着水的洗涤，教会从外邦人中得以圣洁和没有瑕疵}\BibleRef{弗 5:26}，却以骄傲和邪僻的意思提出他们的论点，反对教会本身的祈祷。因为他们这样说，是为了使人相信在圣洗（其中完成所有罪的赦免）之后，教会不再有罪；然而，从日出之地到日落之地，教会的所有成员都相反他们向天主呼喊：\BibleQuote{赦免我们的罪债}\BibleRef{玛 6:12}。但如果他们在这件事上被询问关于他们自己，他们就找不到可回答的。因为若他们说没有罪，\BibleQuote{若望就回答他们，他们欺骗自己，真理不在他们内}\BibleRef{若一 1:8}但若他们承认自己的罪，既然他们希望自己是基督身体的成员，那么那身体（即教会）如何在此时就如他们所想的那样完全没有瑕疵和皱纹呢？如果它的成员真实地承认自己有罪。因此，在圣洗中，所有罪都被赦免，且借由那在圣言中水洗、教会就在基督内毫无瑕疵和皱纹地呈现出来；\BibleRef{弗 5:27}除非它受了洗，否则它无效地呼喊：\BibleQuote{赦免我们的罪债}，直到它被带入光荣中，那时它内绝对毫无瑕疵或皱纹。\par

% structure: p0042 source=h2 target=subsection reason=relative-heading-rank; base=h2

\subsection{圣人对自己本身的看法}

应当承认，\Quote{圣神即使在古代，}不仅\Quote{助佑善的倾向}——连他们自己都承认——甚至使它们成为善的，这是他们所不愿承认的。\Quote{所有那些天主为之作证的先知、宗徒或圣人，无论是福音时期的还是古代的，他们之所以是义人，并非相对于恶人而言，而是按照德行的标准，}这是无可置疑的。这一点反对了摩尼教人，因为他们亵渎圣祖和先知；但反对白拉奇主义者的却是：所有这些圣人，当他们仍生活在肉身内而被问及自身时，会以极其一致的声音回答说：\BibleQuote{如果我们说我们没有罪，便是欺骗自己，真理也不在我们内。}\BibleRef{若一 1:8}\Quote{但在未来的时期，}不可否认\Quote{善行和恶行都将得到报偿，并且没有人会在那里被命令去做他在此处所轻视的诫命，}而是完美的义德——在那里罪失去能力，圣人在这里如饥似渴地追求的义德——在此处于训导中被期望，在那里作为酬报而领受，因着赒济和祈祷的恳求；如此，在此处可能未完全遵守诫命之处，可因罪过的赦免而不受惩罚。\par

% structure: p0044 source=h2 target=subsection reason=relative-heading-rank; base=h2

\subsection{第十九章——白拉奇主义者的诡计。}

如果事情果然如此，那么让白拉奇主义者停止用他们对这五件事——即赞美受造物、赞美婚姻、赞美法律、赞美自由意志、赞美圣人——的极其阴险的赞扬，来假装他们想要将人，如同从摩尼教人的小网罗中拯救出来，以便将他们缠在自己的网中——即否认原罪；拒绝给婴儿基督医生的救助；说天主的恩宠是按照我们的功德赐予的，这样恩宠就不再是恩宠；以及说圣人在今生没有罪，从而使祂赐给没有罪的圣人的祈祷无效——因为藉此祈恳，所有罪过得赦免。他们用对这五件美事的欺骗性赞美，引诱粗心和无知的人陷于这三个邪恶的教义。关于这一切，我想我已经充分谴责了他们极其残忍、邪恶和傲慢的虚妄。\par

% structure: p0046 source=h2 target=subsection reason=relative-heading-rank; base=h2

\subsection{第二十章[VIII.]——古人对白拉奇主义者的证词。}

但是既然他们说\Quote{他们的敌人出于对真理的仇恨而盗用了我们的话，}并抱怨说\Quote{几乎在整个西方，一种既愚蠢又邪恶的教义被采纳，简单的教区主教们在没有主教会议的情况下被强迫签名支持这一教义，}——尽管基督的教会，无论是西方还是东方，都对他们言语的亵渎新颖感到战栗——我认为我有责任不仅使用神圣的正典圣经作为反对他们的见证——这我已经充分做了——而且还要从那些在我们之前以最广泛的名望和极大荣耀处理过这些圣经的圣人们的著作中提出一些证据。我并非要把任何辩论家的权威与正典书籍等同，好像对某一教理，一位天主教徒比另一位同样是天主教徒的思考更正确或更真实；而是为了提醒那些认为这些人在这些主题上的空谈之前，追随天主神谕的全体天主教教师从未如此说过的人，让他们知道，反对白拉奇异端者的后退的狂妄和毒害，我们维护的是真正古老确立的公教信仰。\par

% structure: p0048 source=h2 target=subsection reason=relative-heading-rank; base=h2

\subsection{第二十一章——白拉奇乌斯模仿西彼廉写了一本见证集。}

就连这些人的异端首领\Person{白拉奇}本人，也以应得的敬意提及最有福的\Person{西彼廉}——这位因殉道者的荣冠而最光荣的人，不仅在\Place{非洲教会}和\Place{西方教会}，也在\Place{东方教会}中，因名声的传扬及其著作的广泛传播而为人所知——当他撰写一部证言集时，宣称自己在效法他，说\Quote{他对\Person{罗马努斯}所做的，正如\Person{西彼廉}对\Person{基里诺}所做的一样}。那么，让我们看看\Person{西彼廉}对由一人进入世界的原罪有何看法。在\WorkTitle{论善功与施舍}的书信中，他如此说：\Quote{主在来临时，治愈了亚当所引入的这些创伤，并医治了蛇的旧毒，祂给健全的人颁布了法律，命他不要再犯罪，免得他犯罪后有更糟的事发生。我们因无罪的诫命而被限制和关闭在狭窄的空间内；人的软弱和脆弱没有任何出路，除非再次降来援助的天主仁慈，借着指出正义和仁慈的善工，开辟某种确保救恩的道路，以便我们通过施舍洗净随后沾染的任何污秽。}藉此证词，这位证人驳斥了他们的两个谎言——其一，他们说人类从亚当并未遗传什么需要藉基督治愈和医治的罪；其二，他们说圣人们在领洗后不再有罪。同样，在同一书信中，他说：\Quote{让每个人将魔鬼及其仆役——即丧亡和死亡之民——放在眼前，仿佛冲入中间，在基督——祂亲自临在并审判——的面前，以比较的考验挑战基督的子民，用这些话：\InnerQuote{我，代表你们所看到与我在一起的那些人，既没有接受拳打，也没有忍受鞭打，也没有背十字架，也没有流我的血，也没有以我的苦难和血的代价赎回我的家族；我既不向他们应许天国，也不在重新使他们恢复不死之后，将他们召回乐园。}}让白拉奇派回答并说出我们何时曾处于乐园的不死之境，以及我们如何被逐出那里，以致要被基督的恩宠召回。尽管他们可能无法找到在此情况下为自己的乖谬可以回答什么，但让他们观察\Person{西彼廉}如何理解宗徒所说的\BibleQuote{在他内众人都犯了罪}。并让脱离旧摩尼异端的白拉奇异端者，不要胆敢对一位公教徒提出任何诬蔑，以免被证明对古代殉道者\Person{西彼廉}犯了如此邪恶的错。\par

% structure: p0050 source=h2 target=subsection reason=relative-heading-rank; base=h2

\subsection{第二十二章 再论西彼廉}

因为他也在标题为\WorkTitle{论死亡}的书信中说了这样的话：\Quote{亲爱的弟兄们，天主的王国即将临近；生命的赏报、永生救恩的喜乐与永恒的欢愉，以及从前失落的乐园的拥有，正随着世界的消逝而到来。}同样，在同一书信中，他说：\Quote{让我们迎接那将我们各人分配到自己家中的日子，这日子将我们从这里攫取，从世界的陷阱中释放，并使我们恢复于乐园和王国。}此外，他在关于忍耐的书信中说：\Quote{让我们默想天主的审判，这审判早在世界和人类的开端，亚当——忘记诫命并违犯所赐法律者——便已领受。这样我们就会知道，在此生中我们应当多么忍耐，我们生在这样的状态中，以苦难和争战在此劳苦。因为，祂说：\InnerQuote{你听从了你妻子的声音，吃了我唯独禁止你吃的树上的果子，地因你的作为而受诅咒；你一生日日都要在忧愁和叹息中吃它；它要给你生出荆棘和蒺藜，你要吃田间的蔬菜；你必汗流满面才得糊口，直到你归回所从出的土地：因为你是尘土，你必归于尘土。}我们都在这判决的锁链中被捆绑束缚，直到死亡被消灭，我们离开这个世界。}而且，在同一书信中，他说：\Quote{因为，自从那第一次违犯诫命起，身体的力气与不死性一同离去，软弱随死亡而来，且除非不死性也已被领受，否则力气不能被领受，所以我们在这身体的脆弱和软弱中，必须始终奋斗和争战；而这奋斗和交锋唯有藉忍耐的力量才能承受。}\par

% structure: p0052 source=h2 target=subsection reason=relative-heading-rank; base=h2

\subsection{第二十三章 再论西彼廉}

在和六十六位同僚主教写给非都斯主教的信函中，他因割损律法而受咨询，问婴儿是否可在第八天之前受洗，此事被如此处理，仿佛出于天主上智，天主教会将要在许久之后反驳将要出现的伯拉纠异端者。因为咨询者对于孩子出生时是否继承原罪——他们可藉重生予以洗净——并无怀疑。因为基督徒信仰从未对此事有所怀疑，望之远矣。但他怀疑的是，重生的洗涤——他无疑认为原罪由此被除去——是否应在第八天之前施行。最蒙福的西彼廉回复此咨询说：\Quote{但关于婴儿的案例，你说他们不应在出生后第二或第三天受洗，而应遵守古割损律法，因此你认为出生者不应在第八天之内受洗成圣，我们全体在会议中持非常不同的意见。因为你所认为应采取的道路，无人赞同，我们反而都判断，慈悲天主的恩宠不应拒绝任何生于人者；因为主在祂的福音中说：\BibleQuote{人子来不是要毁灭人的性命，而是要拯救他们。} \BibleRef{路 9:56} 尽我们所能，我们必须努力，若可能，不使任何灵魂丧亡。} 稍后他又说：\Quote{我们谁也不应害怕天主所俯就创造的事物。因为婴儿虽刚出生，但它并非如此，以致任何人应在施恩和缔造和平时害怕亲吻它，因为在亲吻婴儿时，我们每个人都应为着敬虔之故，思量那近乎崭新的天主之手本身，我们仿佛正在亲吻那刚成形和新出生的人中所有的那手，当我们拥抱天主所造之物时。} 稍后他又说：\Quote{但若有什么能阻止人获得恩宠，他们的更重罪恶或许更阻止那些成熟、成长和年长者。但再说，若甚至对最大的罪人，以及那些先前大大得罪天主的人，当他们后来相信时，罪的赦免被赐予，并且无人被阻止领受圣洗和恩宠；那么，我们更应多么不敢阻止一个婴儿，他刚出生，并未犯罪，除非他按肉体出于亚当出生时，在最初出生时就染上了古死的传染；他正因此更容易获得罪的赦免，因为赦免给他的不是他自己的罪，而是别人的罪！}\par

% structure: p0054 source=h2 target=subsection reason=relative-heading-rank; base=h2

\subsection{第24章——向伯拉纠派提出的两难问题}

对于这些事，那些不仅是离弃者，而且是迫害天主恩宠的人，将说什么呢？他们将如何回应这些事？凭什么\Quote{乐园的拥有}被恢复给我们？若我们从未在那里，又如何被恢复至乐园？或者说，我们如何在那里，除非因为我们在亚当内？我们如何属于那对悖逆者所说的\Quote{审判}，若我们不承受悖逆者的伤害？最后，他认为婴儿甚至在第八天之前就应受洗，免得因\Quote{在最初出生时所染上的古死的传染}，婴儿的灵魂丧亡。若连信徒所生的孩子，在基督内重生并被从黑暗权势中拔出、迁入祂的国度之前，不被魔鬼所持有，他们又如何丧亡？谁说那些出生者的灵魂若不重生就会丧亡？不是别人，正是那位如此赞美造物主和受造物、工匠和作品，以致约束并纠正人情的恐惧——人们拒绝亲吻初离母胎的婴儿——而提出对造物主本身的崇敬，说在亲吻那种年龄的婴儿时，应思量那近乎崭新的天主之手！那么，他在承认原罪时，是否定本性或婚姻吗？他因将重生的洁净应用于从亚当而生而有罪的婴儿，就否认天主是所生者的造物主吗？因为他害怕任何年龄的灵魂丧亡，便与其同僚会议决定甚至在第八天之前就应藉圣洗圣事交付他们，他因此控告婚姻吗？实际上，对婴儿——无论生于婚姻或奸淫，但因是人生而为人——他宣称天主之手甚至配享平安之吻。那么，如果圣主教和最光荣的殉道者西彼廉能够认为婴儿的原罪必须以基督的医药治愈，而不否定对受造物的赞美，不否定对婚姻的赞美，为什么一种新奇的瘟疫，虽不敢称像他这样的人为摩尼教徒，却以为应将另一个人的过错归咎于坚持这些事的天主教徒，以掩盖自己的过错？于是，最受赞誉的圣经注释者，甚至在摩尼教瘟疫的丝毫污染触及我们的土地之前，毫无对天主作品和婚姻的指责，承认原罪——并未说基督沾染了任何罪的污点，也未将祂与别人出生时那有罪的肉体相比较，祂藉此有罪肉体的相似而赐予洁净的帮助；他也不被灵魂来源的模糊问题所阻，而不承认那些藉基督的恩宠获得自由的人返回乐园。他是否说死亡的条件从亚当传给人类而没有罪的传染？因为他说婴儿无论多么新鲜出生，都要藉圣洗提供帮助，并非为避免身体的死亡，而是为那由一人进入世界的罪，\BibleRef{罗 5:12}\par

% structure: p0056 source=h2 target=subsection reason=relative-heading-rank; base=h2

\subsection{第25章【IX.】——西彼廉关于天主恩宠的证词}

但现在，西彼廉在论及《主祷文》时，如何向这些人宣讲天主的恩宠，已明显可见。因为他说：\Quote{我们说，\BibleQuote{愿祢的名被尊为圣}，并非我们愿天主因我们的祈祷而被尊为圣，而是我们恳求祂，使祂的名在我们内被尊为圣。但天主怎能被尊为圣呢？因为祂自己就是圣者。然而，因为祂说：\BibleQuote{你们应是圣的，因为我是圣的}，我们祈求并恳求这一点：即我们已在圣洗中成为圣洁的人，能继续持守我们所开始成为的。}在同一封信的另一处，他说：\Quote{我们还加上并说：\BibleQuote{愿祢的旨意承行于地，如于天焉}，并非为使天主做祂所愿的，而是为使我们能做天主所愿的。因为谁抵抗天主，阻止祂行其所愿呢？但是，由于我们在一切事上都被魔鬼阻碍，不能以思想和行为服从天主，我们就祈祷并求天主旨意在我们内承行。为使它在我们内承行，我们需要天主的旨意，即祂的帮助和保护；因为没有人凭自己的力量强大，而是因天主的宽仁和仁慈而安全。}在另一处又说：\Quote{此外，我们求天主的旨意承行于天也承行于地，这两件事都关乎我们安全与救恩的完成。因为我们从地得到身体，从天得到神魂，我们自身就是地与天；因此，我们求天主的旨意在两者中承行，即在身体和神魂中。因为肉体与神魂之间有争斗，它们彼此不相合，每日相争；以致我们不能行我们所愿的，因为神魂寻求天上的和神圣的事，而肉体贪恋地上的和暂时的事。所以，我们祈求，因天主的帮助和援助，这两性之间能达致和谐；这样，当天主的旨意在神魂和肉体中都承行时，由祂重生的灵魂就能得以保存。宗徒保禄以他的话清楚地表明这事。他说：\BibleQuote{肉体的私欲相反圣神，圣神的私欲相反肉体；它们彼此相敌对，以致你们不能行你们所愿意的事。}}稍后他又说：\Quote{最亲爱的弟兄们，可以这样理解：既然主命令并教导我们甚至要爱仇敌，并为迫害我们的人祈祷，我们就应为那些仍是土地、尚未开始成为天上的人祈求，使天主的旨意也在他们身上承行，这旨意就是基督在保存和更新人性时所完成的。}又在另一处他说：\Quote{但我们求这饼日赐给我们，使我们这些在基督内、每日领受圣体作为救恩食粮的人，不致因某种更严重的罪的介入——即被阻止，如那些不领受、不共融的人，不能分享这天上的饼——而与基督的身体分离。}稍后，在同一篇论文中他说：\Quote{但当我们求不陷于诱惑时，我们被提醒自己的软弱和脆弱，我们如此祈求，是叫没有人傲慢地自夸；没有人骄傲狂妄地将任何事归于自己；没有人将宣认或受苦的光荣归于自己，当作自己的，因为主亲自教导谦卑说：\BibleQuote{你们应当警醒祈祷，免得陷于诱惑；心神固然愿意，肉体却是软弱的。}所以，当谦卑顺服的宣认在先，并将一切归于天主时，凡以敬畏和尊崇天主的心谦卑祈求的，都可得蒙祂的慈爱赐予。}此外，在\WorkTitle{致基利诺}的论文中（白拉奇希望自己在这部著作中显得是他的模仿者），他在第三卷中说：\Quote{我们不可夸耀任何事，因为没有任何事是属于我们自己的。}他附上神性见证以支持这一命题，其中特别加上了那使这类人的口尤其必须被堵住的宗徒话语：\Quote{\BibleQuote{你有什么不是领受的呢？既然是领受的，为什么你还要夸耀，好像不是领受的呢？}}同样在\WorkTitle{论忍耐}的书信中他说：\Quote{我们与天主共有这一德行。忍耐从他开始；它的光荣和尊荣从他而来。忍耐的起源和伟大出自天主，他是其作者。}\par

% structure: p0058 source=h2 target=subsection reason=relative-heading-rank; base=h2

\subsection{第二十六章——进一步引用西彼廉的教导}

那位在真理之言中教导众教会的圣洁而令人难忘的导师，是否否认人有自由意志，因为他将你们全部的公义生活归于天主？他是否指责天主的法律，因为他暗示人不能因法律成义，既然他宣称法律所命令的必须通过祈祷从主天主那里获得？他是否以恩宠的名义断言命运，因他说我们不可在任何事上夸耀，因为没有什么是我们自己的？他是否像这些人一样相信，圣神是以这样的方式帮助德行，仿佛它所帮助的那同一个德行源自我们自己，当他断言没有什么是我们自己的时，他在这方面提到宗徒说过：\Quote{你有什么不是你领受的呢？}并且说最卓越的德行，即忍耐，并非从我们开始、随后才由天主的圣神接受帮助，而是从祂本身取得源头，从祂取得起源？最后，他承认若非天主的恩宠，好的意图、对德行的渴慕、善的意向都不会在人心中开始，因为他说：\Quote{我们不可在任何事上夸耀，因为没有什么是我们自己的。}有什么比法律所说\Quote{不可崇拜偶像，不可奸淫，不可杀人}更要建立于自由意志之上呢？不，这些罪行及类似的事，若有人犯了，就被驱逐出基督身体的共融。然而，如果\Person{圣西彼廉}认为我们自己的意志足以不犯这些罪行，他就不会如此理解我们在主祷文中所说的\Quote{求你今天赏给我们日用的食粮}，以至于断言我们祈求\Quote{免得因某些可憎之罪的介入——因被阻止而戒除并不同领，以致不能分食天上的食粮——而与基督的身体分离}。让这些新异端者确实回答：在那些敌视基督之名的人身上，有什么善功先行？因为他们不但没有善功，反而有最坏的功过。然而，\Person{圣西彼廉}甚至如此理解我们在祷词中所说的\Quote{愿你的旨意奉行在人间，如同在天上}：我们也为那些在此方面被称为\Emph{地}的人祈祷。因此，我们不仅为那些不愿意的人祈祷，也为那些反对和抵抗的人祈祷。那么，我们祈求什么，岂不是使不愿意的变为愿意，反对的变为同意，抵抗的变为爱慕？并且藉着谁？岂非藉着那一位，关于祂所记载的：\Quote{意志由天主预备}？\BibleRef{箴 8:36}那么，那些人不愿——倘若他们不作任何恶，倘若他们行任何善——不夸耀自己，而夸耀于主，就让他们学习做天主教徒吧。\par

% structure: p0060 source=h2 target=subsection reason=relative-heading-rank; base=h2

\subsection{第二十七章 [X.]——\Person{西彼廉}关于我们自身成义不完善的证言。}

那么，让我们来看那第三点，这一点在这些人身上，对基督的每一个肢体和祂的整个身体来说，同样令人震惊——他们主张：在今世，或曾有完全无罪的义人。在这种自以为是中，他们最明显地与主祷文相矛盾，在祷文中，基督的所有肢体以诚实的心和日常的话语大声呼喊：\Quote{宽免我们的罪债。} 那么，让我们看看\Person{西彼廉}——在主内最荣耀的——对此有何看法？他不只为教导教会——当然不是摩尼教，而是天主教会——说了这些话，而且还将它们付诸书信并铭记于心。在关于\WorkTitle{善工与施舍}的书信中，他说：\Quote{因此，亲爱的弟兄们，让我们承认天主仁慈的有益恩赐；既然我们不可能没有良心的某种创伤，就让我们借由属灵的疗法来医治我们的创伤，以洗净并清除我们的罪。不要让任何人因自认有一颗纯洁无瑕的心而自我陶醉，以致依赖自己的无辜，认为无需将药物敷在自己的伤口上；因为经上记载：\BibleQuote{谁能夸口说，自己有一颗洁净的心？谁能夸口说，自己纯洁无罪？} \BibleRef{箴 20:9} 再者，\Person{若望}在他的书信中规定并说：\BibleQuote{如果我们说我们没有罪，就是欺骗自己，真理就不在我们内。} \BibleRef{若一 1:8} 但如果没有人能没有罪，而任何人若自称无罪，不是骄傲就是愚蠢，那么天主的仁慈是何等必要、何等慈祥！祂知道在那些已经痊愈的人身上仍然发现一些伤口，就在他们痊愈之后，又赐下有益的疗法来重新治愈和医治他们的伤口！} 再者，在同一篇论文中，他说：\Quote{既然每日在天主眼中犯的罪不会停止，那么每日用来洁净罪恶的祭献也不会停止。} 此外，在论\WorkTitle{死亡}的论文中，他说：\Quote{我们的争战是与贪婪、不洁、愤怒、野心争战；我们与肉体的恶习、世界的诱惑进行艰难而辛勤的角力。被围攻的人心，四面受魔鬼猛攻，几乎不能应对反复的攻击，几乎不能抵抗。如果贪婪被制服，情欲就升起；如果情欲被克服，野心就取而代之；如果野心被藐视，愤怒就激惹，骄傲就膨胀，酗酒就引诱；嫉妒破坏和谐，妒忌割断友谊；你被迫诅咒，这是天主法律禁止的；你被迫起誓，这是不合法的。灵魂每日遭受如此多的迫害，心灵因如此多的危险而疲惫；然而它却喜欢长久留在这里，处在魔鬼的武器中间，尽管我们更应该渴望并愿意借着更快地死亡而奔向基督。} 再者，在同一篇论文中，他说：\Quote{蒙福的宗徒\Person{保禄}在他的书信中规定说：\BibleQuote{对我来说，生活就是基督，死亡就是利益；} \BibleRef{斐 1:21} 他认为最大的利益是，不再被这个世界的网罗束缚，不再受制于肉体的罪和恶习。} 此外，在论\WorkTitle{天主经}中，在解释我们当说\Quote{愿你的名受显扬}时，他和其他事一起说：\Quote{因为我们需要每日的圣化，好使我们这些每日堕落的人，能借着不断的圣化洗去我们的罪。} 再者，在同一篇论文中，当他解释我们的\Quote{宽免我们的罪债}时，他说：\Quote{我们是多么必要、多么有远见与有益地被提醒我们是罪人，因为我们被迫为自己的罪恳求；当向天主恳求宽恕时，灵魂就回忆起自己的罪恶意识。免得有人自认无辜，借着自高而更深地沉沦，他被教导并学习他每日犯罪，因为他被命令每日为自己的罪恳求。这样，\Person{若望}也在他的书信中警告我们，说：\BibleQuote{如果我们说我们没有罪，就是欺骗自己，真理就不在我们内。但如果我们承认我们的罪，祂是忠信正义的，必要赦免我们的罪。} } 他也恰当地在他的致\Person{奎林}书中提出了他对此主题的最绝对判断，并附上了神圣的证言，\Quote{没有人是毫无污秽和没有罪的。} 在那裡，他也写下了确认原罪的证言——这些人企图将它们曲解为我不知道的某种新而邪恶的意义——无论是圣\Person{约伯}所说：\Quote{没有人是纯洁无污的，即使他的生命在地上只有一日，} \BibleRef{约 14:4} 还是圣咏中所读的：\Quote{看，我是在罪孽中怀孕的，我母亲在胎中孕育我时就有了罪。} 对于这些证言，由于那些在成熟年龄已经成圣的人也不例外，他们也不是毫无污秽和罪的，他又加上了最蒙福的\Person{若望}的那句话，除了在其他许多地方，他常提起这句话：\Quote{如果我们说我们没有罪，就是欺骗自己；} \BibleRef{若一 1:8} 以及其他相同含义的段落，这些段落被所有天主教徒所主张，用以反对那些\Quote{欺骗自己，而真理不在他们内}的人。\par

% structure: p0062 source=h2 target=subsection reason=relative-heading-rank; base=h2

\subsection{第二十八章 —— \Person{西彼廉}的正统性毋庸置疑。}

让白拉奇主义者们胆敢说一说，这位天主的人竟被摩尼教的谬误所败坏，以至在赞扬圣人们的同时，又承认此生无人曾达到如此完美的义德，以至于完全没有罪，并以圣经证据的清晰真理和神圣权威来确证他的判断。难道他否认在圣洗中所有罪都被赦免，因为他承认仍有软弱和脆弱，从此他说我们在圣洗后甚至直到此生结束都会犯罪，与肉身的恶习不断争战？或者他不记得宗徒关于无瑕的教会所说的话，他规定没有人应该如此自夸有一颗纯洁无瑕的心，以至于相信自己的无罪，并认为无需给伤口敷药？我想这些新异端者可能会向这位公教人士承认，他知道\Quote{圣神即使在旧时代也帮助善意的倾向}；不仅如此，甚至他们自己也不允许的，即除非通过圣神，他们不可能拥有善意的倾向。我想\Person{西彼廉}知道，所有先知、宗徒或任何时代蒙主喜悦的圣人，都是义人——他们错误地断言我们说\Quote{不是与恶人比较}，而是像他们夸口所说的\Quote{以德行的标准}；然而\Person{西彼廉}说，没有人能没有罪，谁若声称自己无罪，不是骄傲就是愚妄。他理解圣经也不是别的意思：\BibleQuote{谁能夸口说自己有一颗纯洁的心？谁能夸口说自己纯洁无罪？}\BibleRef{箴 20:9}我想\Person{西彼廉}无需被这样的人教导，他很清楚知道\Quote{将来会有善行的赏报和恶行的惩罚，但那时无人能完成在这里可能曾轻视的命令}；然而他理解并断言宗徒\Person{保禄}（他肯定不是轻视天主命令的人）说\BibleQuote{对我来说，生活就是基督，死亡就是收益}\BibleRef{斐 2:21}，没有别的理由，只因为他认为这是最大的收益，在死后不再被世俗的羁绊所束缚，不再受肉身的罪恶和恶习的影响。因此最蒙福的\Person{西彼廉}感到，并在圣经的真理中看到，即使是宗徒们自己的生命，无论多么良善、圣洁、义，也遭受一些世俗羁绊的缠绕，易受一些肉身的罪恶和恶习的影响；他们渴望死亡，以便能摆脱那些邪恶，并能达到那完美的义，这义不会遭受这些事情，并且不再仅仅以命令的方式去成就，而是以赏报的方式去领受。因为甚至当我们祈祷\BibleQuote{愿你的国来临}时，那所要来临的，在天主的国里也没有不义的；因为宗徒说：\BibleQuote{天主的国不在于吃喝，而在于义德、平安和圣神内的喜乐。}\BibleRef{罗 14:17}当然这三件事与其他神圣诫命一同被命令。这里义德被规定给我们，经上说：\BibleQuote{行义！}\BibleRef{依 56:1}平安被规定，经上说：\BibleQuote{你们彼此之间要平安！}\BibleRef{谷 9:49}喜乐被规定，经上说：\BibleQuote{你们要常常在主内喜乐！}\BibleRef{斐 4:4}那么，让白拉奇主义者们否认这些事物会在天主的国里，那里我们将永无止境地生活；或者，如果他们认为好，让他们疯狂到主张义德、平安和喜乐在那里与在这里对义人一样。但如果它们两者都将存在，却又不一样，那么在这里，关于它们的命令，我们要关注实行；在那里，完美要以赏报的方式被希望；那时，不被任何世俗羁绊所阻挡，也不受任何肉身的罪恶和恶习影响（为此宗徒，正如\Person{西彼廉}接受这证言时所说，死亡为他将是收益），我们可以完美地爱天主，他的瞻仰将是面对面；我们也可以完美地爱邻人，因为当心中的思念被显露时，没有人会对任何人产生任何邪恶的怀疑。\par

% structure: p0064 source=h2 target=subsection reason=relative-heading-rank; base=h2

\subsection{第二十九章 [XI]——\Person{安博}反对白拉奇主义者的证言，首先关于原罪。}

如今，为了更充分地驳斥这些人，我再给最光荣的殉道者西彼廉加上最蒙福的安波罗修；因为连白拉奇也如此赞扬他，以至于说在他的著作中甚至连他的敌人都找不到可指责之处。既然白拉奇派声称婴儿出生时没有原罪，并向反对他们、捍卫教会最古老信仰的公教徒扣上摩尼教异端的罪名，那么，就让这位公教的天主之人安波罗修——连白拉奇本人在信仰的真实中也称赞过他——就此事回答他们吧。当他阐释依撒意亚先知时，他说：\Quote{因此，基督没有瑕疵，因为连在出生的通常情况下，祂也没有被玷污。}在同一著作的另一处，谈到宗徒伯多禄时，他说：\Quote{祂奉献了自己——伯多禄先前以为这是罪——并要求为自己不仅洗脚，连头也要洗，因为他直接明白，通过洗脚，对那些在首先的人中跌倒的人，那有害遗传的污秽被消除了。}在同一著作中他还说：\Quote{因此得以保存，即从男人和女人，也就是说通过那种身体的结合，没有人能被视为免于罪恶；但那免于罪恶的，也免于这种受孕方式。}在\WorkTitle{反对诺瓦提安派}的著作中他说：\Quote{我们所有人都是生在罪恶之下。我们的起源本身就在败坏之中，正如你在达味的话中读到的那样：\InnerQuote{看，我是在不义中受孕，我母亲在罪恶中怀了我。}}在\WorkTitle{达味先知的辩护}中他说：\Quote{我们出生之前就染上了传染，在得见光明之前就受到那起源的损害。我们是在不义中受孕。}论及主时，他说：\Quote{那不会拥有身体堕落的罪的那一位，当然不应感受到受生的自然传染。因此，达味以哭泣为自己哀叹这些自然的污秽，以及污点在人生命开始之前就已经存在，这是对的。}同样，在\WorkTitle{诺厄方舟}中他说：\Quote{因此，通过一位主耶稣，将来的救恩被宣告给万民；因为只有祂能是义人，尽管每一世代都迷失，也没有别的理由，只因为祂由童贞女所生，完全不受有罪世代的法令束缚。\InnerQuote{看，}他说，\InnerQuote{我是在不义中受孕，我母亲在罪恶中怀了我。}素来被认为比别人更义的这人如此说。那么，我现在该称谁为义人呢？除非是那脱离这些锁链、不受我们共同本性的束缚所捆绑的那一位？}看，这位圣洁的人，甚至被白拉奇本人证实在公教信仰中是最受认可的，却以如此明显的证据谴责否认原罪的白拉奇派；然而他并不与摩尼教徒一起否认天主是那些出生者的创造者，也不谴责婚姻——这是天主所建立和祝福的。\par

% structure: p0066 source=h2 target=subsection reason=relative-heading-rank; base=h2

\subsection{第三十章——安波罗修关于天主恩宠的见证}

白拉奇派说，功绩由人藉自由意志开始，天主以随后的恩宠助佑作为回报。让可敬的\Person{盎博罗削}也在这里驳斥他们，当他在对依撒意亚先知书的注解中说：\Quote{没有天主的帮助，人的照料无力医治，需要一位神圣的助佑者。}此外，在名为\WorkTitle{论避开世俗}的论著中，他说：\Quote{我们时常谈论避开世俗；但愿我们的心思也能像我们的谈论一样谨慎小心。更糟糕的是，尘世情欲的诱惑常常潜入，虚妄的流淌占据心灵，以致你渴望避免之事，你却在思想中反复思量；这真是人难以提防、更无法摆脱的事。最后，先知证实这更多是愿望而非成就，当他说：\InnerQuote{使我的心倾向祢的诫命，而非贪财。}因为我们的心和思想不在我们手中，它们突然涌出，迷惑心神和灵魂，将它们引向与你所提议的不同方向——它们唤起对时间事物的回忆，提示世俗之事，强行灌输逸乐的思想，编织诱惑的念头，恰恰在我们打算提升心灵的时候，空虚的思想侵入我们，大多时候将我们抛向尘世的事物。谁又能如此幸福，恒常在心中上升？没有天主的帮助，这事怎能成就？绝对不能。最后，旧约圣经也说了同样的话：\InnerQuote{上主，依靠祢的人是有福的，他们的心向上攀登。}}还有什么比这说得更清楚、更充分的呢？但恐怕白拉奇派也许回答说，在祈求天主帮助的这一点上，人的功绩在先，他们说祈祷本身即是功绩，即人藉祈祷渴望天主的恩宠来帮助他，让他们留意同一位圣人在依撒意亚先知书注解中所说的。他说：\Quote{向天主祈祷是一种神恩；因为除非在圣神内，没有人能说耶稣是主。}\BibleRef{格前 12:13}因此，在讲解\BibleBook{路加}{福音}时，他说：\Quote{你确实看到，主的能力处处与人的愿望合作，以致没有主，无人能建造，没有主，无人能承担任何事。}因为像\Person{盎博罗削}这样的人这样说，颂扬天主的恩宠，正如一位应许之子以感恩的虔诚所应当做的，他难道因此而毁灭了自由意志吗？或者，他是否认为恩宠应当被理解为白拉奇派在他们不同的论述中将不得不将恩宠视为仅仅是法律——例如，人们相信天主帮助我们不是去做我们可能知道的事，而是去知道我们可能做的事？如果他们认为这样一位天主的人怀有这种想法，让他们听听他关于法律本身所说的话。在\WorkTitle{论避开世俗}一书中，他说：\Quote{法律能堵住众人的口，却不能改变他们的心意。}在同一论著的另一处，他说：\Quote{法律谴责行为，却不除去其邪恶。}让他们看到这位忠信而公教的人与宗徒一致，宗徒说：\Quote{我们知道：凡法律所说的，都是对那在法律之下的人说的，为叫每一个口都被堵住，使普世都成为天主面前的罪人。因为由于法律，没有一人能在天主面前成义。}\BibleRef{罗 3:19}因为\Person{盎博罗削}正是从宗徒的这一见解中取材并写下了这些事。\par

% structure: p0068 source=h2 target=subsection reason=relative-heading-rank; base=h2

\subsection{第三十一章——\Person{盎博罗削}关于现世义德不完善的见证。}

但如今，既然白拉奇党人声称，在今世中有或曾有义人活过一生而毫无罪过，甚至他们以为来世应盼望作为赏报的生命不能比这更前进或更成全，那么让安波罗削在此也回答他们并驳斥他们吧。因为他在诠释依撒意亚先知论及圣经所写的\BibleQuote{我生育并抚养了儿女，他们却轻视了我}\BibleRef{依 1:2}时，便着手辩论关于出于天主的世代。在那论证中，他引用了若望的证言说\BibleQuote{凡由天主生的，就不犯罪}\BibleRef{若一 3:9}。在处理这同一极难的问题时，他说：\Quote{既然在这世上没有一个人是无罪的；既然若望自己说：\BibleQuote{如果我们说我们没有犯过罪，就是使祂成为说谎者}\BibleRef{若一 1:10}。如果\BibleQuote{凡由天主生的就不犯罪}\BibleRef{若一 3:9}，而这话是指他们中那些在世上的人，那么我们必须视他们为那些藉洗濯的重生获得天主恩宠的无数人群。然而，当先知说：\BibleQuote{这一切都仰望祢，祢按时给它们食物。祢给它们，它们便自己收集；祢张开手，一切都充满美善。但祢转面不顾，它们便惊惶；祢收回它们的气息，它们就断气，归于尘土。祢发出祢的神，它们便被创造；祢更新地面}\BibleRef{咏 104:27-30}时，这样的事似乎不能看为指任何时代，而只能指那未来的时代，那时将有新天新地。因此它们将被惊扰，以便开始。\BibleQuote{祢张开手，一切都充满美善}\BibleRef{咏 104:28}，这不轻易属于此世。因为关于此世，圣经说什么？\BibleQuote{没有行善的，连一个也没有}\BibleRef{咏 14:3}。因此，如果有不同的世代——并且这里进入此生的入口如此接纳罪恶，甚至那生养者也被轻视；而另一个世代不接纳罪恶——让我们考虑是否可能在我们此生结束后有一种重生——关于这重生，《圣经》说：\BibleQuote{在重生时，当人子坐在祂光荣的宝座上}\BibleRef{玛 19:28}。因为正如那称为洗濯的重生，藉此我们从被洗去的罪污中得以更新，同样，那被称为重生的，似乎是由此我们从身体物质的每一污迹中被净化，并在灵魂的纯正感觉中为永生而重生；以致这重生的每一品质都比那洗濯更为纯净，使得对罪恶的猜疑不能落在人的行为上，甚至不能落在人的思想本身。} 此外，在同一著作的另一处，他说：\Quote{我们看到，任何在肉身中被造的人不可能绝对无玷，因为连保禄也说他自己是不成全的。因为他这样说：\BibleQuote{这不是说我已经获得，或已成为成全的了}\BibleRef{斐 3:12}；然而过了一会儿他又说：\BibleQuote{所以我们凡是成全的}\BibleRef{斐 3:15}。除非，或许在此世有一种成全，在此生完结后有另一种成全，正如他对格林多人所说：\BibleQuote{等到那成全的到来}\BibleRef{格前 13:10}；以及别处：\BibleQuote{直到我们众人达到信德的一致，并认识天主子，成为成全的人，达到基督圆满年龄的尺度}\BibleRef{弗 4:13}。因此，正如宗徒说，许多人同他一起被安置在此世，是成全的，但若你论及真正的成全，他们不能是成全的，因为他所说：\BibleQuote{我们现在是藉着镜子观看，模糊不清；到那时，就要面对面了。我现在所认识的，只是局部的；到那时，就要认清了，如同我被认清一样}\BibleRef{格前 13:12}；同样，在此世也有那些\InnerQuote{无玷}的人，在天主的国里也将有\InnerQuote{无玷}的人，尽管确实，若你准确地考虑，无人能是无玷的，因为无人是无罪的。} 在同一著作中他也说：\Quote{我们看到，当我们活在此生时，我们应当净化自己并寻求天主；并从我们灵魂的净化开始，如同建立德行的基础，以便我们配得来生得到我们净化的成全。} 并且再次，在同一著作中他说：\Quote{但重担下呻吟，谁不说：\BibleQuote{我真不幸！谁能救我脱离这死亡的身体}\BibleRef{罗 7:24}？因此，同一位导师，我们给出了各种解释。因为如果那承认自己陷于肉身邪恶中的人是不幸的，当然每个人都是不幸的；因为我不称那人有福，他因头脑的某种黑暗而混乱，不知自己的状况。此外，那并非荒谬地得以理解；因为如果自知的人是不幸的，肯定所有人都是悲惨的，因为每个人或藉智慧承认自己的软弱，或藉愚昧而不知。} 此外，在\WorkTitle{论死亡的好处}一文中，他说：\Quote{让死亡在我们内工作，好让那也产生生命——死后的好生命，即胜利后的好生命，竞赛结束后的好生命；从而不再有肉身的法律知道如何反抗心神的法律，不再与死亡的身体有任何争战。} 再次，在同一篇论文中他说：\Quote{因此，因为义人有这赏报，即看见天主的面，以及那光照每个人的光，让我们从此戴上对这种赏报的渴望，好使我们的灵魂亲近天主，我们的祈祷亲近祂，我们的渴望紧贴祂，使我们不与祂分离。并在此处而被安置，让我们藉默想、阅读、寻求，与天主结合。让我们尽所能认识祂。因为在此处我们只认识祂的一部分；因为此处一切都不成全，那里一切成全；此处我们是婴儿，那里我们将是强健的人。\BibleQuote{我们现在是藉着镜子观看，模糊不清；到那时，就要面对面了}\BibleRef{格前 13:12}。那时，祂的面容被揭示，我们将得以瞻仰天主的荣耀，而我们的灵魂，因被这身体的固渣所缠绕，并被此肉身的一些污点和秽物所遮蔽，现在不能清楚地看见。\BibleQuote{因为谁看见我的面而能生活？}\BibleRef{出 33:20}，并且对。因为若我们的眼睛不能承受太阳的光线——且若任何人长久注视太阳的区域，据说他就失明了——若一个受造物不能无欺骗地注视另一个受造物，那么，一个人怎能不危及自己而注视永恒创造主闪烁的面容，他仍被这身体的衣裳所遮盖？因为在天主面前谁算成义？即使一天的婴儿也不能免于罪，没有人可以夸耀自己的正直和纯洁的心。}\par

% structure: p0070 source=h2 target=subsection reason=relative-heading-rank; base=h2

\subsection{第三十二章［XII］——伯拉纠异端是在盎博罗削之后很久才兴起的。}

倘若我要寻求提及圣盎博罗削针对后来才兴起的伯拉纠异端所说所写的一切，那就太长了。他这样做，不是为了反驳他们，而是为了宣示公教信仰，并在其中建立人。此外，我也不能也不应提及那些最荣耀于主的西彼廉在其书信中所写的一切，这些书信表明我们所持守的信仰是真实、真正基督徒和公教的信仰，正如它自古以来由圣经所传授，并由我们的教父所保存和持守，直到现今这些异端试图摧毁它的时候，并且将来藉天主的恩佑也将被保存和持守。因为这些事情及其同类是如此传授给西彼廉并由西彼廉传达的，这已从他的书信中引用的证词得到证实；它们在我们的时代之前一直如此被持守，这由盎博罗削在这些异端开始肆虐之前就这些事情所写的表明了，公教听众对各地出现的亵渎新论感到战栗；而且，它们将来也必如此被持守，这已部分由对这些观点的谴责，部分由对它们的纠正得到充分有力的宣告。因为无论他们胆敢如何嗡嗡低语反对西彼廉和盎博罗削的健全信仰，我不认为他们会疯狂到胆敢称那些著名和值得纪念的天主之人为摩尼教派。\par

% structure: p0072 source=h2 target=subsection reason=relative-heading-rank; base=h2

\subsection{第三十三章——摩尼教与公教教义的对立。}

那么，他们在其狂怒的心灵盲目中现在散布的是什么呢？说\Quote{几乎在整个西方，一种既不比愚蠢更少也不比不虔敬更次的教条已被采纳}；然而，因天主的仁慈和对祂教会的仁慈治理，公教信仰一直如此警醒，以致\Quote{既不比愚蠢更少也不比邪恶更次}的教条，像摩尼教派的，也像这些异端的，都不应被采纳。如此圣洁和有学问的公教人士——他们藉整个教会的证言被确认为如此——赞扬天主的创造、祂所设立的婚姻、祂藉圣梅瑟所赐的法律、植入人类本性的自由意志，以及圣祖和先知们，并给予适当的宣讲；所有这五件事，摩尼教派部分通过否认、部分也通过憎恶来谴责。由此可见，这些公教博士远离摩尼教派的概念，然而他们断言原罪；他们断言天主的恩宠优先于自由意志，先于一切功德，以致真正提供无偿的神圣助佑；他们断言圣徒们在这肉身中如此正义地生活，以至需要祈祷的帮助，以便每日的罪过得赦免；而且，一种完善而无罪的义德，将是那些在此生正义生活的人在另一个生命的赏报。\par

% structure: p0074 source=h2 target=subsection reason=relative-heading-rank; base=h2

\subsection{第三十四章——谴责异端并非总是需要召集大公会议。}

那么，他们说什么呢？说\Quote{未经任何大公会议集合，就向坐在自己位置上的单纯主教强取了签名}？在如此异端出现之前，最蒙福和信仰卓越的人西彼廉和盎博罗削是否就被强取了反对这类异端的签名？——他们如此清晰地推翻这些异端的邪恶教义，以至我们几乎找不到任何更明显的话来反驳他们。或者，难道需要召集一次大公会议来谴责这公开的祸害，好像不经过大公会议集合就绝不能谴责任何异端？——实际上，我们很少发现某些异端需要为此产生这种必要；而许多且多得无可比拟的异端，在其出现的地方就应受指控和谴责，并从此可以在其余地方被认识并避免。但这些人的骄傲——他们如此高举自己反对天主，以致不愿在祂内夸耀，反而在自由意志中夸耀——也被理解为要攫取这份荣耀，即为了他们的缘故要召集一次东西方大公会议。事实上，他们试图扰乱公教世界，因为主反对他们，他们无法颠覆它；而他们本应被那些牧人的警惕和牧灵勤奋，在任何出现这些狼的地方，在作出充分足够的判断之后，就被践踏；或为了他们能被治愈和改变，或为了其他人能藉牧羊人——祂也在小羊中寻找迷失的羊，祂白白地使羊成圣和成义；祂既明智地教导他们（尽管已被圣化和成义），却因他们的软弱和脆弱，要为他们每日的罪过祈求每日的赦免（没有这赦免，无人能在这世界生活，即使他可能生活得好）；并仁慈地垂听他们的祈祷——的帮助而安全健康地避开他们。\par

\SourceRecord{Translated by Peter Holmes and Robert Ernest Wallis, and revised by Benjamin B. Warfield.；Nicene and Post-Nicene Fathers, First Series；Vol. 5.；Edited by Philip Schaff.；Buffalo, NY: Christian Literature Publishing Co.,；1887.；<http://www.newadvent.org/fathers/15094.htm>.}
